\documentclass[10pt,letterpaper]{book}
\usepackage[margin=0.3in,includeheadfoot]{geometry}
\setlength{\headheight}{14pt}
\usepackage[T1]{fontenc}
\usepackage[utf8]{inputenc}
\usepackage{lmodern}
\usepackage{microtype}
\usepackage{xcolor}
\usepackage{listings}
\usepackage{booktabs}
\usepackage{emptypage}
\usepackage[colorlinks,linkcolor=refcol,urlcolor=refcol]{hyperref}
\definecolor{refcol}{RGB}{20,60,130}
\makeatletter                            % TOC: room for two-digit section numbers (18.10)
\renewcommand*\l@section{\@dottedtocline{1}{1.5em}{3.2em}}
\makeatother

\definecolor{kw}{RGB}{20,60,130}
\definecolor{cm}{RGB}{80,120,80}
\definecolor{st}{RGB}{150,50,50}
\definecolor{rulecol}{RGB}{170,175,190}
\definecolor{numcol}{RGB}{150,150,150}

\lstdefinestyle{code}{language=Python,basicstyle=\ttfamily\small,
  keywordstyle=\color{kw}\bfseries,commentstyle=\color{cm}\itshape,
  stringstyle=\color{st},showstringspaces=false,keepspaces=true,
  columns=fullflexible,breaklines=true,breakatwhitespace=true,
  postbreak=\mbox{\textcolor{numcol}{$\hookrightarrow$}\space},
  upquote=true,aboveskip=6pt,belowskip=6pt,xleftmargin=1.1em,
  frame=leftline,framerule=0.8pt,rulecolor=\color{rulecol}}
\lstdefinestyle{ex}{basicstyle=\ttfamily\small,keepspaces=true,
  columns=fullflexible,breaklines=true,breakatwhitespace=true,
  postbreak=\mbox{\textcolor{numcol}{$\hookrightarrow$}\space},
  upquote=true,aboveskip=4pt,belowskip=7pt,xleftmargin=1.1em,
  frame=leftline,framerule=0.8pt,rulecolor=\color{rulecol},language={}}
\lstdefinestyle{file}{language=Python,basicstyle=\ttfamily\footnotesize,
  keywordstyle=\color{kw}\bfseries,commentstyle=\color{cm}\itshape,
  stringstyle=\color{st},showstringspaces=false,keepspaces=true,
  columns=fullflexible,breaklines=true,upquote=true,numbers=left,
  numberstyle=\tiny\color{numcol},numbersep=7pt,aboveskip=5pt,belowskip=5pt}

\setcounter{tocdepth}{1}
\setlength{\parskip}{2pt}
\title{\Huge\bfseries Prelude Student\\[6pt]\Large A Complete Course in Functional Python\\[14pt]\normalsize sixteen chapters, 193 homework problems, 74 applied labs}
\author{J.~L.~Clements~III}
\date{\today}
\begin{document}
\frontmatter
\maketitle

\chapter{How to Take This Course}
This book teaches \texttt{prelude.py} -- a Haskell-style prelude in pure
Python, no imports anywhere -- until it is not reference material but
reflex. The order is the mastery order, not the file order: the spine is
\textbf{fold $\rightarrow$ scan $\rightarrow$ unfold $\rightarrow$ dict fold
$\rightarrow$ memo} -- one concept in five costumes -- and every chapter
hangs off it.

Each chapter has three movements. \textbf{Concepts}: every function
introduced with its purpose, worked examples (all of them runnable
doctests), deep dives where the idea is subtle, and a ``roll your own''
walkthrough of the implementation itself. \textbf{Homework}: twelve
problems -- five drills, five applied scenarios, two interview-grade
challenges. \textbf{Discipline}: solve in a scratch file with the printed
doctests pasted in; \texttt{python3 -m doctest} is the referee, and an
answer is not done until it is N/N.

Part III then applies the whole toolkit to seventy-four classic problems
organised by class -- the HackerRank/LeetCode canon reworked in prelude
style. Solutions live in the Teacher's Edition and in the repository -- fight each problem first; the doctests will tell you when you have won.

\section*{The recognition ladder}
Before any problem, name the INPUT shape, the OUTPUT shape, and the quantity
word -- then walk this ladder; the first row whose trigger words match names
your tools.

\begin{center}\small
\begin{tabular}{@{}p{0.13\textwidth}p{0.37\textwidth}p{0.41\textwidth}@{}}
\toprule
\textbf{Shape} & \textbf{Trigger words} & \textbf{Reach for} \\
\midrule
SORT FIRST & ``sorted'', ``k-th'', ``closest pair'', ``intervals'' & \texttt{sorted, sortOn, bisect\_left, merge} \\
FOLD & ``total'', ``count'', ``largest'', ``best single \ldots'' & \texttt{foldl, minOn / maxOn} \\
SCAN & ``running'', ``so far'', ``at each step'', ``prefix sums'' & \texttt{scanl, scanl1} \\
DICT FOLD & ``group by'', ``frequency'', ``most common'' & \texttt{fromListWith, Counter, unionWith} \\
WINDOW & ``substring'', ``contiguous'', ``longest run'' & \texttt{longest\_window, pairwise} \\
STACK & ``nested'', ``matching pairs'', ``next greater'' & \texttt{list as stack; foldl with stack acc} \\
UNFOLD & ``digits of'', ``split into pieces'', ``repeat until'' & \texttt{unfoldr, iterate + take} \\
ENUMERATE & ``all subsets'', ``every way to pick'' & \texttt{subsequences, replicateM, cross} \\
TREE / PATHS & ``nested'', ``prefix'' (trie), ``directories'' & \texttt{ITree, paths, setpath, tfold} \\
GRAPH & ``prerequisites'', ``network'', ``shortest route'' & \texttt{Tree(1,list) adjacency, deque BFS, heap, until} \\
MEMOIZE & ``fewest'', ``max value'', ``count the ways'' & \texttt{memo over indices} \\
\bottomrule
\end{tabular}
\end{center}

\tableofcontents
\mainmatter
\part{The Sixteen Chapters}
\chapter{Grammar, Pairs and Maybe}

Every course opens with its grammar. Here it is the conventions the whole
prelude leans on: functions come first and data second in every call;
\texttt{None} is Nothing -- the universal ``no answer'' -- and
\texttt{fromMaybe} is how a pipeline lands a default, and
\texttt{isJust}/\texttt{isNothing} make the None test itself a predicate you
can hand to other tools; \texttt{NOTHING} is the other absence, the ``no
argument supplied'' sentinel. The pair accessors look
trivial, and are: that is what makes them readable inside sort keys and maps
where an index would be noise.
\section{Concepts}
\subsection{\texttt{fst}}\label{def:fst}
the first element of a pair. Reads better than p[0] in code that is all tuples, and slots into higher-order positions.
\begin{lstlisting}[style=ex]
>>> fst((1, 2))
1
>>> map_(fst, [(1, 'a'), (2, 'b')])
[1, 2]
\end{lstlisting}

Roll your own --- the definition:

\begin{lstlisting}[style=ex]
fst = lambda p: p[0]
\end{lstlisting}
\noindent{\small\itshape Used before it is rolled: \texttt{map\_} in \S\ref{def:map-} (p.~\pageref{def:map-}).}

\subsection{\texttt{snd}}\label{def:snd}
the second element of a pair. The classic sort key for (key, count) items.
\begin{lstlisting}[style=ex]
>>> snd((1, 2))
2
>>> sortOn(snd, [('a', 3), ('b', 1)])
[('b', 1), ('a', 3)]
\end{lstlisting}

Roll your own --- the definition:

\begin{lstlisting}[style=ex]
snd = lambda p: p[1]
\end{lstlisting}
\noindent{\small\itshape Used before it is rolled: \texttt{sortOn} in \S\ref{def:sortOn} (p.~\pageref{def:sortOn}).}

\subsection{\texttt{swap}}\label{def:swap}
exchanges a pair's elements. One honest use: inverting a dict by swapping its items.
\begin{lstlisting}[style=ex]
>>> swap((1, 2))
(2, 1)
>>> dict(map_(swap, [('a', 1), ('b', 2)]))
{1: 'a', 2: 'b'}
\end{lstlisting}

Roll your own --- the definition:

\begin{lstlisting}[style=ex]
swap = lambda p: (p[1], p[0])  # Data.Tuple
\end{lstlisting}
\noindent{\small\itshape Used before it is rolled: \texttt{map\_} in \S\ref{def:map-} (p.~\pageref{def:map-}).}

\subsection{\texttt{NOTHING}}\label{def:NOTHING}
a private sentinel object meaning "no argument was supplied". Why not None? Because None is a legitimate VALUE a caller might pass; NOTHING can only mean "omitted", since nobody else has a reference to it. Always test with 'is'. This is how foldl and accumulate tell an absent initial value from a real one.
\begin{lstlisting}[style=ex]
>>> x = NOTHING
>>> x is NOTHING
True
>>> def greet(name=NOTHING):
...     return "hello, stranger" if name is NOTHING else "hello, " + name
>>> greet()
'hello, stranger'
>>> greet("ada")
'hello, ada'
\end{lstlisting}

How it works: object() creates a fresh, featureless object whose only property is its IDENTITY -- nothing else in the program can ever be it. Testing x is NOTHING therefore asks exactly one question: "was I handed this precise sentinel?", which makes it a safe flag for "no argument was supplied".


Why not None, demonstrated: None is data (fromMaybe's whole job is handling it), so using it to also mean "argument omitted" would blur two different situations. The sentinel keeps them apart: None is Nothing in the DATA; NOTHING is Nothing in the ARGUMENT LIST.

\begin{lstlisting}[style=ex]
>>> def describe(x=NOTHING):
...     if x is NOTHING:
...         return "no argument"
...     return "got " + repr(x)
>>> describe()
'no argument'
>>> describe(None)
'got None'
\end{lstlisting}

Where the prelude leans on it: foldl and accumulate take an optional seed. With NOTHING as the default they can tell "seed omitted -- use the first element" from "the seed IS None", which a None default could not.


Remember it as: None is a value; NOTHING is the absence of one.


Roll your own --- the definition:

\begin{lstlisting}[style=ex]
NOTHING = object()  # Maybe's Nothing: "no arg given"; test with `is`
\end{lstlisting}
\subsection{\texttt{fromMaybe}}\label{def:fromMaybe}
lands a default: fromMaybe(d, x) is x unless x is None, in which case it is d. The standard way to finish a pipeline built on the None-is-Nothing convention (find, lookup, getpath, stripPrefix all return None on failure).
\begin{lstlisting}[style=ex]
>>> fromMaybe(0, None), fromMaybe(0, 5)
(0, 5)
>>> fromMaybe('_', find(lambda c: c == 'q', "abc"))
'_'
>>> fromMaybe(99, lookup('b', [('a', 1), ('b', 2)]))
2
\end{lstlisting}

Start from a problem you already have: find looks for something that might not be there, and when it is not, you are left holding None -- which you cannot print to a user, add, or put in a report. So you write a backup plan:

\begin{lstlisting}[style=ex]
>>> result = find(even, [3, 5, 7])
>>> if result is None:
...     result = 0
>>> result
0
\end{lstlisting}

fromMaybe is those three lines as one word. Argument order: the backup plan first, then the risky thing -- said aloud it parses: "use 0 if this comes back empty".

\begin{lstlisting}[style=ex]
>>> fromMaybe(0, find(even, [3, 5, 7]))
0
>>> fromMaybe(0, find(even, [3, 5, 8]))
8
\end{lstlisting}

Versus x or d: Python's or falls back on ANYTHING falsy -- 0, "", [], False -- not just None, so it silently destroys legitimate zero-ish answers. fromMaybe asks the only correct question: is this the absence marker?

\begin{lstlisting}[style=ex]
>>> fromMaybe(99, 0), (0 or 99)
(0, 99)
\end{lstlisting}

Remember it as: ifNoneUse -- the three-line backup plan as one word; or-with-a-conscience.


Roll your own --- the definition:

\begin{lstlisting}[style=ex]
fromMaybe = lambda d, x: d if x is None else x  # Data.Maybe: default for every None-returning tool
\end{lstlisting}
\noindent{\small\itshape Used before it is rolled: \texttt{even} in \S\ref{def:even} (p.~\pageref{def:even}), \texttt{lookup} in \S\ref{def:lookup} (p.~\pageref{def:lookup}), \texttt{find} in \S\ref{def:find} (p.~\pageref{def:find}).}

\subsection{\texttt{isJust}}\label{def:isJust}
the None test as a predicate you can hand to other tools. The lambda it replaces, 'lambda v: v is not None', kept reappearing wherever a Maybe-producing step feeds a predicate slot; naming it also locks in the discipline: identity against None, never truthiness, so 0 and '' count as present.
\begin{lstlisting}[style=ex]
>>> isJust(0), isJust(''), isJust(None)
(True, True, False)
>>> find(isJust, [None, None, 7, None])
7
>>> all(map(isJust, [1, 0, 'x']))
True
\end{lstlisting}

Roll your own --- the definition:

\begin{lstlisting}[style=ex]
isJust = lambda x: x is not None  # Data.Maybe: None test as a handable predicate
\end{lstlisting}
\noindent{\small\itshape Used before it is rolled: \texttt{find} in \S\ref{def:find} (p.~\pageref{def:find}).}

\subsection{\texttt{isNothing}}\label{def:isNothing}
the mirror: True only for None. Counting failures reads as one word.
\begin{lstlisting}[style=ex]
>>> isNothing(None), isNothing(0)
(True, False)
>>> len(filter_(isNothing, [None, 3, None]))
2
\end{lstlisting}

Together they finish the Maybe family: fromMaybe LANDS an optional value, mapMaybe/catMaybes FILTER optional values, isJust/isNothing TEST one. Haskell's Data.Maybe has all five under the same names.

\begin{lstlisting}[style=ex]
>>> readings = [(1, None), (2, 7.5), (3, None)]
>>> fromMaybe(-1, find(isJust, map(snd, readings)))
7.5
\end{lstlisting}

Roll your own --- the definition:

\begin{lstlisting}[style=ex]
isNothing = lambda x: x is None
\end{lstlisting}
\noindent{\small\itshape Used before it is rolled: \texttt{filter\_} in \S\ref{def:filter-} (p.~\pageref{def:filter-}), \texttt{find} in \S\ref{def:find} (p.~\pageref{def:find}).}

\section{Homework}
Work in order: drills cement each tool, applies combine them,
challenges are interview-grade. Write each solution in a scratch file with
the given doctests pasted in, and let \texttt{python3 -m doctest} referee.
\subsection*{1.1\quad Mirror a pair\hfill\textnormal{\textit{drill}}}
Write mirror(p): the pair with its elements exchanged. Applying it twice must give the original pair back.

\noindent\texttt{mirror(p) -> pair}
\begin{lstlisting}[style=ex]
>>> mirror((1, 2))
(2, 1)
>>> mirror(('lo', 'hi'))
('hi', 'lo')
>>> mirror(mirror((3, 4)))
(3, 4)
\end{lstlisting}
\subsection*{1.2\quad Record projections\hfill\textnormal{\textit{drill}}}
A leaderboard record is a (name, score) pair. Write name\_of(r) and score\_of(r) using the pair accessors, not indexing.

\noindent\texttt{name\_of(r) -> value; score\_of(r) -> value}
\begin{lstlisting}[style=ex]
>>> r = ('ada', 95)
>>> name_of(r)
'ada'
>>> score_of(r)
95
>>> name_of(swap(r))
95
\end{lstlisting}
\subsection*{1.3\quad Default when missing\hfill\textnormal{\textit{drill}}}
Write or\_default(d, x): x itself unless x is None, in which case d. Falsy values like 0 and '' are real data and must survive.

\noindent\texttt{or\_default(d, x) -> value}
\begin{lstlisting}[style=ex]
>>> or_default(0, None)
0
>>> or_default(0, 5)
5
>>> or_default(9, 0)
0
>>> or_default('-', None)
'-'
\end{lstlisting}
\subsection*{1.4\quad The sentinel argument\hfill\textnormal{\textit{drill}}}
Write greet(name=NOTHING) returning 'hello, stranger' when NO argument was passed, and 'hello, ' + str(name) otherwise -- including when the caller passes None on purpose.

\noindent\texttt{greet(name=NOTHING) -> str}
\begin{lstlisting}[style=ex]
>>> greet()
'hello, stranger'
>>> greet('ada')
'hello, ada'
>>> greet(None)
'hello, None'
\end{lstlisting}
\subsection*{1.5\quad Manufactured defaults\hfill\textnormal{\textit{drill}}}
Write make\_default(v): a function that returns a one-argument function answering v no matter what it is asked.

\noindent\texttt{make\_default(v) -> function}
\begin{lstlisting}[style=ex]
>>> make_default(7)('anything')
7
>>> f = make_default(0)
>>> f(99)
0
>>> f('x')
0
\end{lstlisting}
\noindent{\small\itshape Rolled later: \texttt{const} in \S\ref{def:const} (p.~\pageref{def:const}).}

\subsection*{1.6\quad Leaderboard by name\hfill\textnormal{\textit{apply}}}
The scoring service emits (score, name) pairs, but the display wants (name, score). Write by\_name(pairs) converting the whole list, order preserved.

\noindent\texttt{by\_name(pairs) -> list of pairs}
\begin{lstlisting}[style=ex]
>>> by_name([(95, 'ada'), (87, 'bob')])
[('ada', 95), ('bob', 87)]
>>> by_name([])
[]
>>> by_name([(1, 'solo')])
[('solo', 1)]
\end{lstlisting}
\subsection*{1.7\quad Config with fallbacks\hfill\textnormal{\textit{apply}}}
Settings live in a dict; an absent key means 'use the default'. Write setting(cfg, key, default). A stored falsy value (False, 0, '') is a deliberate choice and must be returned as-is.

\noindent\texttt{setting(cfg, key, default) -> value}
\begin{lstlisting}[style=ex]
>>> setting({'retries': 3}, 'retries', 1)
3
>>> setting({}, 'retries', 1)
1
>>> setting({'debug': False}, 'debug', True)
False
\end{lstlisting}
\subsection*{1.8\quad Invert the phone book\hfill\textnormal{\textit{apply}}}
A phone book maps names to numbers, and numbers are unique. Produce the reverse index mapping each number back to its name.

\noindent\texttt{invert\_book(d) -> dict}
\begin{lstlisting}[style=ex]
>>> invert_book({'ada': '555-01', 'bob': '555-02'})
{'555-01': 'ada', '555-02': 'bob'}
>>> invert_book({})
{}
>>> invert_book({'x': 1})
{1: 'x'}
\end{lstlisting}
\subsection*{1.9\quad First good sensor fix\hfill\textnormal{\textit{apply}}}
A sensor log is a list of (timestamp, reading) pairs where a failed sample records None. Return the first non-None reading, or the supplied default when every sample failed (or the log is empty). A reading of 0 is a real fix.

\noindent\texttt{first\_fix(default, log) -> value}
\begin{lstlisting}[style=ex]
>>> first_fix(-1, [(1, None), (2, 7.5), (3, 8.0)])
7.5
>>> first_fix(-1, [(1, None), (2, None)])
-1
>>> first_fix(-1, [])
-1
>>> first_fix(-1, [(1, 0)])
0
\end{lstlisting}
\noindent{\small\itshape Rolled later: \texttt{find} in \S\ref{def:find} (p.~\pageref{def:find}).}

\subsection*{1.10\quad Effective badge ids\hfill\textnormal{\textit{apply}}}
Each visitor holds (badge\_or\_None, temp\_id). The effective id is the badge when one was issued, else the temp id. Map a whole visitor list to effective ids; an empty-string badge is a real (blank) badge, not a missing one.

\noindent\texttt{effective\_ids(pairs) -> list}
\begin{lstlisting}[style=ex]
>>> effective_ids([(None, 'T1'), ('B2', 'T2'), (None, 'T3')])
['T1', 'B2', 'T3']
>>> effective_ids([])
[]
>>> effective_ids([('', 'T9')])
['']
\end{lstlisting}
\subsection*{1.11\quad Ledger reconciliation\hfill\textnormal{\textit{challenge}}}
Ledger rows are (txn\_id, debit\_or\_None, credit\_or\_None); a None side simply was not recorded and counts as 0 toward the balance. Return the pair (net, incomplete): net is total credits minus total debits; incomplete counts rows where EITHER side is missing.

\noindent\texttt{reconcile(rows) -> (net, incomplete)}
\begin{lstlisting}[style=ex]
>>> reconcile([('a', 5, None), ('b', None, 10), ('c', 2, 3)])
(6, 2)
>>> reconcile([])
(0, 0)
>>> reconcile([('x', 1, 4)])
(3, 0)
>>> reconcile([('y', None, None)])
(0, 1)
\end{lstlisting}
\subsection*{1.12\quad Three-level settings cascade\hfill\textnormal{\textit{challenge}}}
Configuration resolves through a cascade: an environment dict overrides a config-file dict, which overrides a built-in default. Absence at a level means None from .get; stored falsy values win like any other value. Write resolve(key, env, cfg, default).

\noindent\texttt{resolve(key, env, cfg, default) -> value}
\begin{lstlisting}[style=ex]
>>> resolve('p', {'p': 1}, {'p': 2}, 9)
1
>>> resolve('p', {}, {'p': 2}, 9)
2
>>> resolve('p', {}, {}, 9)
9
>>> resolve('debug', {}, {'debug': 0}, 1)
0
\end{lstlisting}
\clearpage
\chapter{Arithmetic and Logic}

Small tools with sharp edges. The floor-versus-truncate split
(\texttt{div}/\texttt{mod} against \texttt{quot}/\texttt{rem}) decides real
interview questions the moment a negative number appears; \texttt{signum}
compresses three-way comparison into arithmetic; \texttt{isqrt} shows Newton's
method surviving entirely inside the integers. Master the edges here and
modular-arithmetic problems stop being scary.
\section{Concepts}
\subsection{\texttt{succ}}\label{def:succ}
successor: one more; on a character (Haskell's Char instance of Enum), the next codepoint.
\begin{lstlisting}[style=ex]
>>> succ(41)
42
>>> succ('a')
'b'
>>> map_(succ, [1, 2, 3])
[2, 3, 4]
\end{lstlisting}

Roll your own --- the definition:

\begin{lstlisting}[style=ex]
succ = lambda x: chr(ord(x) + 1) if isinstance(x, str) else x + 1
\end{lstlisting}
\noindent{\small\itshape Used before it is rolled: \texttt{map\_} in \S\ref{def:map-} (p.~\pageref{def:map-}).}

\subsection{\texttt{pred}}\label{def:pred}
predecessor: one less; characters step back one codepoint.
\begin{lstlisting}[style=ex]
>>> pred(43)
42
>>> pred('b')
'a'
\end{lstlisting}

Roll your own --- the definition:

\begin{lstlisting}[style=ex]
pred = lambda x: chr(ord(x) - 1) if isinstance(x, str) else x - 1  # Enum: numbers and Chars
\end{lstlisting}
\subsection{\texttt{even}}\label{def:even}
is n divisible by two? A predicate shaped for filter\_ and friends.
\begin{lstlisting}[style=ex]
>>> even(4), even(7)
(True, False)
>>> filter_(even, [1, 2, 3, 4])
[2, 4]
\end{lstlisting}

Roll your own --- the definition:

\begin{lstlisting}[style=ex]
even = lambda n: n % 2 == 0
\end{lstlisting}
\noindent{\small\itshape Used before it is rolled: \texttt{filter\_} in \S\ref{def:filter-} (p.~\pageref{def:filter-}).}

\subsection{\texttt{odd}}\label{def:odd}
the complement of even.
\begin{lstlisting}[style=ex]
>>> odd(4), odd(7)
(False, True)
>>> partition(odd, [1, 2, 3, 4])
([1, 3], [2, 4])
\end{lstlisting}

Roll your own --- the definition:

\begin{lstlisting}[style=ex]
odd = lambda n: n % 2 == 1
\end{lstlisting}
\noindent{\small\itshape Used before it is rolled: \texttt{partition} in \S\ref{def:partition} (p.~\pageref{def:partition}).}

\subsection{\texttt{not\_}}\label{def:not-}
logical negation as a function (''not'' is a keyword, so it cannot be passed around; not\_ can). Composes predicates.
\begin{lstlisting}[style=ex]
>>> not_(True)
False
>>> is_odd = compose(not_, even)
>>> is_odd(3)
True
\end{lstlisting}

Roll your own --- the definition:

\begin{lstlisting}[style=ex]
not_ = lambda b: not b  # `not` is a Python keyword
\end{lstlisting}
\noindent{\small\itshape Used before it is rolled: \texttt{compose} in \S\ref{def:compose} (p.~\pageref{def:compose}).}

\subsection{\texttt{otherwise}}\label{def:otherwise}
just True, wearing Haskell's name for the catch-all guard. Its habitat is the RULES TABLE: a long if/elif ladder rewritten as data -- (predicate, result) pairs tried in order, find picking the first match. One care: Haskell guards are EVALUATED, so bare otherwise works there; a table's guards are APPLIED to the argument, so wrap it as const(otherwise) -- the guard that ignores its input and says yes. Result slots follow the same rule: fixed answers wear const, computed ones (like str below) go bare.
\begin{lstlisting}[style=ex]
>>> otherwise
True
>>> RULES = [(lambda n: n % 15 == 0, const("FizzBuzz")),
...          (lambda n: n % 3 == 0,  const("Fizz")),
...          (lambda n: n % 5 == 0,  const("Buzz")),
...          (const(otherwise),      str)]
>>> speak = lambda n: snd(find(lambda rule: fst(rule)(n), RULES))(n)
>>> map_(speak, [1, 3, 5, 15])
['1', 'Fizz', 'Buzz', 'FizzBuzz']
\end{lstlisting}

Roll your own --- the definition:

\begin{lstlisting}[style=ex]
otherwise = True
\end{lstlisting}
\noindent{\small\itshape Used before it is rolled: \texttt{map\_} in \S\ref{def:map-} (p.~\pageref{def:map-}), \texttt{find} in \S\ref{def:find} (p.~\pageref{def:find}), \texttt{const} in \S\ref{def:const} (p.~\pageref{def:const}).}

\subsection{\texttt{signum}}\label{def:signum}
the sign of a number as -1, 0 or 1, computed by arithmetic on two comparisons -- no branches. Also a ready-made comparator result.
\begin{lstlisting}[style=ex]
>>> signum(-7), signum(0), signum(3)
(-1, 0, 1)
>>> sortBy(lambda a, b: signum(a - b), [3, 1, 2])
[1, 2, 3]
\end{lstlisting}

Roll your own --- the definition:

\begin{lstlisting}[style=ex]
signum = lambda x: (x > 0) - (x < 0)
\end{lstlisting}
\noindent{\small\itshape Used before it is rolled: \texttt{sortBy} in \S\ref{def:sortBy} (p.~\pageref{def:sortBy}).}

\subsection{\texttt{div}}\label{def:div}
integer division that FLOORS (rounds toward negative infinity), like Haskell's div and Python's //. The div/mod pair and the quot/rem pair agree on positives and split on negatives -- interviews live in that gap.
\begin{lstlisting}[style=ex]
>>> div(7, 2), div(-7, 2)
(3, -4)
\end{lstlisting}

Roll your own --- the definition:

\begin{lstlisting}[style=ex]
div = lambda a, b: a // b  # floors, like Haskell div
\end{lstlisting}
\subsection{\texttt{mod}}\label{def:mod}
the remainder that matches div: its sign follows the DIVISOR, and div(a,b)*b + mod(a,b) == a always holds.
\begin{lstlisting}[style=ex]
>>> mod(7, 2), mod(-7, 2), mod(7, -2)
(1, 1, -1)
\end{lstlisting}

Roll your own --- the definition:

\begin{lstlisting}[style=ex]
mod = lambda a, b: a % b
\end{lstlisting}
\subsection{\texttt{quot}}\label{def:quot}
integer division that TRUNCATES (rounds toward zero), like Haskell's quot and C's /. Same as div for positives; differs by one for negatives.
\begin{lstlisting}[style=ex]
>>> quot(7, 2), quot(-7, 2)
(3, -3)
\end{lstlisting}

Roll your own --- the definition:

\begin{lstlisting}[style=ex]
quot = lambda a, b: -(-a // b) if (a < 0) != (b < 0) else a // b  # truncates, like Haskell quot
\end{lstlisting}
\subsection{\texttt{rem}}\label{def:rem}
the remainder that matches quot: its sign follows the DIVIDEND.
\begin{lstlisting}[style=ex]
>>> rem(7, 2), rem(-7, 2), rem(7, -2)
(1, -1, 1)
\end{lstlisting}

Roll your own --- the definition:

\begin{lstlisting}[style=ex]
rem = lambda a, b: a - b * quot(a, b)
\end{lstlisting}
\subsection{\texttt{quotRem}}\label{def:quotRem}
truncating quotient and remainder at once.
\begin{lstlisting}[style=ex]
>>> quotRem(-7, 2)
(-3, -1)
\end{lstlisting}

Roll your own --- the definition:

\begin{lstlisting}[style=ex]
quotRem = lambda a, b: (quot(a, b), rem(a, b))
\end{lstlisting}
\subsection{\texttt{gcd}}\label{def:gcd}
greatest common divisor by Euclid's recursion, always >= 0 like Haskell's. Recursion is safe here: the depth is at most \~{}90 even for 64-bit inputs (worst case: consecutive Fibonaccis).
\begin{lstlisting}[style=ex]
>>> gcd(48, 18)
6
>>> gcd(-4, 6), gcd(17, 5)
(2, 1)
\end{lstlisting}

Roll your own --- the definition:

\begin{lstlisting}[style=ex]
gcd = lambda a, b: abs(a) if b == 0 else gcd(b, a % b)  # >= 0 like Haskell; depth <= 90
\end{lstlisting}
\subsection{\texttt{lcm}}\label{def:lcm}
least common multiple via gcd, also always >= 0; defined as 0 when either input is 0.
\begin{lstlisting}[style=ex]
>>> lcm(4, 6)
12
>>> lcm(-4, 6), lcm(0, 5)
(12, 0)
\end{lstlisting}

Roll your own --- the definition:

\begin{lstlisting}[style=ex]
lcm = lambda a, b: abs(a // gcd(a, b) * b) if a and b else 0  # >= 0 like Haskell
\end{lstlisting}
\subsection{\texttt{hypot}}\label{def:hypot}
the Euclidean distance sqrt(x\^{}2 + y\^{}2). Note it returns a float.
\begin{lstlisting}[style=ex]
>>> hypot(3, 4)
5.0
\end{lstlisting}

Roll your own --- the definition:

\begin{lstlisting}[style=ex]
hypot = lambda x, y: (x*x + y*y) ** 0.5
\end{lstlisting}
\subsection{\texttt{isqrt}}\label{def:isqrt}
floor square root using Newton's method on INTEGERS -- no floats, so no precision loss on huge numbers; a negative input raises rather than lying.
\begin{lstlisting}[style=ex]
>>> isqrt(16), isqrt(17)
(4, 4)
>>> isqrt(10**18)
1000000000
\end{lstlisting}

Roll your own: Newton's iteration on integers -- repeatedly replace x with the average of x and n // x. It converges from above, so loop while the new estimate still shrinks; all arithmetic is //, so no float (and no float rounding lie) ever appears.

\begin{lstlisting}[style=ex]
def isqrt(n):  # floor sqrt without floats (Newton)
    if n < 0:
        raise ValueError("isqrt of negative")
    x = n
    y = (x + 1) // 2
    while y < x:
        x, y = y, (y + n // y) // 2
    return x if n else 0
\end{lstlisting}
\section{Homework}
Work in order: drills cement each tool, applies combine them,
challenges are interview-grade. Write each solution in a scratch file with
the given doctests pasted in, and let \texttt{python3 -m doctest} referee.
\subsection*{2.1\quad Parity census\hfill\textnormal{\textit{drill}}}
Count how many values are even and how many are odd, as a pair (evens, odds). Zero is even; negatives carry their usual parity.

\noindent\texttt{parity\_counts(xs) -> (evens, odds)}
\begin{lstlisting}[style=ex]
>>> parity_counts([1, 2, 3, 4])
(2, 2)
>>> parity_counts([])
(0, 0)
>>> parity_counts([-2, -3])
(1, 1)
>>> parity_counts([0])
(1, 0)
\end{lstlisting}
\subsection*{2.2\quad Floor versus truncate\hfill\textnormal{\textit{drill}}}
Write both\_divs(a, b) returning (div(a, b), quot(a, b)) so the two integer divisions can be compared side by side. They agree on positives and split on mixed signs.

\noindent\texttt{both\_divs(a, b) -> (floored, truncated)}
\begin{lstlisting}[style=ex]
>>> both_divs(7, 2)
(3, 3)
>>> both_divs(-7, 2)
(-4, -3)
>>> both_divs(7, -2)
(-4, -3)
>>> both_divs(-7, -2)
(3, 3)
\end{lstlisting}
\subsection*{2.3\quad Character neighbours\hfill\textnormal{\textit{drill}}}
Using the Enum behaviour of succ and pred on characters, return the pair of codepoint neighbours of a character: (previous, next). Works on letters and digits alike.

\noindent\texttt{char\_neighbours(c) -> (prev\_char, next\_char)}
\begin{lstlisting}[style=ex]
>>> char_neighbours('b')
('a', 'c')
>>> char_neighbours('m')
('l', 'n')
>>> char_neighbours('5')
('4', '6')
\end{lstlisting}
\subsection*{2.4\quad Perfect square test\hfill\textnormal{\textit{drill}}}
Decide whether a non-negative integer is a perfect square, with no floating point anywhere -- float sqrt lies for large inputs, isqrt does not.

\noindent\texttt{is\_square(n) -> bool}
\begin{lstlisting}[style=ex]
>>> is_square(16)
True
>>> is_square(17)
False
>>> is_square(0)
True
>>> is_square(10**18)
True
>>> is_square(10**18 + 1)
False
\end{lstlisting}
\subsection*{2.5\quad Coprime pairs\hfill\textnormal{\textit{drill}}}
Two integers are coprime when they share no factor beyond 1 -- that is, their gcd is exactly 1. Signs do not matter.

\noindent\texttt{coprime(a, b) -> bool}
\begin{lstlisting}[style=ex]
>>> coprime(8, 15)
True
>>> coprime(6, 9)
False
>>> coprime(1, 99)
True
>>> coprime(-3, 4)
True
\end{lstlisting}
\subsection*{2.6\quad Meeting end on a 24-hour clock\hfill\textnormal{\textit{apply}}}
A meeting starts at start\_hour (0..23) and runs for duration hours, possibly negative (rescheduled earlier) or longer than a day. Return the ending hour on the 24-hour clock. Flooring mod does all the wrap-around work.

\noindent\texttt{end\_hour(start, duration) -> hour in 0..23}
\begin{lstlisting}[style=ex]
>>> end_hour(23, 2)
1
>>> end_hour(9, 8)
17
>>> end_hour(0, 48)
0
>>> end_hour(10, -12)
22
\end{lstlisting}
\subsection*{2.7\quad Pages needed\hfill\textnormal{\textit{apply}}}
A report has n items and each page holds k. How many pages are needed? This is ceiling division, and the branch-free trick is to floor-divide the NEGATION: -div(-n, k).

\noindent\texttt{pages(n, k) -> int}
\begin{lstlisting}[style=ex]
>>> pages(10, 3)
4
>>> pages(9, 3)
3
>>> pages(0, 5)
0
>>> pages(1, 10)
1
\end{lstlisting}
\subsection*{2.8\quad Reduced aspect ratio\hfill\textnormal{\textit{apply}}}
Given a screen's width and height in pixels, report its aspect ratio in lowest terms as a pair. 1920x1080 is famously 16:9.

\noindent\texttt{ratio(w, h) -> (w\_reduced, h\_reduced)}
\begin{lstlisting}[style=ex]
>>> ratio(1920, 1080)
(16, 9)
>>> ratio(100, 100)
(1, 1)
>>> ratio(7, 3)
(7, 3)
\end{lstlisting}
\subsection*{2.9\quad Cron co-firing\hfill\textnormal{\textit{apply}}}
Two cron jobs fire every a minutes and every b minutes, and both just fired together. In how many minutes do they next fire together again? That is the least common multiple.

\noindent\texttt{next\_together(a, b) -> minutes}
\begin{lstlisting}[style=ex]
>>> next_together(4, 6)
12
>>> next_together(5, 7)
35
>>> next_together(10, 10)
10
>>> next_together(1, 9)
9
\end{lstlisting}
\subsection*{2.10\quad Signal triage\hfill\textnormal{\textit{apply}}}
Sensor deltas arrive as integers. Triage them into (negative, zero, positive) counts using signum so the branching lives in arithmetic, not in if-chains.

\noindent\texttt{sign\_census(xs) -> (negatives, zeros, positives)}
\begin{lstlisting}[style=ex]
>>> sign_census([-2, 0, 3, 4])
(1, 1, 2)
>>> sign_census([])
(0, 0, 0)
>>> sign_census([0, 0])
(0, 2, 0)
\end{lstlisting}
\noindent{\small\itshape Rolled later: \texttt{count} in \S\ref{def:count} (p.~\pageref{def:count}).}

\subsection*{2.11\quad Digital root\hfill\textnormal{\textit{challenge}}}
The digital root of a non-negative integer repeatedly replaces the number by the sum of its digits until a single digit remains: 38 -> 11 -> 2. Compute it with integer arithmetic only -- no string conversion.

\noindent\texttt{digital\_root(n) -> digit}
\begin{lstlisting}[style=ex]
>>> digital_root(0)
0
>>> digital_root(9)
9
>>> digital_root(38)
2
>>> digital_root(12345)
6
>>> digital_root(999)
9
\end{lstlisting}
\subsection*{2.12\quad Visible lattice points\hfill\textnormal{\textit{challenge}}}
Standing at the origin of an integer grid, the point (x, y) is visible when no other lattice point lies on the straight segment between you and it -- which happens exactly when gcd(x, y) == 1. The origin itself (gcd 0) is not visible. Count the visible points in a list.

\noindent\texttt{visible\_count(points) -> int}
\begin{lstlisting}[style=ex]
>>> visible_count([(1, 1), (2, 2), (0, 1), (0, 0), (3, 5)])
3
>>> visible_count([])
0
>>> visible_count([(0, 0)])
0
>>> visible_count([(1, 0), (0, 1)])
2
\end{lstlisting}
\clearpage
\chapter{Lists I: The Basics}

The list is the prelude's universal substrate, and this chapter is its
anatomy: take lists apart at the ends, cut them at positions, pair them
against each other -- and against their own tails, which is what
\texttt{pairwise} does and why every rule about ``adjacent elements'' starts
there. \texttt{nub}, \texttt{transpose} and the prefix family each hide one
idea worth owning outright.
\section{Concepts}
\subsection{\texttt{head}}\label{def:head}
the first element. A PARTIAL function: on an empty list it raises, so guard with null() or reach for find/fromMaybe when emptiness is possible.
\begin{lstlisting}[style=ex]
>>> head([1, 2, 3])
1
>>> head([])
Traceback (most recent call last):
  ...
IndexError: list index out of range
\end{lstlisting}

Roll your own --- the definition:

\begin{lstlisting}[style=ex]
head = lambda xs: xs[0]
\end{lstlisting}
\noindent{\small\itshape Used before it is rolled: \texttt{last} in \S\ref{def:last} (p.~\pageref{def:last}).}

\subsection{\texttt{tail}}\label{def:tail}
everything after the first element. head and tail together are the grammar of structural recursion: do something with head, recurse on tail.
\begin{lstlisting}[style=ex]
>>> tail([1, 2, 3])
[2, 3]
>>> tail([1])
[]
\end{lstlisting}

Roll your own --- the definition:

\begin{lstlisting}[style=ex]
tail = lambda xs: xs[1:]
\end{lstlisting}
\subsection{\texttt{init}}\label{def:init}
everything before the last element (the mirror of tail).
\begin{lstlisting}[style=ex]
>>> init([1, 2, 3])
[1, 2]
\end{lstlisting}

Roll your own --- the definition:

\begin{lstlisting}[style=ex]
init = lambda xs: xs[:-1]
\end{lstlisting}
\subsection{\texttt{last}}\label{def:last}
the final element. Partial, like head.
\begin{lstlisting}[style=ex]
>>> last([1, 2, 3])
3
\end{lstlisting}

Roll your own --- the definition:

\begin{lstlisting}[style=ex]
last = lambda xs: xs[-1]
\end{lstlisting}
\subsection{\texttt{null}}\label{def:null}
is the list empty? The guard that makes head/tail safe to use.
\begin{lstlisting}[style=ex]
>>> null([]), null([1])
(True, False)
\end{lstlisting}

Roll your own --- the definition:

\begin{lstlisting}[style=ex]
null = lambda xs: len(xs) == 0
\end{lstlisting}
\subsection{\texttt{elem}}\label{def:elem}
membership test, argument order (needle, haystack).
\begin{lstlisting}[style=ex]
>>> elem(2, [1, 2, 3])
True
>>> elem('q', "abc")
False
\end{lstlisting}

Roll your own --- the definition:

\begin{lstlisting}[style=ex]
elem = lambda x, xs: x in xs
\end{lstlisting}
\subsection{\texttt{notElem}}\label{def:notElem}
non-membership.
\begin{lstlisting}[style=ex]
>>> notElem(9, [1, 2, 3])
True
\end{lstlisting}

Roll your own --- the definition:

\begin{lstlisting}[style=ex]
notElem = lambda x, xs: x not in xs
\end{lstlisting}
\subsection{\texttt{replicate}}\label{def:replicate}
n copies of one value, as a list.
\begin{lstlisting}[style=ex]
>>> replicate(3, 'x')
['x', 'x', 'x']
>>> replicate(0, 'x')
[]
\end{lstlisting}

Roll your own --- the definition:

\begin{lstlisting}[style=ex]
replicate = lambda n, x: [x] * n
\end{lstlisting}
\subsection{\texttt{drop}}\label{def:drop}
discard the first n elements, keep the rest; finite lists only (it materialises the input), and a negative n drops nothing.
\begin{lstlisting}[style=ex]
>>> drop(2, [1, 2, 3, 4])
[3, 4]
>>> drop(10, [1, 2]), drop(-3, [1, 2])
([], [1, 2])
\end{lstlisting}

Roll your own --- the definition:

\begin{lstlisting}[style=ex]
drop = lambda n, xs: list(xs)[max(n, 0):]
\end{lstlisting}
\subsection{\texttt{splitAt}}\label{def:splitAt}
cut a list into (first n, rest) in one call; a negative n splits at the front. The pieces glue back with +, which is why rotations and chunking fall out of it.
\begin{lstlisting}[style=ex]
>>> splitAt(2, [1, 2, 3, 4])
([1, 2], [3, 4])
>>> splitAt(0, [1, 2]), splitAt(-2, [1, 2])
(([], [1, 2]), ([], [1, 2]))
\end{lstlisting}

Roll your own --- the definition:

\begin{lstlisting}[style=ex]
splitAt = lambda n, xs: (xs[:n], xs[n:]) if n >= 0 else (xs[:0], xs)  # n<0 splits at the front
\end{lstlisting}
\subsection{\texttt{nub}}\label{def:nub}
deduplicate, FIRST occurrence wins, order preserved. The trick: dict.fromkeys builds a dict (whose keys are unique and insertion-ordered) and throws away the values. Elements must be hashable. Compare set(xs), which destroys order.
\begin{lstlisting}[style=ex]
>>> nub([3, 1, 3, 2, 1])
[3, 1, 2]
>>> nub("mississippi")
['m', 'i', 's', 'p']
\end{lstlisting}

Roll your own --- the definition:

\begin{lstlisting}[style=ex]
nub = lambda xs: list(dict.fromkeys(xs))
\end{lstlisting}
\subsection{\texttt{lookup}}\label{def:lookup}
the first value paired with a key in an association list (a list of (key, value) tuples), or None if absent. First match wins, so an assoc list can shadow like a scope chain.
\begin{lstlisting}[style=ex]
>>> lookup('b', [('a', 1), ('b', 2)])
2
>>> lookup('a', [('a', 1), ('a', 99)])
1
>>> lookup('z', [('a', 1)]) is None
True
\end{lstlisting}

Roll your own --- the definition:

\begin{lstlisting}[style=ex]
lookup = lambda k, pairs: next((v for kk, v in pairs if kk == k), None)  # Nothing -> None
\end{lstlisting}
\subsection{\texttt{zip\_}}\label{def:zip-}
pair two iterables positionally, stopping at the shorter; returns a real list (the builtin zip returns a lazy iterator), and either side may be any iterable, streams included.
\begin{lstlisting}[style=ex]
>>> zip_([1, 2], [10, 20, 30])
[(1, 10), (2, 20)]
>>> zip_("ab", count(0))
[('a', 0), ('b', 1)]
\end{lstlisting}

Roll your own --- the definition:

\begin{lstlisting}[style=ex]
zip_ = lambda a, b: list(zip(a, b))  # shortest wins, any iterable
\end{lstlisting}
\noindent{\small\itshape Used before it is rolled: \texttt{count} in \S\ref{def:count} (p.~\pageref{def:count}).}

\subsection{\texttt{zip3}}\label{def:zip3}
the three-list version.
\begin{lstlisting}[style=ex]
>>> zip3([1, 2], ['a', 'b'], [True, False])
[(1, 'a', True), (2, 'b', False)]
\end{lstlisting}

Roll your own --- the definition:

\begin{lstlisting}[style=ex]
zip3 = lambda a, b, c: list(zip(a, b, c))
\end{lstlisting}
\subsection{\texttt{unzip}}\label{def:unzip}
a list of pairs back into a pair of lists; the inverse of zip\_.
\begin{lstlisting}[style=ex]
>>> unzip([(1, 'a'), (2, 'b')])
([1, 2], ['a', 'b'])
>>> unzip([])
([], [])
\end{lstlisting}

Roll your own --- the definition:

\begin{lstlisting}[style=ex]
unzip = lambda ps: tuple(map(list, zip(*ps))) if ps else ([], [])
\end{lstlisting}
\subsection{\texttt{transpose}}\label{def:transpose}
rows become columns, RAGGED-SAFE exactly like Haskell's: short rows simply drop out of later columns; nothing is silently truncated.
\begin{lstlisting}[style=ex]
>>> transpose([[1, 2, 3], [4, 5, 6]])
[[1, 4], [2, 5], [3, 6]]
>>> transpose([[1, 2, 3], [4, 5]])
[[1, 4], [2, 5], [3]]
>>> transpose([])
[]
\end{lstlisting}

Rows may be any iterables, even one-shot ones -- the rows are materialised first, so a generator of iterators transposes like a list of lists (without that step the generator is spent by the length scan):

\begin{lstlisting}[style=ex]
>>> transpose(iter([1, 2]) for _ in range(2))
[[1, 1], [2, 2]]
>>> transpose(["abc", "de"])
[['a', 'd'], ['b', 'e'], ['c']]
\end{lstlisting}

Roll your own: materialise the rows, find the longest, then for each column index keep xs[i] only from rows that reach it -- the 'if i < len(xs)' filter is precisely where short rows drop out of later columns instead of truncating everyone.

\begin{lstlisting}[style=ex]
def transpose(xss):  # Data.List: rows <-> cols, RAGGED-SAFE like Haskell
    xss = [list(xs) for xs in xss]  # any iterables in (rows may be one-shot); materialise once
    n = max(map(len, xss), default=0)  # short rows just drop out of later columns (no truncation)
    return [[xs[i] for xs in xss if i < len(xs)] for i in range(n)]
\end{lstlisting}
\subsection{\texttt{enum}}\label{def:enum}
pair every element with its index: a list-returning enumerate. The shape that feeds index-aware folds (two\_sum's seen-dict, next\_greater's stack).
\begin{lstlisting}[style=ex]
>>> enum(['a', 'b'])
[(0, 'a'), (1, 'b')]
>>> enum(['a', 'b'], start=1)
[(1, 'a'), (2, 'b')]
\end{lstlisting}

Roll your own --- the definition:

\begin{lstlisting}[style=ex]
enum = lambda xs, start=0: zip_(list(range(start, start + len(xs))), xs)
\end{lstlisting}
\subsection{\texttt{pairwise}}\label{def:pairwise}
every element with its RIGHT neighbour: zip xs (tail xs). The tool for any rule about adjacency -- deltas, local comparisons, Roman numeral subtraction.
\begin{lstlisting}[style=ex]
>>> pairwise([1, 4, 9])
[(1, 4), (4, 9)]
>>> map_(lambda p: p[1] - p[0], pairwise([3, 10, 6]))
[7, -4]
\end{lstlisting}

Roll your own --- the definition:

\begin{lstlisting}[style=ex]
pairwise = lambda xs: list(zip(xs, xs[1:]))  # zip xs (tail xs)
\end{lstlisting}
\noindent{\small\itshape Used before it is rolled: \texttt{map\_} in \S\ref{def:map-} (p.~\pageref{def:map-}).}

\subsection{\texttt{isPrefixOf}}\label{def:isPrefixOf}
does xs start with p? Works on strings and lists alike (both are normalised through list()).
\begin{lstlisting}[style=ex]
>>> isPrefixOf("re", "rebuild")
True
>>> isPrefixOf([1, 2], [1, 2, 3])
True
>>> isPrefixOf("", "anything")
True
\end{lstlisting}

Roll your own --- the definition:

\begin{lstlisting}[style=ex]
isPrefixOf = lambda p, xs: list(xs[:len(p)]) == list(p)
\end{lstlisting}
\subsection{\texttt{isSuffixOf}}\label{def:isSuffixOf}
does xs end with p? Implemented by slicing from len(xs) - len(p), NOT by xs[-len(p):], because when p is empty that slice is xs[-0:] == the WHOLE list. A scar worth keeping visible.
\begin{lstlisting}[style=ex]
>>> isSuffixOf("ing", "folding")
True
>>> isSuffixOf("", "x")
True
>>> isSuffixOf("abc", "bc")
False
\end{lstlisting}

Roll your own --- the definition:

\begin{lstlisting}[style=ex]
isSuffixOf = lambda p, xs: list(xs[len(xs) - len(p):]) == list(p)
\end{lstlisting}
\subsection{\texttt{stripPrefix}}\label{def:stripPrefix}
Just the rest after the prefix, or Nothing: returns xs[len(p):] when the prefix matches, None when it does not. Pairs with fromMaybe and mapMaybe.
\begin{lstlisting}[style=ex]
>>> stripPrefix("re", "rebuild")
'build'
>>> stripPrefix("un", "rebuild") is None
True
>>> mapMaybe(partial(stripPrefix, "err:"), ["err:a", "ok:b", "err:c"])
['a', 'c']
\end{lstlisting}

Roll your own --- the definition:

\begin{lstlisting}[style=ex]
stripPrefix = lambda p, xs: xs[len(p):] if isPrefixOf(p, xs) else None  # Just the rest, or Nothing
\end{lstlisting}
\noindent{\small\itshape Used before it is rolled: \texttt{mapMaybe} in \S\ref{def:mapMaybe} (p.~\pageref{def:mapMaybe}), \texttt{partial} in \S\ref{def:partial} (p.~\pageref{def:partial}).}

\section{Homework}
Work in order: drills cement each tool, applies combine them,
challenges are interview-grade. Write each solution in a scratch file with
the given doctests pasted in, and let \texttt{python3 -m doctest} referee.
\subsection*{3.1\quad Swap the ends\hfill\textnormal{\textit{drill}}}
Return a copy of a list (length >= 2) with its first and last elements swapped and everything in between left alone. Build it from the end-taking tools, not by index.

\noindent\texttt{swap\_ends(xs: list) -> list}
\begin{lstlisting}[style=ex]

>>> swap_ends([1, 2, 3, 4])
[4, 2, 3, 1]
>>> swap_ends([1, 2])
[2, 1]
>>> swap_ends([5, 9, 9, 5])
[5, 9, 9, 5]

\end{lstlisting}
\subsection*{3.2\quad Deal the deck\hfill\textnormal{\textit{drill}}}
Deal a list into two halves. With an odd length the extra element goes to the SECOND half. Return the pair of halves.

\noindent\texttt{deal(xs: list) -> tuple}
\begin{lstlisting}[style=ex]

>>> deal([1, 2, 3, 4])
([1, 2], [3, 4])
>>> deal([1, 2, 3])
([1], [2, 3])
>>> deal([])
([], [])

\end{lstlisting}
\subsection*{3.3\quad Adjacent duplicate\hfill\textnormal{\textit{drill}}}
Report whether any two neighbouring elements of a list are equal. Empty and single-element lists have no neighbours, so the answer is False.

\noindent\texttt{has\_adjacent\_dup(xs: list) -> bool}
\begin{lstlisting}[style=ex]

>>> has_adjacent_dup([1, 2, 2, 3])
True
>>> has_adjacent_dup([1, 2, 3])
False
>>> has_adjacent_dup([])
False
>>> has_adjacent_dup([5])
False

\end{lstlisting}
\subsection*{3.4\quad Count the distinct\hfill\textnormal{\textit{drill}}}
Count how many distinct values a list contains. Order does not matter; duplicates count once.

\noindent\texttt{count\_distinct(xs: list) -> int}
\begin{lstlisting}[style=ex]

>>> count_distinct([1, 2, 2, 3, 1])
3
>>> count_distinct([])
0
>>> count_distinct([7, 7, 7])
1

\end{lstlisting}
\subsection*{3.5\quad Phone book lookup\hfill\textnormal{\textit{drill}}}
A phone book is a list of (name, number) pairs. Look up a name and return its number, or a supplied default string when the name is absent.

\noindent\texttt{lookup\_or(book: list, name, default) -> value}
\begin{lstlisting}[style=ex]

>>> book = [("ada", "123"), ("bob", "456")]
>>> lookup_or(book, "bob", "unknown")
'456'
>>> lookup_or(book, "cy", "unknown")
'unknown'
>>> lookup_or([], "ada", "unknown")
'unknown'

\end{lstlisting}
\subsection*{3.6\quad Spreadsheet column totals\hfill\textnormal{\textit{apply}}}
A spreadsheet is a list of rows, each row a list of numbers. Rows may be ragged (some shorter than others). Return the total of each column; a short row simply contributes nothing to columns it does not reach. An empty sheet has no columns.

\noindent\texttt{column\_sums(grid: list) -> list}
\begin{lstlisting}[style=ex]

>>> column_sums([[1, 2, 3], [4, 5, 6]])
[5, 7, 9]
>>> column_sums([[1, 2, 3], [4, 5]])
[5, 7, 3]
>>> column_sums([])
[]

\end{lstlisting}
\subsection*{3.7\quad Mastermind hits\hfill\textnormal{\textit{apply}}}
In Mastermind a guess and the secret code are equal-length sequences. Count the positions where the guess matches the secret exactly (the 'exact hits').

\noindent\texttt{matches(guess: list, secret: list) -> int}
\begin{lstlisting}[style=ex]

>>> matches([1, 2, 3, 4], [1, 0, 3, 0])
2
>>> matches("abc", "abc")
3
>>> matches([], [])
0

\end{lstlisting}
\subsection*{3.8\quad Strip a log tag\hfill\textnormal{\textit{apply}}}
Log lines sometimes begin with a fixed tag such as 'ERROR:'. Remove that leading tag if the line starts with it; otherwise return the line unchanged.

\noindent\texttt{strip\_tag(tag: str, line: str) -> str}
\begin{lstlisting}[style=ex]

>>> strip_tag("ERROR:", "ERROR:disk full")
'disk full'
>>> strip_tag("ERROR:", "all good")
'all good'
>>> strip_tag("", "unchanged")
'unchanged'

\end{lstlisting}
\subsection*{3.9\quad Filter by prefix and suffix\hfill\textnormal{\textit{apply}}}
Given a list of filenames, keep exactly those that both start with a required prefix and end with a required extension. Empty prefix or suffix matches everything.

\noindent\texttt{by\_affix(names: list, prefix: str, suffix: str) -> list}
\begin{lstlisting}[style=ex]

>>> by_affix(["test_a.py", "test_b.txt", "main.py"], "test_", ".py")
['test_a.py']
>>> by_affix(["a", "b"], "", "")
['a', 'b']
>>> by_affix(["main.py"], "test_", ".py")
[]

\end{lstlisting}
\subsection*{3.10\quad Pad a scoreboard\hfill\textnormal{\textit{apply}}}
A scoreboard must show exactly n slots. Given the current entries, pad on the right with a fill value until the list has length n. If it is already n or longer, leave it unchanged (never truncate).

\noindent\texttt{pad\_to(xs: list, n: int, fill) -> list}
\begin{lstlisting}[style=ex]

>>> pad_to([1, 2], 4, 0)
[1, 2, 0, 0]
>>> pad_to([1, 2, 3], 2, 0)
[1, 2, 3]
>>> pad_to([], 3, "-")
['-', '-', '-']

\end{lstlisting}
\subsection*{3.11\quad Roster reconciliation\hfill\textnormal{\textit{challenge}}}
An event keeps a list of previous attendees and a list of this time's check-ins (which may contain repeats). Produce three things: the DISTINCT current attendees who also attended before ('returning'), the distinct current attendees who did not ('newcomers'), and a flag that is True when there were no previous attendees at all. Preserve the first-seen order of the current check-ins in both lists.

\noindent\texttt{roster\_diff(prev: list, curr: list) -> tuple}
\begin{lstlisting}[style=ex]

>>> roster_diff(["ada", "bob"], ["bob", "cy", "cy", "ada", "dee"])
(['bob', 'ada'], ['cy', 'dee'], False)
>>> roster_diff([], ["ada", "ada", "bob"])
([], ['ada', 'bob'], True)
>>> roster_diff(["x"], [])
([], [], False)

\end{lstlisting}
\subsection*{3.12\quad Shipping manifest\hfill\textnormal{\textit{challenge}}}
A shipping manifest is built from three parallel arrays: item names, weights, and bin labels (equal length). Combine them into records, skip the first 'skip' records (already shipped), then number the remaining records from 1. Return the pair (numbered, total) where numbered is a list of (rank, name) and total is the summed weight of the kept records. Skipping past the end yields ([], 0).

\noindent\texttt{manifest(names: list, weights: list, bins: list, skip: int) -> tuple}
\begin{lstlisting}[style=ex]

>>> manifest(["a", "b", "c"], [10, 20, 30], ["x", "y", "z"], 1)
([(1, 'b'), (2, 'c')], 50)
>>> manifest(["a", "b", "c"], [10, 20, 30], ["x", "y", "z"], 0)
([(1, 'a'), (2, 'b'), (3, 'c')], 60)
>>> manifest(["a"], [10], ["x"], 5)
([], 0)

\end{lstlisting}
\clearpage
\chapter{Lists II: Higher-Order}

Where loops go to be named. \texttt{map\_} transforms, \texttt{filter\_}
selects, \texttt{concatMap} enumerates, \texttt{zipWith} combines --
and two tools do what loops do badly: \texttt{find} stops at the first hit
(even on an infinite stream), and \texttt{mapMaybe} runs a parser over messy
input and keeps only what parsed, one pass, no flags. This chapter is the
single highest-frequency vocabulary in the whole course.
\section{Concepts}
\subsection{\texttt{map\_}}\label{def:map-}
apply a function to every element; returns a real list, eagerly (the builtin map is lazy). The workhorse transformation.
\begin{lstlisting}[style=ex]
>>> map_(lambda x: x * x, [1, 2, 3])
[1, 4, 9]
>>> map_(str.upper, ["ab", "cd"])
['AB', 'CD']
\end{lstlisting}

Roll your own --- the definition:

\begin{lstlisting}[style=ex]
map_ = lambda f, xs: [f(x) for x in xs]
\end{lstlisting}
\subsection{\texttt{filter\_}}\label{def:filter-}
keep the elements satisfying a predicate. Scans the WHOLE list; if you only want the first hit, that is find.
\begin{lstlisting}[style=ex]
>>> filter_(even, [1, 2, 3, 4])
[2, 4]
>>> filter_(str.isdigit, "a1b22")
['1', '2', '2']
\end{lstlisting}

Roll your own --- the definition:

\begin{lstlisting}[style=ex]
filter_ = lambda pred, xs: [x for x in xs if pred(x)]
\end{lstlisting}
\subsection{\texttt{find}}\label{def:find}
the first element satisfying the predicate, else None. Lazy: it stops at the first hit, which folds cannot do, and therefore works even on infinite streams. Finish with fromMaybe if you need a default.
\begin{lstlisting}[style=ex]
>>> find(even, [1, 3, 4, 6])
4
>>> find(even, [1, 3]) is None
True
>>> find(lambda x: x % 7 == 0, count(1))
7
\end{lstlisting}

How it works: find wraps a generator expression in next() with a None default. The generator is the point -- elements are tested one at a time, on demand, and the search STOPS at the first hit. Nothing after the match is ever examined.


Why it earns its place: folds and filter\_ consume the whole list; they cannot stop early. find can -- which makes it the right tool the moment the word "first" appears in a statement, and the only tool that survives an infinite stream. Land defaults with fromMaybe.

\begin{lstlisting}[style=ex]
>>> find(lambda x: x * x > 50, count(1))
8
>>> fromMaybe('?', find(str.isdigit, "abc7de"))
'7'
\end{lstlisting}

Versus filter\_: filter\_ answers "which ones?"; find answers "who was first, if anyone?". Taking filter\_(p, xs)[0] instead does wasted work and explodes when nothing matches; find does neither.


Remember it as: find is the early exit that folds cannot perform.


Roll your own: a generator expression handed to next() with a default. The generator is what makes it lazy -- elements are produced one at a time, so the first hit stops everything -- and the default is what makes it total.

\begin{lstlisting}[style=ex]
find = lambda pred, xs: next((x for x in xs if pred(x)), None)  # first match, else None
\end{lstlisting}
\noindent{\small\itshape Used before it is rolled: \texttt{count} in \S\ref{def:count} (p.~\pageref{def:count}).}

\subsection{\texttt{partition}}\label{def:partition}
split into (satisfy, don't) in ONE pass, both sides in original order; the predicate runs exactly once per element, so it may be expensive or even stateful.
\begin{lstlisting}[style=ex]
>>> partition(even, [1, 2, 3, 4])
([2, 4], [1, 3])
>>> partition(str.isalpha, "a1b2")
(['a', 'b'], ['1', '2'])
\end{lstlisting}

Roll your own: two buckets and one conditional append -- (yes if pred(x) else no).append(x). The conditional expression picks the LIST, then the append mutates it; one pass, exactly one predicate call per element.

\begin{lstlisting}[style=ex]
def partition(pred, xs):  # (keepers, rest) in ONE pass; pred called once per element
    yes, no = [], []
    for x in xs:
        (yes if pred(x) else no).append(x)
    return (yes, no)
\end{lstlisting}
\subsection{\texttt{mapMaybe}}\label{def:mapMaybe}
map a function that returns a value or None, and keep only the values: parse-and-filter in ONE pass. The shape for messy input -- write a parser that answers None for garbage, then mapMaybe it over everything. dict.get is already such a function.
\begin{lstlisting}[style=ex]
>>> mapMaybe(lambda s: int(s) if s.isdigit() else None, ["3", "x", "7"])
[3, 7]
>>> mapMaybe({'a': 1, 'b': 2}.get, ['a', 'x', 'b'])
[1, 2]
\end{lstlisting}

How it works: run f over every element; f answers a VALUE when it can make sense of the input and None when it cannot; keep only the values. Map and filter fused into one pass, with the parser itself deciding what survives.

\begin{lstlisting}[style=ex]
>>> parse_int = lambda s: int(s) if s.lstrip('-').isdigit() else None
>>> mapMaybe(parse_int, ["12", "x", "-3", ""])
[12, -3]
\end{lstlisting}

Why it earns its place: messy input. The alternative is a loop with an if, a flag and an append -- three chances for a bug. Here the contract is crisp: write parse(item) -> value-or-None once, and mapMaybe handles every item. dict.get is already such a function, so lookups that may miss compose directly.

\begin{lstlisting}[style=ex]
>>> legend = {'a': 1, 'b': 2}
>>> mapMaybe(legend.get, "cabbage")
[1, 2, 2, 1]
\end{lstlisting}

The Haskell connection: Data.Maybe's mapMaybe, with None standing in for Nothing -- the "parse, don't validate" discipline: turn questionable data into clean data at the boundary, and everything downstream stops worrying.


Remember it as: map with a parser, drop the Nothings, one pass.


Roll your own --- the definition:

\begin{lstlisting}[style=ex]
mapMaybe = lambda f, xs: [y for y in (f(x) for x in xs) if y is not None]
\end{lstlisting}
\subsection{\texttt{catMaybes}}\label{def:catMaybes}
keep the non-None values; mapMaybe with the identity.
\begin{lstlisting}[style=ex]
>>> catMaybes([1, None, 2, None])
[1, 2]
\end{lstlisting}

Roll your own --- the definition:

\begin{lstlisting}[style=ex]
catMaybes = lambda xs: [x for x in xs if x is not None]  # mapMaybe id
\end{lstlisting}
\noindent{\small\itshape Used before it is rolled: \texttt{id} in \S\ref{def:id} (p.~\pageref{def:id}).}

\subsection{\texttt{concat}}\label{def:concat}
flatten ONE level of nesting.
\begin{lstlisting}[style=ex]
>>> concat([[1, 2], [3], []])
[1, 2, 3]
>>> concat(["ab", "cd"])
['a', 'b', 'c', 'd']
\end{lstlisting}

Roll your own --- the definition:

\begin{lstlisting}[style=ex]
concat = lambda xss: [x for xs in xss for x in xs]
\end{lstlisting}
\subsection{\texttt{concatMap}}\label{def:concatMap}
map each element to a LIST, and pool all the results: concat after map\_, fused. The backbone of enumeration -- "for every choice, produce all outcomes".
\begin{lstlisting}[style=ex]
>>> concatMap(lambda w: [w, w.upper()], ["a", "b"])
['a', 'A', 'b', 'B']
>>> concatMap(lambda x: [x] * x, [1, 2, 3])
[1, 2, 2, 3, 3, 3]
>>> concatMap(lambda x: [x, -x], [1, 2]) == concat(map_(lambda x: [x, -x], [1, 2]))
True
\end{lstlisting}

Roll your own --- the definition:

\begin{lstlisting}[style=ex]
concatMap = lambda f, xs: [y for x in xs for y in f(x)]
\end{lstlisting}
\subsection{\texttt{starmap}}\label{def:starmap}
map a MULTI-argument function over a list of argument tuples: map\_ composed with uncurry.
\begin{lstlisting}[style=ex]
>>> starmap(lambda a, b: a * b, [(2, 3), (4, 5)])
[6, 20]
>>> starmap(pow, [(2, 3), (3, 2)])
[8, 9]
\end{lstlisting}

Roll your own --- the definition:

\begin{lstlisting}[style=ex]
starmap = lambda f, pairs: [f(*p) for p in pairs]  # map . uncurry
\end{lstlisting}
\noindent{\small\itshape Used before it is rolled: \texttt{uncurry} in \S\ref{def:uncurry} (p.~\pageref{def:uncurry}).}

\subsection{\texttt{zipWith}}\label{def:zipWith}
combine two lists elementwise with a function; stops at the shorter. zip\_ is just zipWith(make-a-pair).
\begin{lstlisting}[style=ex]
>>> zipWith(add, [1, 2], [10, 20])
[11, 22]
>>> zipWith(lambda a, b: a > b, [3, 1], [2, 5, 9])
[True, False]
\end{lstlisting}

Roll your own --- the definition:

\begin{lstlisting}[style=ex]
zipWith = lambda f, a, b: [f(x, y) for x, y in zip(a, b)]
\end{lstlisting}
\subsection{\texttt{zipWith3}}\label{def:zipWith3}
the three-list version.
\begin{lstlisting}[style=ex]
>>> zipWith3(lambda a, b, c: a + b + c, [1], [2], [3])
[6]
\end{lstlisting}

Roll your own --- the definition:

\begin{lstlisting}[style=ex]
zipWith3 = lambda f, a, b, c: [f(x, y, z) for x, y, z in zip(a, b, c)]
\end{lstlisting}
\subsection{\texttt{cross}}\label{def:cross}
every pairing of two lists (the cartesian product). For the product of a list with ITSELF n times, see replicateM.
\begin{lstlisting}[style=ex]
>>> cross([1, 2], "ab")
[(1, 'a'), (1, 'b'), (2, 'a'), (2, 'b')]
>>> len(cross(range(3), range(4)))
12
\end{lstlisting}

Roll your own --- the definition:

\begin{lstlisting}[style=ex]
cross = lambda a, b: [(x, y) for x in a for y in b]  # cartesian; `product` is the fold
\end{lstlisting}
\noindent{\small\itshape Used before it is rolled: \texttt{product} in \S\ref{def:product} (p.~\pageref{def:product}).}

\section{Homework}
Work in order: drills cement each tool, applies combine them,
challenges are interview-grade. Write each solution in a scratch file with
the given doctests pasted in, and let \texttt{python3 -m doctest} referee.
\subsection*{4.1\quad Squares of the evens\hfill\textnormal{\textit{drill}}}
Given a list of integers, return the squares of just the even ones, in order. This is the filter-then-map pipeline in its smallest form.

\noindent\texttt{squares\_of\_evens(xs: list) -> list}
\begin{lstlisting}[style=ex]

>>> squares_of_evens([1, 2, 3, 4])
[4, 16]
>>> squares_of_evens([1, 3, 5])
[]
>>> squares_of_evens([])
[]

\end{lstlisting}
\subsection*{4.2\quad First over the line\hfill\textnormal{\textit{drill}}}
Return the first element of a list strictly greater than a threshold, or None if no element qualifies. Stop at the first hit -- do not scan the whole list.

\noindent\texttt{first\_over(xs: list, t) -> value | None}
\begin{lstlisting}[style=ex]

>>> first_over([1, 5, 3, 9], 4)
5
>>> first_over([1, 2, 3], 9) is None
True
>>> first_over([], 0) is None
True

\end{lstlisting}
\subsection*{4.3\quad Pass / fail split\hfill\textnormal{\textit{drill}}}
Split exam scores into (passed, failed) against a cutoff, where a score >= cutoff passes. Keep both groups in their original order and make a single pass.

\noindent\texttt{split\_pass\_fail(scores: list, cutoff: int) -> tuple}
\begin{lstlisting}[style=ex]

>>> split_pass_fail([50, 80, 60, 90], 70)
([80, 90], [50, 60])
>>> split_pass_fail([], 70)
([], [])
>>> split_pass_fail([100, 100], 50)
([100, 100], [])

\end{lstlisting}
\subsection*{4.4\quad Salvage the integers\hfill\textnormal{\textit{drill}}}
Given a list of string tokens, return the integers among them, in order, dropping any token that is not a (possibly negative) integer. Empty strings are dropped.

\noindent\texttt{clean\_ints(tokens: list) -> list}
\begin{lstlisting}[style=ex]

>>> clean_ints(["12", "x", "-3", ""])
[12, -3]
>>> clean_ints(["0", "00"])
[0, 0]
>>> clean_ints(["nope"])
[]

\end{lstlisting}
\subsection*{4.5\quad Stutter\hfill\textnormal{\textit{drill}}}
Return a new list in which every element of the input appears twice in a row, in order.

\noindent\texttt{duplicate\_each(xs: list) -> list}
\begin{lstlisting}[style=ex]

>>> duplicate_each([1, 2])
[1, 1, 2, 2]
>>> duplicate_each([])
[]
>>> duplicate_each(["a"])
['a', 'a']

\end{lstlisting}
\subsection*{4.6\quad Clean the sensor feed\hfill\textnormal{\textit{apply}}}
A sensor feed is a list of readings in which dropped samples appear as None. Return the pair (good, dropped): the list of surviving readings in order, and how many were None.

\noindent\texttt{clean\_readings(readings: list) -> tuple}
\begin{lstlisting}[style=ex]

>>> clean_readings([1, None, 2, None, 3])
([1, 2, 3], 2)
>>> clean_readings([None, None])
([], 2)
>>> clean_readings([4, 5])
([4, 5], 0)

\end{lstlisting}
\subsection*{4.7\quad Line-item totals\hfill\textnormal{\textit{apply}}}
An order has parallel arrays of unit prices and quantities. Return the total for each line item (price times quantity), position by position. If the arrays differ in length, the shorter one governs.

\noindent\texttt{line\_totals(prices: list, quantities: list) -> list}
\begin{lstlisting}[style=ex]

>>> line_totals([2, 3, 5], [4, 1, 2])
[8, 3, 10]
>>> line_totals([2, 3, 5], [4])
[8]
>>> line_totals([], [])
[]

\end{lstlisting}
\subsection*{4.8\quad Three-column totals\hfill\textnormal{\textit{apply}}}
An invoice has three parallel columns: base price, tax, and tip. Return the grand total for each row by summing the three columns position by position.

\noindent\texttt{row\_totals(base: list, tax: list, tip: list) -> list}
\begin{lstlisting}[style=ex]

>>> row_totals([10, 20], [1, 2], [2, 3])
[13, 25]
>>> row_totals([10], [1], [2])
[13]
>>> row_totals([], [], [])
[]

\end{lstlisting}
\subsection*{4.9\quad Rectangle areas\hfill\textnormal{\textit{apply}}}
Given a list of rectangles, each a (width, height) pair, return the area of each rectangle in order.

\noindent\texttt{areas(rects: list) -> list}
\begin{lstlisting}[style=ex]

>>> areas([(2, 3), (4, 5)])
[6, 20]
>>> areas([(0, 9)])
[0]
>>> areas([])
[]

\end{lstlisting}
\subsection*{4.10\quad Merge the playlists\hfill\textnormal{\textit{apply}}}
Given several playlists, each a list of song titles, merge them into one playlist in order, then remove duplicate songs keeping the first appearance of each.

\noindent\texttt{merge\_playlists(playlists: list) -> list}
\begin{lstlisting}[style=ex]

>>> merge_playlists([["a", "b"], ["b", "c"], ["a"]])
['a', 'b', 'c']
>>> merge_playlists([])
[]
>>> merge_playlists([[], ["solo"]])
['solo']

\end{lstlisting}
\subsection*{4.11\quad Affordable outfits\hfill\textnormal{\textit{challenge}}}
A wardrobe has shirts and pants, each an (name, price) pair. Generate every shirt-and-pant outfit, keep only those whose combined price is within a budget, and return them as (shirt\_name, pant\_name) pairs in generation order (shirts outer, pants inner).

\noindent\texttt{affordable\_outfits(shirts: list, pants: list, budget: int) -> list}
\begin{lstlisting}[style=ex]

>>> shirts = [("red", 20), ("blue", 30)]
>>> pants = [("a", 40), ("b", 10)]
>>> affordable_outfits(shirts, pants, 45)
[('red', 'b'), ('blue', 'b')]
>>> affordable_outfits(shirts, pants, 5)
[]
>>> affordable_outfits([], pants, 100)
[]

\end{lstlisting}
\subsection*{4.12\quad Inventory reconciliation\hfill\textnormal{\textit{challenge}}}
An inventory log is a list of transaction strings. Each is an item name optionally led by a sign: '+apple' or 'apple' is an addition, '-banana' is a removal. Blank strings are skipped. Parse the log, then return the pair (added, removed): the item names that were added and the item names that were removed, each in the order they appeared in the log.

\noindent\texttt{reconcile(txns: list) -> tuple}
\begin{lstlisting}[style=ex]

>>> reconcile(["+apple", "-banana", "cherry", "", "-apple"])
(['apple', 'cherry'], ['banana', 'apple'])
>>> reconcile(["", "  "])
([], [])
>>> reconcile(["milk", "-milk"])
(['milk'], ['milk'])

\end{lstlisting}
\clearpage
\chapter{The Fold}

The trunk of the tree. A fold is a for-loop whose state has been given a
contract: an accumulator, a combiner, one pass. Everything after this chapter
is a fold in costume -- scans keep the fold's history, unfolds run it
backwards, dict folds aim it at a mapping, \texttt{memo} folds a recursion
into a cache. Learn the two argument disciplines cold: \texttt{foldl}'s
combiner takes \texttt{(acc, x)}, \texttt{foldr}'s takes \texttt{(x, acc)}.
The enumeration folds close the chapter with the brute-force licenses:
$2^n$ subsets and $|xs|^n$ words, legal when you say the bound out loud.
\section{Concepts}
\subsection{\texttt{foldl}}\label{def:foldl}
THE left fold, and the prelude's most important definition. It is exactly this loop: acc = init; for x in xs: acc = f(acc, x); return acc. The combiner takes (ACCUMULATOR, element) -- accumulator first. Argument order of foldl itself: (f, xs, init), with init optional -- omitted, the first element seeds the accumulator (and an empty list is then an error).
\begin{lstlisting}[style=ex]
>>> foldl(lambda acc, x: acc * 10 + x, [1, 2, 3], 0)
123
>>> foldl(max, [3, 9, 2])               # no init: 3 seeds it
9
>>> foldl(add, [], 0)                   # empty + init: the init comes back
0
>>> foldl(lambda acc, w: acc + len(w), ["ab", "c"], 0)
3
\end{lstlisting}

How it works: foldl is precisely this loop, given a name and a contract:

\begin{lstlisting}[style=ex]
>>> def foldl_spelled_out(f, xs, init):
...     acc = init
...     for x in xs:
...         acc = f(acc, x)
...     return acc
>>> foldl_spelled_out(add, [1, 2, 3], 0) == foldl(add, [1, 2, 3], 0)
True
\end{lstlisting}

Two disciplines hide in the signature. The combiner takes (ACCUMULATOR, element) -- accumulator first, always; reversing them is the classic fold bug. And the optional init is guarded by NOTHING, so the fold can tell "no seed -- use the first element" from "the seed is None".


Why it earns its place: the fold is the universal list consumer. sum, max, reverse, "build a dict", "run a state machine over events" -- each is one choice of combiner and seed. When you can say "carry THIS state, update it THIS way per element", the loop is already written.

\begin{lstlisting}[style=ex]
>>> foldl(lambda acc, x: [x] + acc, [1, 2, 3], [])       # reverse is a fold
[3, 2, 1]
>>> foldl(lambda acc, w: acc + len(w), ["ab", "cde"], 0)
5
\end{lstlisting}

Remember it as: a fold is a for-loop whose state has been given a contract.


Roll your own: seed the accumulator -- init if one was given, else the first element (the NOTHING guard is what lets None be a legitimate seed) -- then a single for-loop folding acc = f(acc, x). The entire contract is which side the accumulator sits on.

\begin{lstlisting}[style=ex]
def foldl(f, xs, init=NOTHING):  # THE left fold; Python buried its own in functools as reduce
    it = iter(xs)
    if init is NOTHING:
        try:
            acc = next(it)
        except StopIteration:
            raise TypeError("fold of empty sequence with no initial value")
    else:
        acc = init
    for x in it:
        acc = f(acc, x)
    return acc
\end{lstlisting}
\subsection{\texttt{foldl1}}\label{def:foldl1}
foldl seeded by the first element, spelled out. Use when the accumulator has the same type as the elements and emptiness cannot happen.
\begin{lstlisting}[style=ex]
>>> foldl1(add, [1, 2, 3])
6
>>> foldl1(min, [3, 1, 2])
1
\end{lstlisting}

Roll your own --- the definition:

\begin{lstlisting}[style=ex]
foldl1 = lambda f, xs: foldl(f, xs)
\end{lstlisting}
\subsection{\texttt{foldr}}\label{def:foldr}
the RIGHT fold: nests from the right, f takes (element, ACCUMULATOR) -- element first, mirror of foldl. Implemented as foldl of the flipped function over the reversed list: iterative and stack-safe, an identity that only holds because Python is strict. The parenthesisation example makes the direction visible.
\begin{lstlisting}[style=ex]
>>> foldr(lambda x, acc: '(' + x + '+' + acc + ')', "abc", "z")
'(a+(b+(c+z)))'
>>> foldl(lambda acc, x: '(' + acc + '+' + x + ')', "abc", "z")
'(((z+a)+b)+c)'
>>> foldr(lambda x, acc: [x] + acc, [1, 2, 3], [])
[1, 2, 3]
\end{lstlisting}

How it works: the mirror image -- the combiner takes (element, ACCUMULATOR) and the nesting grows from the RIGHT: f(a, f(b, f(c, z))). The prelude does not write a second recursion: foldr = foldl(flip(f), reversed(xs), init). Flip fixes the argument order, reversing fixes the traversal, and the existing engine runs backwards.


The honesty clause: that identity only holds because Python is STRICT. Haskell's foldr can process infinite lists -- laziness lets f decide whether the rest is ever needed; ours materialises and reverses the list first, so it cannot. Same name, same finite behaviour, different superpower.


Why it earns its place: building right-to-left. Consing onto a front is natural from the right -- section 13's from\_list IS a foldr of ListNode, and here it is, derived by hand:

\begin{lstlisting}[style=ex]
>>> foldr(lambda x, acc: [x] + acc, [1, 2, 3], [])       # rebuild, order kept
[1, 2, 3]
>>> to_list(foldr(ListNode, [1, 2, 3], None))            # from_list, spelled out
[1, 2, 3]
\end{lstlisting}

Remember it as: foldl eats left-to-right; foldr thinks right-to-left -- and in strict Python it is foldl wearing flip and a reversal.


Roll your own --- the definition:

\begin{lstlisting}[style=ex]
foldr = lambda f, xs, init: foldl(flip(f), reversed(list(xs)), init)
\end{lstlisting}
\noindent{\small\itshape Used before it is rolled: \texttt{ListNode} in \S\ref{def:ListNode} (p.~\pageref{def:ListNode}), \texttt{from\_list} in \S\ref{def:from-list} (p.~\pageref{def:from-list}), \texttt{to\_list} in \S\ref{def:to-list} (p.~\pageref{def:to-list}), \texttt{flip} in \S\ref{def:flip} (p.~\pageref{def:flip}).}

\subsection{\texttt{foldr1}}\label{def:foldr1}
right fold seeded by the LAST element.
\begin{lstlisting}[style=ex]
>>> foldr1(lambda x, acc: x - acc, [1, 2, 3])   # 1 - (2 - 3)
2
\end{lstlisting}

Roll your own --- the definition:

\begin{lstlisting}[style=ex]
foldr1 = lambda f, xs: foldl(flip(f), reversed(list(xs)))
\end{lstlisting}
\noindent{\small\itshape Used before it is rolled: \texttt{flip} in \S\ref{def:flip} (p.~\pageref{def:flip}).}

\subsection{\texttt{product}}\label{def:product}
multiply everything: the numeric fold, seeded with 1 so an empty product is 1.
\begin{lstlisting}[style=ex]
>>> product([2, 3, 4])
24
>>> product([])
1
\end{lstlisting}

Roll your own --- the definition:

\begin{lstlisting}[style=ex]
product = lambda xs: foldl(lambda a, x: a * x, xs, 1)  # the numeric fold
\end{lstlisting}
\subsection{\texttt{subsequences}}\label{def:subsequences}
every subset of the elements (the powerset), built by a fold in which each element DOUBLES the accumulator: keep every existing subset, and every existing subset extended by the newcomer. 2\^{}n results -- a brute-force license for n up to about 20, and you should say that bound out loud before using it.
\begin{lstlisting}[style=ex]
>>> subsequences([1, 2])
[[], [1], [2], [1, 2]]
>>> subsequences("ab")
[[], ['a'], ['b'], ['a', 'b']]
>>> len(subsequences("abcd"))
16
\end{lstlisting}

How it works: watch the accumulator double. Start from [[]] -- one subset, the empty one. Each element x rewrites acc as acc + [s + [x] for s in acc]: every existing subset survives (x left out), and every existing subset reappears extended by x (x taken). n elements, n doublings, 2\^{}n subsets, in a stable order.

\begin{lstlisting}[style=ex]
>>> acc = [[]]
>>> acc = acc + [s + [1] for s in acc]; acc
[[], [1]]
>>> acc = acc + [s + [2] for s in acc]; acc
[[], [1], [2], [1, 2]]
\end{lstlisting}

Why it earns its place: it is a brute-force LICENSE. At n <= 20, checking every subset is a few million cheap operations -- honest, exhaustive, and unbeatable to debug. Say the bound out loud ("2\^{}15 is 32768 -- fine") and enumerate; cleverness can come later if the bound grows.

\begin{lstlisting}[style=ex]
>>> best = maxOn(sum, [s for s in subsequences([3, -1, 4, -2]) if sum(s) <= 6])
>>> sum(best)
6
\end{lstlisting}

Remember it as: the powerset fold -- each element doubles the world.


Roll your own --- the definition:

\begin{lstlisting}[style=ex]
subsequences = lambda xs: foldl(lambda acc, x: acc + [s + [x] for s in acc], list(xs), [[]])
\end{lstlisting}
\noindent{\small\itshape Used before it is rolled: \texttt{maxOn} in \S\ref{def:maxOn} (p.~\pageref{def:maxOn}).}

\subsection{\texttt{replicateM}}\label{def:replicateM}
every length-n word over an alphabet: cross applied n times, |xs|\^{}n results. The other brute-force license -- all bitstrings, all dice rolls, all assignments of n slots.
\begin{lstlisting}[style=ex]
>>> replicateM(2, "ab")
[['a', 'a'], ['a', 'b'], ['b', 'a'], ['b', 'b']]
>>> len(replicateM(3, [0, 1]))
8
\end{lstlisting}

Roll your own --- the definition:

\begin{lstlisting}[style=ex]
replicateM = lambda n, xs: foldl(lambda acc, _: [s + [x] for s in acc for x in xs], range(n), [[]])
\end{lstlisting}
\section{Homework}
Work in order: drills cement each tool, applies combine them,
challenges are interview-grade. Write each solution in a scratch file with
the given doctests pasted in, and let \texttt{python3 -m doctest} referee.
\subsection*{5.1\quad Sum of squares\hfill\textnormal{\textit{drill}}}
Using foldl, compute the sum of the squares of a list of integers. The accumulator carries the running total; seed it with 0 so the empty list answers 0.

\noindent\texttt{sum\_squares(xs: list[int]) -> int}
\begin{lstlisting}[style=ex]
>>> sum_squares([1, 2, 3])
14
>>> sum_squares([])
0
>>> sum_squares([-2])
4
\end{lstlisting}
\subsection*{5.2\quad Longest word, first on ties\hfill\textnormal{\textit{drill}}}
Using foldl1, return the longest word in a non-empty list. When two words tie in length, the EARLIER one must win -- choose your comparison so the accumulator survives ties.

\noindent\texttt{longest\_word(ws: list[str]) -> str}
\begin{lstlisting}[style=ex]
>>> longest_word(["hi", "abc", "de"])
'abc'
>>> longest_word(["aa", "bb"])
'aa'
>>> longest_word(["x"])
'x'
\end{lstlisting}
\subsection*{5.3\quad Cons cells by hand\hfill\textnormal{\textit{drill}}}
Using foldr, convert a list into nested pairs ending in None, Lisp-style: [1, 2, 3] becomes (1, (2, (3, None))). The nesting grows from the right, which is exactly what foldr does; note its combiner receives (element, accumulator).

\noindent\texttt{to\_pairs(xs: list) -> tuple | None}
\begin{lstlisting}[style=ex]
>>> to_pairs([1, 2, 3])
(1, (2, (3, None)))
>>> to_pairs(['a'])
('a', None)
>>> to_pairs([]) is None
True
\end{lstlisting}
\subsection*{5.4\quad Factorial\hfill\textnormal{\textit{drill}}}
Using product, compute n factorial. Remember product of an empty sequence is 1 -- which is precisely why 0! works with no special case.

\noindent\texttt{factorial(n: int) -> int}
\begin{lstlisting}[style=ex]
>>> factorial(5)
120
>>> factorial(0)
1
>>> factorial(1)
1
\end{lstlisting}
\subsection*{5.5\quad Every substring team name\hfill\textnormal{\textit{drill}}}
Using subsequences, list all subsets of s's characters as strings -- every string formable by deleting zero or more characters from s (keeping order): powerset\_strings("ab") is ['', 'a', 'b', 'ab']. Join each subsequence back into a string; keep subsequences' own order.

\noindent\texttt{powerset\_strings(s: str) -> list[str]}
\begin{lstlisting}[style=ex]
>>> powerset_strings("ab")
['', 'a', 'b', 'ab']
>>> powerset_strings("")
['']
>>> len(powerset_strings("abc"))
8
\end{lstlisting}
\subsection*{5.6\quad Bank statement audit\hfill\textnormal{\textit{apply}}}
A bank account starts at 'start' (never negative) and processes a list of signed transactions in order. Return (final\_balance, overdrawn) where overdrawn is True if the running balance ever dipped below zero at any point -- even if later deposits recovered it. One fold: the accumulator carries both the balance and the flag.

\noindent\texttt{balance\_audit(start: int, txns: list[int]) -> tuple[int, bool]}
\begin{lstlisting}[style=ex]
>>> balance_audit(100, [-150, 75])
(25, True)
>>> balance_audit(100, [])
(100, False)
>>> balance_audit(0, [-1, 1])
(0, True)
>>> balance_audit(10, [5])
(15, False)
\end{lstlisting}
\subsection*{5.7\quad Ways to roll it\hfill\textnormal{\textit{apply}}}
You roll n six-sided dice. Using replicateM, count every way the ordered dice can sum to exactly target. n is small (n <= 6), so enumerating all 6\^{}n outcomes is the honest play.

\noindent\texttt{ways\_to\_roll(n: int, target: int) -> int}
\begin{lstlisting}[style=ex]
>>> ways_to_roll(2, 7)
6
>>> ways_to_roll(1, 3)
1
>>> ways_to_roll(2, 13)
0
>>> ways_to_roll(3, 18)
1
\end{lstlisting}
\subsection*{5.8\quad Decompress the telemetry\hfill\textnormal{\textit{apply}}}
A sensor stream arrives run-length encoded as (char, count) pairs. Using foldl, decompress it back into the original string: [('a', 3), ('b', 2)] becomes 'aaabb'.

\noindent\texttt{rle\_decode(pairs: list[tuple[str, int]]) -> str}
\begin{lstlisting}[style=ex]
>>> rle_decode([('a', 3), ('b', 2)])
'aaabb'
>>> rle_decode([])
''
>>> rle_decode([('x', 1)])
'x'
\end{lstlisting}
\subsection*{5.9\quad Config templating\hfill\textnormal{\textit{apply}}}
A deployment file is produced by applying find-and-replace edits IN ORDER: each edit is a (find, replace) pair, and later edits see the result of earlier ones. Using foldl with the text as the accumulator, apply them all.

\noindent\texttt{apply\_edits(text: str, edits: list[tuple[str, str]]) -> str}
\begin{lstlisting}[style=ex]
>>> apply_edits("hello world", [("world", "there"), ("hello", "hi")])
'hi there'
>>> apply_edits("aa", [("a", "b")])
'bb'
>>> apply_edits("cfg:$X", [("$X", "1")])
'cfg:1'
>>> apply_edits("keep", [])
'keep'
\end{lstlisting}
\subsection*{5.10\quad Exact change\hfill\textnormal{\textit{apply}}}
You hold a small pouch of coins, each usable at most once (n <= 20). Using subsequences, decide whether there is any way to pick a subset of them summing to exactly target. The empty subset makes 0, so target 0 is always reachable.

\noindent\texttt{exact\_change(coins: list[int], target: int) -> bool}
\begin{lstlisting}[style=ex]
>>> exact_change([3, 7, 2], 9)
True
>>> exact_change([5, 10], 7)
False
>>> exact_change([], 0)
True
>>> exact_change([4], 0)
True
\end{lstlisting}
\subsection*{5.11\quad Minimal hiring committee\hfill\textnormal{\textit{challenge}}}
A project needs a set of skills. Each candidate brings a list of skills. Hiring is expensive, so find the SIZE of the smallest group of candidates that together covers every needed skill, or -1 if no group can. Candidate count is small (n <= 12): enumerate groups with subsequences, keep the covering ones, minimise the size. An empty requirement needs zero hires.

\noindent\texttt{min\_team(required: list[str], candidates: list[list[str]]) -> int}
\begin{lstlisting}[style=ex]
>>> min_team(["py", "sql"], [["py"], ["sql"], ["py", "sql"]])
1
>>> min_team(["a", "b"], [["a"], ["b"]])
2
>>> min_team(["go"], [["py"]])
-1
>>> min_team([], [["x"]])
0
\end{lstlisting}
\subsection*{5.12\quad Fair loot split\hfill\textnormal{\textit{challenge}}}
Two players split a pile of treasure items (values, n <= 20) into two groups; every item goes to one side. Return the smallest possible absolute difference between the two groups' totals. Observe that choosing one side's subset s fixes the difference to abs(total - 2*sum(s)), then let subsequences try every side.

\noindent\texttt{min\_partition\_diff(xs: list[int]) -> int}
\begin{lstlisting}[style=ex]
>>> min_partition_diff([1, 6, 11, 5])
1
>>> min_partition_diff([3, 1])
2
>>> min_partition_diff([2, 2])
0
>>> min_partition_diff([])
0
>>> min_partition_diff([7])
7
\end{lstlisting}
\subsection*{5.13\quad Next greater reading\hfill\textnormal{\textit{challenge}}}
Daily sensor readings arrive as a list. For each position report the NEXT GREATER reading to its right -- the first later value strictly larger -- or -1 when no later reading beats it. A double loop is O(n\^{}2); a list used as a stack of still-unanswered positions, driven by one left-to-right fold, answers everything in O(n).

\noindent\texttt{next\_greater(xs: list[int]) -> list[int]}
\begin{lstlisting}[style=ex]
>>> next_greater([2, 1, 3])
[3, 3, -1]
>>> next_greater([1, 2, 3, 4])
[2, 3, 4, -1]
>>> next_greater([4, 3, 2, 1])
[-1, -1, -1, -1]
>>> next_greater([2, 2])
[-1, -1]
>>> next_greater([])
[]
\end{lstlisting}
\clearpage
\chapter{Scans}

A scan is a fold that keeps its history -- every intermediate accumulator,
not just the last. That one change turns ``the state at every step'' from n
folds into one pass, makes prefix sums (and with them $O(1)$ range queries)
a single line, and reveals celebrated one-pass algorithms -- Kadane's most
famously -- as scans in a trench coat. When a statement says ``running'',
``so far'' or ``at each step'', it is dictating this chapter.
\section{Concepts}
\subsection{\texttt{accumulate}}\label{def:accumulate}
the generator behind the scans: yields the running combination of everything seen so far. With no function it sums; 'initial' seeds it. This is Python's itertools name with Haskell's scan semantics.
\begin{lstlisting}[style=ex]
>>> list(accumulate([1, 2, 3]))
[1, 3, 6]
>>> list(accumulate([1, 2], initial=10))
[10, 11, 13]
>>> list(accumulate([3, 1, 2], min))
[3, 1, 1]
\end{lstlisting}

Roll your own: the same skeleton as foldl, but a GENERATOR -- yield the accumulator once before the loop, then yield again after every combine. The early bare return on an empty, seedless input is what makes an empty scan come back empty instead of raising.

\begin{lstlisting}[style=ex]
def accumulate(xs, f=None, initial=NOTHING):  # scanl / scanl1
    if f is None:
        f = lambda a, b: a + b
    it = iter(xs)
    if initial is NOTHING:
        try:
            acc = next(it)
        except StopIteration:
            return
    else:
        acc = initial
    yield acc
    for x in it:
        acc = f(acc, x)
        yield acc
\end{lstlisting}
\noindent{\small\itshape Used before it is rolled: \texttt{scanl} in \S\ref{def:scanl} (p.~\pageref{def:scanl}), \texttt{scanl1} in \S\ref{def:scanl1} (p.~\pageref{def:scanl1}).}

\subsection{\texttt{scanl}}\label{def:scanl}
a fold that KEEPS ITS HISTORY: returns every intermediate accumulator, starting with the seed, so the result has len(xs)+1 elements and its last element equals the foldl. Argument order (f, z, xs) -- seed in the middle, Haskell's order, unlike foldl's trailing init. Prefix sums are the canonical scan.
\begin{lstlisting}[style=ex]
>>> scanl(add, 0, [3, 1, 4])
[0, 3, 4, 8]
>>> last(scanl(add, 0, [1, 2, 3])) == foldl(add, [1, 2, 3], 0)
True
\end{lstlisting}

The DP secret: many celebrated one-pass algorithms are scans in a trench coat, because a scan IS a dynamic-programming table -- entry i holds the answer for the prefix ending at i. Kadane's famous maximum-subarray "trick" is one scanl1 and a max:

\begin{lstlisting}[style=ex]
>>> xs = [-2, 1, -3, 4, -1, 2, 1, -5, 4]
>>> max(scanl1(lambda best, x: max(x, best + x), xs))
6
\end{lstlisting}

Prefix sums, the other superpower: scan once with (+) and every range sum becomes two reads, pref[j] - pref[i] -- O(n) window questions collapse to O(1).

\begin{lstlisting}[style=ex]
>>> pref = scanl(add, 0, [3, 1, 4, 1, 5])
>>> pref[4] - pref[1]                     # sum of elements 1..3
6
\end{lstlisting}

Remember it as: when you need the state at EVERY step, one scan replaces n folds.


Roll your own --- the definition:

\begin{lstlisting}[style=ex]
scanl = lambda f, z, xs: list(accumulate(xs, f, initial=z))
\end{lstlisting}
\noindent{\small\itshape Used before it is rolled: \texttt{scanl1} in \S\ref{def:scanl1} (p.~\pageref{def:scanl1}).}

\subsection{\texttt{scanl1}}\label{def:scanl1}
the history-keeping fold seeded by the first element; len(xs) results. Running maxima and minima live here.
\begin{lstlisting}[style=ex]
>>> scanl1(min, [3, 1, 2])
[3, 1, 1]
>>> scanl1(max, [1, 3, 2, 5])
[1, 3, 3, 5]
\end{lstlisting}

Roll your own --- the definition:

\begin{lstlisting}[style=ex]
scanl1 = lambda f, xs: list(accumulate(xs, f))
\end{lstlisting}
\subsection{\texttt{scanr}}\label{def:scanr}
the suffix scan: folds from the RIGHT, showing every suffix's result; the first element is the full fold, the last is the seed. Suffix sums for product\_except\_self-style problems.
\begin{lstlisting}[style=ex]
>>> scanr(add, 0, [1, 2, 3])
[6, 5, 3, 0]
\end{lstlisting}

Roll your own --- the definition:

\begin{lstlisting}[style=ex]
scanr = lambda f, z, xs: list(reversed(scanl(flip(f), z, list(reversed(list(xs))))))
\end{lstlisting}
\noindent{\small\itshape Used before it is rolled: \texttt{flip} in \S\ref{def:flip} (p.~\pageref{def:flip}).}

\subsection{\texttt{scanr1}}\label{def:scanr1}
suffix scan seeded by the last element.
\begin{lstlisting}[style=ex]
>>> scanr1(add, [1, 2, 3])
[6, 5, 3]
\end{lstlisting}

Roll your own --- the definition:

\begin{lstlisting}[style=ex]
scanr1 = lambda f, xs: list(reversed(scanl1(flip(f), list(reversed(list(xs))))))
\end{lstlisting}
\noindent{\small\itshape Used before it is rolled: \texttt{flip} in \S\ref{def:flip} (p.~\pageref{def:flip}).}

\section{Homework}
Work in order: drills cement each tool, applies combine them,
challenges are interview-grade. Write each solution in a scratch file with
the given doctests pasted in, and let \texttt{python3 -m doctest} referee.
\subsection*{6.1\quad Running maximum\hfill\textnormal{\textit{drill}}}
Using scanl1, produce the running maximum of a list: element i of the answer is the largest value seen among the first i+1 inputs. An empty list yields an empty scan.

\noindent\texttt{running\_max(xs: list[int]) -> list[int]}
\begin{lstlisting}[style=ex]
>>> running_max([3, 1, 4, 1, 5])
[3, 3, 4, 4, 5]
>>> running_max([7])
[7]
>>> running_max([])
[]
\end{lstlisting}
\subsection*{6.2\quad Account history\hfill\textnormal{\textit{drill}}}
An account opens at balance 'start' and applies signed transactions in order. Using scanl, return the FULL balance history -- the balance at each step, opening balance first -- len(txns) + 1 entries.

\noindent\texttt{account\_history(start: int, txns: list[int]) -> list[int]}
\begin{lstlisting}[style=ex]
>>> account_history(100, [-20, 50])
[100, 80, 130]
>>> account_history(0, [])
[0]
>>> account_history(5, [-10])
[5, -5]
\end{lstlisting}
\subsection*{6.3\quad Suffix sums\hfill\textnormal{\textit{drill}}}
Using scanr, compute suffix sums: element i is the sum of xs[i:], and the final element is the seed 0 (the empty suffix). [1, 2, 3] becomes [6, 5, 3, 0].

\noindent\texttt{suffix\_sums(xs: list[int]) -> list[int]}
\begin{lstlisting}[style=ex]
>>> suffix_sums([1, 2, 3])
[6, 5, 3, 0]
>>> suffix_sums([])
[0]
>>> suffix_sums([5])
[5, 0]
\end{lstlisting}
\subsection*{6.4\quad Best price still ahead\hfill\textnormal{\textit{drill}}}
Shopping day: prices[i] is today's price at shop i, and you may only buy at the CURRENT shop or a later one. Using scanr1, compute for each position the cheapest price from there to the end of the street.

\noindent\texttt{best\_remaining(prices: list[int]) -> list[int]}
\begin{lstlisting}[style=ex]
>>> best_remaining([5, 3, 8, 2, 4])
[2, 2, 2, 2, 4]
>>> best_remaining([7])
[7]
>>> best_remaining([4, 6])
[4, 6]
\end{lstlisting}
\subsection*{6.5\quad Running product\hfill\textnormal{\textit{drill}}}
Using accumulate with a custom combiner, produce the running product of a list: [2, 3, 4] becomes [2, 6, 24]. Remember accumulate is a generator -- deliver a list.

\noindent\texttt{running\_product(xs: list[int]) -> list[int]}
\begin{lstlisting}[style=ex]
>>> running_product([2, 3, 4])
[2, 6, 24]
>>> running_product([])
[]
>>> running_product([5])
[5]
\end{lstlisting}
\subsection*{6.6\quad Record-breaking days\hfill\textnormal{\textit{apply}}}
A weather station logs one temperature per day. A day sets a RECORD if it is strictly hotter than every previous day; the first day always counts. Count the record days. Hint: on the running-maximum scan, records are exactly where consecutive entries strictly rise -- pairwise makes those visible.

\noindent\texttt{new\_records(temps: list[int]) -> int}
\begin{lstlisting}[style=ex]
>>> new_records([3, 1, 4, 4, 5])
3
>>> new_records([2, 2, 2])
1
>>> new_records([1, 2, 3])
3
>>> new_records([])
0
\end{lstlisting}
\subsection*{6.7\quad Range-sum service\hfill\textnormal{\textit{apply}}}
A dashboard must answer many sum queries over the same fixed list: each query (i, j) asks for the sum of xs[i:j]. Precompute prefix sums ONCE with scanl, then answer every query in O(1) as pre[j] - pre[i].

\noindent\texttt{range\_sums(xs: list[int], queries: list[tuple[int, int]]) -> list[int]}
\begin{lstlisting}[style=ex]
>>> range_sums([3, 1, 4, 1], [(0, 2), (1, 3)])
[4, 5]
>>> range_sums([2, 2, 2], [(0, 3)])
[6]
>>> range_sums([1, 2], [])
[]
>>> range_sums([5], [(0, 0)])
[0]
\end{lstlisting}
\subsection*{6.8\quad Fuel gauge report\hfill\textnormal{\textit{apply}}}
A rover starts with 'start' units of fuel and applies signed fuel deltas (burns and top-ups) in order. Report (lowest\_level\_ever, mission\_ok) where mission\_ok means the level -- including the starting level -- never went negative. Scan the history once, then read both answers off it.

\noindent\texttt{fuel\_report(start: int, deltas: list[int]) -> tuple[int, bool]}
\begin{lstlisting}[style=ex]
>>> fuel_report(10, [-4, -5, 3, -6])
(-2, False)
>>> fuel_report(5, [-1])
(4, True)
>>> fuel_report(0, [])
(0, True)
\end{lstlisting}
\subsection*{6.9\quad Maximum drawdown\hfill\textnormal{\textit{apply}}}
Risk teams measure a stock's worst fall from a prior peak: the maximum of peak\_so\_far - price over all days. Compute it by zipping the running maximum against the prices themselves; a never-falling series has drawdown 0.

\noindent\texttt{drawdown(prices: list[int]) -> int}
\begin{lstlisting}[style=ex]
>>> drawdown([5, 3, 8, 4, 7])
4
>>> drawdown([1, 2, 3])
0
>>> drawdown([10])
0
>>> drawdown([])
0
\end{lstlisting}
\subsection*{6.10\quad Peak concurrency\hfill\textnormal{\textit{apply}}}
Sessions occupy overlapping intervals of time; the log records them as +1 for each opening and -1 for each closing, in time order. Report the highest number of simultaneously open sessions; with no events the answer is 0. Scan the deltas from 0 and take the highest point.

\noindent\texttt{peak\_concurrency(events: list[int]) -> int}
\begin{lstlisting}[style=ex]
>>> peak_concurrency([1, 1, -1, 1, -1])
2
>>> peak_concurrency([1, -1, 1, -1])
1
>>> peak_concurrency([])
0
\end{lstlisting}
\subsection*{6.11\quad Trapped rainwater\hfill\textnormal{\textit{challenge}}}
After a storm, water sits on a skyline of bar heights: above each bar, water fills up to the LOWER of the tallest bar to its left and the tallest to its right (itself included), minus the bar's own height. Total the trapped water. Two scans -- a running max from the left and one from the right -- plus one zipWith3 finish it in O(n). The classic interview problem, solved the scan way.

\noindent\texttt{trapped(heights: list[int]) -> int}
\begin{lstlisting}[style=ex]
>>> trapped([0, 1, 0, 2, 1, 0, 1, 3, 2, 1, 2, 1])
6
>>> trapped([2, 0, 2])
2
>>> trapped([1, 2, 3])
0
>>> trapped([])
0
\end{lstlisting}
\subsection*{6.12\quad Best sightseeing pair\hfill\textnormal{\textit{challenge}}}
A road trip visits spots with scores. Choosing spots i < j earns score[i] + score[j] + i - j (distance decays the fun). Maximise over all pairs, in O(n): note the formula splits into (score[i] + i) + (score[j] - j), so a running maximum of the first part, one step behind, meets each j. At least two spots are guaranteed.

\noindent\texttt{best\_pair(scores: list[int]) -> int}
\begin{lstlisting}[style=ex]
>>> best_pair([8, 1, 5, 2, 6])
11
>>> best_pair([1, 2])
2
>>> best_pair([1, 1, 1])
1
\end{lstlisting}
\clearpage
\chapter{Streams and Unfolds}

Production instead of consumption. Generators let a sequence be infinite --
\texttt{count}, \texttt{cycle}, \texttt{iterate} -- because nothing is
computed until a consumer asks; \texttt{unfoldr} is the finite twin, growing
a list from a seed until the step function says stop. Anything that PEELS
(digits, chunks, trajectories, parent chains) is an unfold, and recognising
that retires a whole family of while-loops with threaded state.
\section{Concepts}
\subsection{\texttt{count}}\label{def:count}
the infinite arithmetic sequence start, start+step, ... Never consume it whole: pair it with take, takewhile or find, which stop.
\begin{lstlisting}[style=ex]
>>> take(3, count(10))
[10, 11, 12]
>>> take(3, count(0, 5))
[0, 5, 10]
\end{lstlisting}

Roll your own --- the definition:

\begin{lstlisting}[style=ex]
def count(start=0, step=1):  # [start, start+step ..]
    n = start
    while True:
        yield n
        n += step
\end{lstlisting}
\noindent{\small\itshape Used before it is rolled: \texttt{take} in \S\ref{def:take} (p.~\pageref{def:take}).}

\subsection{\texttt{repeat}}\label{def:repeat}
one value forever, or exactly n times with the second argument.
\begin{lstlisting}[style=ex]
>>> take(2, repeat('x'))
['x', 'x']
>>> list(repeat('x', 3))
['x', 'x', 'x']
>>> list(repeat('x', 0))
[]
\end{lstlisting}

Roll your own --- the definition:

\begin{lstlisting}[style=ex]
def repeat(x, n=None):
    if n is None:
        while True:
            yield x
    else:
        for _ in range(n):
            yield x
\end{lstlisting}
\noindent{\small\itshape Used before it is rolled: \texttt{take} in \S\ref{def:take} (p.~\pageref{def:take}).}

\subsection{\texttt{cycle}}\label{def:cycle}
a finite list looped forever; an EMPTY list raises immediately (Haskell errors too -- the alternative is a silent infinite hang).
\begin{lstlisting}[style=ex]
>>> take(5, cycle([1, 2]))
[1, 2, 1, 2, 1]
\end{lstlisting}

Roll your own --- the definition:

\begin{lstlisting}[style=ex]
def cycle(xs):
    xs = list(xs)
    if not xs:
        raise ValueError("cycle: empty list")  # Haskell errors too; the alternative is a silent hang
    while True:
        yield from xs
\end{lstlisting}
\noindent{\small\itshape Used before it is rolled: \texttt{take} in \S\ref{def:take} (p.~\pageref{def:take}).}

\subsection{\texttt{iterate}}\label{def:iterate}
x, f(x), f(f(x)), ... forever: repeated application as a stream. Any deterministic state machine becomes iterate on its step function -- Fibonacci is iterate on a pair.
\begin{lstlisting}[style=ex]
>>> take(4, iterate(lambda x: 2 * x, 1))
[1, 2, 4, 8]
>>> step = lambda p: (snd(p), fst(p) + snd(p))
>>> map_(fst, take(6, iterate(step, (0, 1))))
[0, 1, 1, 2, 3, 5]
\end{lstlisting}

Roll your own --- the definition:

\begin{lstlisting}[style=ex]
def iterate(f, x):  # iterate f x = [x, f x, f (f x), ..]
    while True:
        yield x
        x = f(x)
\end{lstlisting}
\noindent{\small\itshape Used before it is rolled: \texttt{take} in \S\ref{def:take} (p.~\pageref{def:take}).}

\subsection{\texttt{unfoldr}}\label{def:unfoldr}
iterate's FINITE twin, and the cure for threading (state, acc) tuples through loops. You supply f(seed) returning either (value, seed') to emit a value and continue, or None to stop; unfoldr collects the values. Anything that "peels" a structure -- digits, chunks, parent chains -- is an unfoldr.
\begin{lstlisting}[style=ex]
>>> unfoldr(lambda m: None if m == 0 else (m % 10, m // 10), 407)
[7, 0, 4]
>>> unfoldr(lambda n: None if n == 0 else (n, n - 1), 5)
[5, 4, 3, 2, 1]
>>> unfoldr(const(None), "seed")        # stop immediately: empty
[]
\end{lstlisting}

How it works: the step function carries a two-outcome contract -- f(seed) returns (value, next\_seed) to emit and continue, or None to stop -- and unfoldr collects what was emitted. It is the fold REVERSED: a fold consumes a list into a summary; an unfold grows a list out of a seed. (Category theory says anamorphism; interviewers say "generate the sequence".)


Why it earns its place: it retires the (state, accumulator) tuple you would otherwise thread through a while loop. Anything that PEELS -- digits off a number, chunks off a list, parents up a tree, a Collatz trajectory -- states its entire logic in the step function.

\begin{lstlisting}[style=ex]
>>> unfoldr(lambda n: None if n == 1 else (n, n // 2 if even(n) else 3 * n + 1), 6)
[6, 3, 10, 5, 16, 8, 4, 2]
\end{lstlisting}

Versus iterate: iterate is the infinite cousin -- an unfold whose step never says None, consumed with take. Reach for iterate when the stream is conceptually endless, unfoldr when the seed itself knows where the end is.


Remember it as: fold tears down; unfold builds up; the seed carries the future.


More unfolds, same two-outcome contract -- peeling chunks off a list, bits off a number, walking a parent chain to its root:

\begin{lstlisting}[style=ex]
>>> unfoldr(lambda xs: None if xs == [] else (xs[:2], xs[2:]), [1, 2, 3, 4, 5])
[[1, 2], [3, 4], [5]]
>>> unfoldr(lambda m: None if m == 0 else (m % 2, m // 2), 13)       # bits, LSB first
[1, 0, 1, 1]
>>> parents = {'c': 'b', 'b': 'a', 'a': None}
>>> unfoldr(lambda n: None if n is None else (n, parents[n]), 'c')   # ancestry
['c', 'b', 'a']
\end{lstlisting}

And the seed can carry several facts at once -- here (current, next, remaining) generates Fibonacci with its own countdown built in:

\begin{lstlisting}[style=ex]
>>> unfoldr(lambda s: None if s[2] == 0 else (s[0], (s[1], s[0] + s[1], s[2] - 1)), (0, 1, 8))
[0, 1, 1, 2, 3, 5, 8, 13]
\end{lstlisting}

Roll your own: prime the pump with step = f(seed), then loop while step is not None -- append the value, feed the new seed back to f. The two-outcome contract lives entirely in f; unfoldr is just the loop that honours it, plus the list that collects what was emitted.

\begin{lstlisting}[style=ex]
def unfoldr(f, seed):  # Data.List unfoldr: f(seed) -> None | (value, seed')
    out = []  # iterate's finite twin: grows a list until f says stop --
    step = f(seed)  # no more threading (state, acc) tuples through until
    while step is not None:
        val, seed = step
        out.append(val)
        step = f(seed)
    return out
\end{lstlisting}
\noindent{\small\itshape Used before it is rolled: \texttt{until} in \S\ref{def:until} (p.~\pageref{def:until}), \texttt{const} in \S\ref{def:const} (p.~\pageref{def:const}).}

\subsection{\texttt{chain}}\label{def:chain}
concatenate iterables lazily, one after another: a generator expression, so the tail is never touched until asked for (chain over count() is fine). concat is the strict, list-building cousin.
\begin{lstlisting}[style=ex]
>>> list(chain([1], [2, 3]))
[1, 2, 3]
>>> take(4, chain([1, 2], count(10)))
[1, 2, 10, 11]
\end{lstlisting}

Roll your own --- the definition:

\begin{lstlisting}[style=ex]
chain = lambda *iterables: (x for it in iterables for x in it)  # concat, lazily
\end{lstlisting}
\noindent{\small\itshape Used before it is rolled: \texttt{take} in \S\ref{def:take} (p.~\pageref{def:take}).}

\section{Homework}
Work in order: drills cement each tool, applies combine them,
challenges are interview-grade. Write each solution in a scratch file with
the given doctests pasted in, and let \texttt{python3 -m doctest} referee.
\subsection*{7.1\quad First multiples\hfill\textnormal{\textit{drill}}}
Return the first n positive multiples of k, in order. Build the infinite stream of multiples with count and consume it with take -- no arithmetic on indices.

\noindent\texttt{multiples(k: int, n: int) -> list[int]}
\begin{lstlisting}[style=ex]
>>> multiples(3, 4)
[3, 6, 9, 12]
>>> multiples(5, 1)
[5]
>>> multiples(2, 0)
[]
\end{lstlisting}
\noindent{\small\itshape Rolled later: \texttt{take} in \S\ref{def:take} (p.~\pageref{def:take}).}

\subsection*{7.2\quad Banner line\hfill\textnormal{\textit{drill}}}
Build a banner string of n copies of a single character using repeat (the bounded, second-argument form) and join.

\noindent\texttt{banner(ch: str, n: int) -> str}
\begin{lstlisting}[style=ex]
>>> banner('*', 3)
'***'
>>> banner('-', 0)
''
>>> banner('=', 5)
'====='
\end{lstlisting}
\subsection*{7.3\quad Powers table\hfill\textnormal{\textit{drill}}}
Return the first n powers of b starting at b**0, using iterate -- repeated multiplication as a stream, no ** operator.

\noindent\texttt{powers(b: int, n: int) -> list[int]}
\begin{lstlisting}[style=ex]
>>> powers(2, 5)
[1, 2, 4, 8, 16]
>>> powers(10, 3)
[1, 10, 100]
>>> powers(3, 0)
[]
\end{lstlisting}
\noindent{\small\itshape Rolled later: \texttt{take} in \S\ref{def:take} (p.~\pageref{def:take}).}

\subsection*{7.4\quad Countdown\hfill\textnormal{\textit{drill}}}
Return [n, n-1, ..., 1] with unfoldr: the seed is the current number, the unfold stops at zero. countdown(0) is empty.

\noindent\texttt{countdown(n: int) -> list[int]}
\begin{lstlisting}[style=ex]
>>> countdown(4)
[4, 3, 2, 1]
>>> countdown(1)
[1]
>>> countdown(0)
[]
\end{lstlisting}
\subsection*{7.5\quad Splice three feeds\hfill\textnormal{\textit{drill}}}
Concatenate three iterables into one list using chain -- lists, strings, anything iterable, in the order given.

\noindent\texttt{splice(a, b, c) -> list}
\begin{lstlisting}[style=ex]
>>> splice([1], [2, 3], [])
[1, 2, 3]
>>> splice('ab', 'cd', 'e')
['a', 'b', 'c', 'd', 'e']
>>> splice([], [], [])
[]
\end{lstlisting}
\subsection*{7.6\quad Round-robin card deal\hfill\textnormal{\textit{apply}}}
A dealer hands cards to players in strict rotation: first card to the first player, second to the second, wrapping around until the cards run out. Return the (player, card) pairs in deal order. At least one player is guaranteed.

\noindent\texttt{deal(players: list[str], cards: list) -> list[tuple]}
\begin{lstlisting}[style=ex]
>>> deal(['ann', 'bo'], [1, 2, 3, 4, 5])
[('ann', 1), ('bo', 2), ('ann', 3), ('bo', 4), ('ann', 5)]
>>> deal(['solo'], ['x', 'y'])
[('solo', 'x'), ('solo', 'y')]
>>> deal(['a', 'b'], [])
[]
\end{lstlisting}
\subsection*{7.7\quad Colony doubling\hfill\textnormal{\textit{apply}}}
A bacteria colony starts at 'start' cells and doubles every day. Return how many days pass before the colony first reaches at least 'goal' cells (0 if it already has). Model the growth as an iterate stream and count the days spent below goal.

\noindent\texttt{days\_to\_reach(start: int, goal: int) -> int}
\begin{lstlisting}[style=ex]
>>> days_to_reach(1, 8)
3
>>> days_to_reach(5, 5)
0
>>> days_to_reach(3, 100)
6
\end{lstlisting}
\noindent{\small\itshape Rolled later: \texttt{takeWhile} in \S\ref{def:takeWhile} (p.~\pageref{def:takeWhile}).}

\subsection*{7.8\quad LED bit pattern\hfill\textnormal{\textit{apply}}}
An LED row shows a number in binary, most significant bit first. Return the bit list for n >= 1: peel bits least-significant-first with unfoldr, then reverse.

\noindent\texttt{bits(n: int) -> list[int]}
\begin{lstlisting}[style=ex]
>>> bits(13)
[1, 1, 0, 1]
>>> bits(1)
[1]
>>> bits(8)
[1, 0, 0, 0]
\end{lstlisting}
\subsection*{7.9\quad Chain of command\hfill\textnormal{\textit{apply}}}
An org chart maps each employee to their manager, with the top boss mapped to None. Return the chain of command from a given employee up to and including the top: an unfold whose seed is the current person and whose next seed is their manager.

\noindent\texttt{chain\_of\_command(boss: dict, name: str) -> list[str]}
\begin{lstlisting}[style=ex]
>>> boss = {'amy': 'raj', 'raj': 'kim', 'kim': None}
>>> chain_of_command(boss, 'amy')
['amy', 'raj', 'kim']
>>> chain_of_command(boss, 'kim')
['kim']
>>> chain_of_command({'x': None}, 'x')
['x']
\end{lstlisting}
\subsection*{7.10\quad When the buses meet\hfill\textnormal{\textit{apply}}}
Bus A stops at minute a, a+da, a+2*da, ...; bus B at b, b+db, .... Return the first minute at or after both start times when the two schedules coincide. A shared minute is guaranteed to exist. Search A's infinite schedule with find.

\noindent\texttt{sync(a: int, da: int, b: int, db: int) -> int}
\begin{lstlisting}[style=ex]
>>> sync(3, 4, 1, 6)
7
>>> sync(10, 1, 4, 7)
11
>>> sync(0, 5, 0, 3)
0
\end{lstlisting}
\subsection*{7.11\quad Look-and-say\hfill\textnormal{\textit{challenge}}}
The look-and-say sequence starts '1'; each term reads the previous aloud: '1' is one 1 -> '11'; '11' is two 1s -> '21'; '21' -> '1211'; and so on. Return the nth term (1-indexed). Peel each term into runs of equal digits with unfoldr, encode each run as count then digit, and drive the whole sequence with iterate.

\noindent\texttt{look\_and\_say(n: int) -> str}
\begin{lstlisting}[style=ex]
>>> look_and_say(1)
'1'
>>> look_and_say(2)
'11'
>>> look_and_say(3)
'21'
>>> look_and_say(5)
'111221'
\end{lstlisting}
\noindent{\small\itshape Rolled later: \texttt{take} in \S\ref{def:take} (p.~\pageref{def:take}).}

\subsection*{7.12\quad Hailstone champion\hfill\textnormal{\textit{challenge}}}
The hailstone (Collatz) trajectory of n repeatedly halves even numbers and maps odd m to 3*m + 1, ending at 1. Among all starts 1..limit, return the start with the LONGEST trajectory (counting every term including the final 1); ties go to the smaller start. Express one trajectory as an unfoldr and measure it.

\noindent\texttt{hail\_champion(limit: int) -> int}
\begin{lstlisting}[style=ex]
>>> hail_champion(1)
1
>>> hail_champion(10)
9
>>> hail_champion(30)
27
\end{lstlisting}
\clearpage
\chapter{Spans and Views}

The consumers and finite views that make streams usable: \texttt{take} and
friends stop lazily; \texttt{span}/\texttt{break\_} cut a sequence at a
condition -- a lexer in one call; \texttt{inits}/\texttt{tails} enumerate
every prefix and suffix, which is where substring problems begin; and
\texttt{chunksOf} pages anything. Small tools, constant use.
\section{Concepts}
\subsection{\texttt{islice}}\label{def:islice}
at most the first n items of ANY iterable, lazily; the safe window onto an infinite stream. Runs short without complaint.
\begin{lstlisting}[style=ex]
>>> list(islice(count(0), 3))
[0, 1, 2]
>>> list(islice([1, 2], 5))
[1, 2]
\end{lstlisting}

Roll your own --- the definition:

\begin{lstlisting}[style=ex]
def islice(iterable, stop):  # take
    it = iter(iterable)
    for _ in range(stop):
        try:
            yield next(it)
        except StopIteration:
            return
\end{lstlisting}
\noindent{\small\itshape Used before it is rolled: \texttt{take} in \S\ref{def:take} (p.~\pageref{def:take}).}

\subsection{\texttt{take}}\label{def:take}
the first n, as a real list. take/drop are the positional pair; takeWhile/dropWhile are the conditional pair.
\begin{lstlisting}[style=ex]
>>> take(3, count(1))
[1, 2, 3]
>>> take(2, [7, 8, 9])
[7, 8]
\end{lstlisting}

Roll your own --- the definition:

\begin{lstlisting}[style=ex]
take = lambda n, xs: list(islice(iter(xs), n))  # works on infinite streams
\end{lstlisting}
\noindent{\small\itshape Used before it is rolled: \texttt{on} in \S\ref{def:on} (p.~\pageref{def:on}).}

\subsection{\texttt{takewhile}}\label{def:takewhile}
yield elements WHILE the predicate holds, then stop at the first failure (a generator). Contrast filter\_: filter\_ scans everything; takewhile stops -- on [2, 4, 5, 6], filter\_ keeps the 6, takewhile does not.
\begin{lstlisting}[style=ex]
>>> list(takewhile(even, [2, 4, 5, 6]))
[2, 4]
>>> filter_(even, [2, 4, 5, 6])
[2, 4, 6]
\end{lstlisting}

Roll your own --- the definition:

\begin{lstlisting}[style=ex]
def takewhile(pred, xs):
    for x in xs:
        if not pred(x):
            return
        yield x
\end{lstlisting}
\subsection{\texttt{takeuntil}}\label{def:takeuntil}
like takewhile(not p), but INCLUDES the first element where p holds -- the do-while flavour: "read up to and including the terminator".
\begin{lstlisting}[style=ex]
>>> list(takeuntil(lambda x: x == 3, [1, 2, 3, 4]))
[1, 2, 3]
>>> list(takeuntil(lambda x: x > 0, [5, 6]))
[5]
\end{lstlisting}

Roll your own --- the definition:

\begin{lstlisting}[style=ex]
def takeuntil(pred, xs):  # like takewhile(not . pred), but INCLUDES the first hit
    for x in xs:
        yield x
        if pred(x):
            return
\end{lstlisting}
\subsection{\texttt{dropwhile}}\label{def:dropwhile}
discard while the predicate holds, then yield EVERYTHING after the first failure, no more testing. takewhile's complement.
\begin{lstlisting}[style=ex]
>>> list(dropwhile(lambda x: x < 3, [1, 2, 3, 1]))
[3, 1]
\end{lstlisting}

Roll your own --- the definition:

\begin{lstlisting}[style=ex]
def dropwhile(pred, xs):
    it = iter(xs)
    for x in it:
        if not pred(x):
            yield x
            break
    yield from it
\end{lstlisting}
\subsection{\texttt{takeWhile}}\label{def:takeWhile}
takewhile, delivered as a list (capital W = list version).
\begin{lstlisting}[style=ex]
>>> takeWhile(lambda x: x < 3, [1, 2, 3, 1])
[1, 2]
\end{lstlisting}

Roll your own --- the definition:

\begin{lstlisting}[style=ex]
takeWhile = lambda p, xs: list(takewhile(p, xs))
\end{lstlisting}
\subsection{\texttt{dropWhile}}\label{def:dropWhile}
dropwhile, delivered as a list.
\begin{lstlisting}[style=ex]
>>> dropWhile(lambda x: x < 3, [1, 2, 3, 1])
[3, 1]
\end{lstlisting}

Roll your own --- the definition:

\begin{lstlisting}[style=ex]
dropWhile = lambda p, xs: list(dropwhile(p, xs))
\end{lstlisting}
\subsection{\texttt{span}}\label{def:span}
split at the first failure: (longest satisfying prefix, the rest), computed in ONE pass so it is safe on one-shot iterators (the old two-scan version silently broke on generators -- this is the fix). A lexer in one call.
\begin{lstlisting}[style=ex]
>>> span(lambda x: x < 3, [1, 2, 3, 4, 1])
([1, 2], [3, 4, 1])
>>> span(str.isdigit, "123abc")
(['1', '2', '3'], ['a', 'b', 'c'])
\end{lstlisting}

Roll your own: materialise xs once (that is the one-pass guarantee), locate the first failing index with next() over enumerate -- defaulting to len(xs) when the predicate never fails -- and slice twice. The previous two-scan version ran takewhile and dropwhile separately and silently broke on one-shot generators; this is the repair.

\begin{lstlisting}[style=ex]
def span(p, xs):  # split at first failure, ONE pass; safe on one-shot iters
    xs = list(xs)
    i = next((i for i, x in enumerate(xs) if not p(x)), len(xs))
    return (xs[:i], xs[i:])
\end{lstlisting}
\noindent{\small\itshape Used before it is rolled: \texttt{on} in \S\ref{def:on} (p.~\pageref{def:on}).}

\subsection{\texttt{break\_}}\label{def:break-}
span with the predicate negated: split where the condition first HOLDS.
\begin{lstlisting}[style=ex]
>>> break_(lambda x: x > 2, [1, 2, 3, 1])
([1, 2], [3, 1])
\end{lstlisting}

Roll your own --- the definition:

\begin{lstlisting}[style=ex]
break_ = lambda p, xs: span(lambda x: not p(x), xs)
\end{lstlisting}
\subsection{\texttt{inits}}\label{def:inits}
every prefix, from empty to the whole thing; len(xs)+1 of them.
\begin{lstlisting}[style=ex]
>>> inits([1, 2])
[[], [1], [1, 2]]
>>> inits("ab")
['', 'a', 'ab']
\end{lstlisting}

Roll your own --- the definition:

\begin{lstlisting}[style=ex]
inits = lambda xs: [xs[:i] for i in range(len(xs) + 1)]  # Data.List: every prefix, [] first
\end{lstlisting}
\subsection{\texttt{tails}}\label{def:tails}
every suffix, whole thing first, down to empty. Substring problems start here: every substring is a prefix of some suffix.
\begin{lstlisting}[style=ex]
>>> tails("abc")
['abc', 'bc', 'c', '']
>>> tails([1, 2])
[[1, 2], [2], []]
\end{lstlisting}

Roll your own --- the definition:

\begin{lstlisting}[style=ex]
tails = lambda xs: [xs[i:] for i in range(len(xs) + 1)]  # every suffix; substrings start here
\end{lstlisting}
\subsection{\texttt{chunksOf}}\label{def:chunksOf}
consecutive pieces of length n; the last may be short. Works on lists and strings alike since it is pure slicing.
\begin{lstlisting}[style=ex]
>>> chunksOf(2, [1, 2, 3, 4, 5])
[[1, 2], [3, 4], [5]]
>>> chunksOf(3, "abcdefg")
['abc', 'def', 'g']
\end{lstlisting}

Roll your own --- the definition:

\begin{lstlisting}[style=ex]
chunksOf = lambda n, xs: [xs[i:i + n] for i in range(0, len(xs), n)]  # Data.List.Split; short last ok
\end{lstlisting}
\section{Homework}
Work in order: drills cement each tool, applies combine them,
challenges are interview-grade. Write each solution in a scratch file with
the given doctests pasted in, and let \texttt{python3 -m doctest} referee.
\subsection*{8.1\quad Leading run\hfill\textnormal{\textit{drill}}}
Return the longest prefix of xs whose elements are all strictly positive, using span. The rest of the list is discarded.

\noindent\texttt{leading\_positive(xs: list) -> list}
\begin{lstlisting}[style=ex]
>>> leading_positive([2, 5, 1, -1, 9])
[2, 5, 1]
>>> leading_positive([-1, 2])
[]
>>> leading_positive([4, 4])
[4, 4]
\end{lstlisting}
\subsection*{8.2\quad Peek first k\hfill\textnormal{\textit{drill}}}
Return the first k elements of a possibly-infinite stream with take. If the stream is shorter than k, return all of it -- take never complains.

\noindent\texttt{peek(k: int, xs) -> list}
\begin{lstlisting}[style=ex]
>>> peek(3, count(0))
[0, 1, 2]
>>> peek(5, [1, 2])
[1, 2]
>>> peek(0, count(9))
[]
\end{lstlisting}
\subsection*{8.3\quad Skip the header\hfill\textnormal{\textit{drill}}}
A list begins with comment lines starting '\#'. Drop that leading block and return the remaining lines, using dropWhile. Once a non-comment line appears, keep everything.

\noindent\texttt{skip\_header(lines: list) -> list}
\begin{lstlisting}[style=ex]
>>> skip_header(['# a', '# b', 'x', '# c'])
['x', '# c']
>>> skip_header(['data'])
['data']
>>> skip_header(['# only'])
[]
\end{lstlisting}
\noindent{\small\itshape Rolled later: \texttt{lines} in \S\ref{def:lines} (p.~\pageref{def:lines}).}

\subsection*{8.4\quad Every prefix\hfill\textnormal{\textit{drill}}}
Return every prefix of xs from empty to the whole list, using inits.

\noindent\texttt{prefixes(xs: list) -> list[list]}
\begin{lstlisting}[style=ex]
>>> prefixes([1, 2, 3])
[[], [1], [1, 2], [1, 2, 3]]
>>> prefixes([])
[[]]
>>> prefixes(['a'])
[[], ['a']]
\end{lstlisting}
\subsection*{8.5\quad Deal into pages\hfill\textnormal{\textit{drill}}}
Split a list into consecutive pages of size k; the final page may be short. Use chunksOf.

\noindent\texttt{paginate(k: int, xs: list) -> list[list]}
\begin{lstlisting}[style=ex]
>>> paginate(2, [1, 2, 3, 4, 5])
[[1, 2], [3, 4], [5]]
>>> paginate(3, [1, 2, 3])
[[1, 2, 3]]
>>> paginate(2, [])
[]
\end{lstlisting}
\subsection*{8.6\quad Read the leading integer\hfill\textnormal{\textit{apply}}}
A token string begins with an optional run of digits, then arbitrary characters. Return the leading integer, or 0 if no digit starts it. Use span to split the digit run from the tail, then interpret it.

\noindent\texttt{read\_int(s: str) -> int}
\begin{lstlisting}[style=ex]
>>> read_int('123abc')
123
>>> read_int('7')
7
>>> read_int('abc')
0
>>> read_int('')
0
\end{lstlisting}
\subsection*{8.7\quad Split on the first blank\hfill\textnormal{\textit{apply}}}
An email is a header block followed by a blank line ('') and then the body. Return (headers, body) where headers excludes the blank separator and body is the lines after it. Use break\_ to cut at the first blank line. If there is no blank line, the body is empty.

\noindent\texttt{split\_email(lines: list) -> tuple}
\begin{lstlisting}[style=ex]
>>> split_email(['To: a', 'Re: b', '', 'hello', 'bye'])
(['To: a', 'Re: b'], ['hello', 'bye'])
>>> split_email(['To: a', 'Re: b'])
(['To: a', 'Re: b'], [])
>>> split_email(['', 'body'])
([], ['body'])
\end{lstlisting}
\noindent{\small\itshape Rolled later: \texttt{lines} in \S\ref{def:lines} (p.~\pageref{def:lines}).}

\subsection*{8.8\quad Read up to the sentinel\hfill\textnormal{\textit{apply}}}
A sensor emits readings terminated by a sentinel value -1. Return the readings up to AND including the first -1 (the do-while flavour: the terminator is part of the record). Use takeuntil. If no -1 appears, return the whole list.

\noindent\texttt{read\_record(xs: list) -> list}
\begin{lstlisting}[style=ex]
>>> read_record([4, 7, -1, 9])
[4, 7, -1]
>>> read_record([1, 2, 3])
[1, 2, 3]
>>> read_record([-1, 5])
[-1]
\end{lstlisting}
\subsection*{8.9\quad Longest good prefix length\hfill\textnormal{\textit{apply}}}
Given a list of daily temperatures, a warm streak from day 0 runs while each day is at least 'floor'. Return how many opening days form that streak, using span to isolate the prefix and measuring it. A first day below floor gives 0.

\noindent\texttt{warm\_streak(temps: list, floor: int) -> int}
\begin{lstlisting}[style=ex]
>>> warm_streak([70, 72, 68, 90], 70)
2
>>> warm_streak([60, 80], 70)
0
>>> warm_streak([75, 75, 75], 70)
3
\end{lstlisting}
\subsection*{8.10\quad All substrings\hfill\textnormal{\textit{apply}}}
Return every non-empty contiguous substring of s, in the order: by start position, then by length. The tool is tails -- every substring is a non-empty prefix of some suffix. So take each suffix's inits, drop the empty one, and concatenate.

\noindent\texttt{substrings(s: str) -> list[str]}
\begin{lstlisting}[style=ex]
>>> substrings('ab')
['a', 'ab', 'b']
>>> substrings('a')
['a']
>>> substrings('')
[]
>>> substrings('xyz')
['x', 'xy', 'xyz', 'y', 'yz', 'z']
\end{lstlisting}
\subsection*{8.11\quad Run-length split\hfill\textnormal{\textit{challenge}}}
Split a list into maximal runs of equal consecutive elements, WITHOUT using group or groupBy -- do it directly with span. Peel the leading run (all elements equal to the head) off the front repeatedly until the list is empty.

\noindent\texttt{runs(xs: list) -> list[list]}
\begin{lstlisting}[style=ex]
>>> runs([1, 1, 2, 3, 3, 3])
[[1, 1], [2], [3, 3, 3]]
>>> runs([])
[]
>>> runs(['a'])
[['a']]
>>> runs([5, 5, 5])
[[5, 5, 5]]
\end{lstlisting}
\subsection*{8.12\quad Justify to width\hfill\textnormal{\textit{challenge}}}
Greedily wrap words into lines no longer than 'width' characters, words joined by single spaces; each line packs as many words as fit. Every word is guaranteed to fit alone. Peel one line at a time: take the longest prefix of the remaining words whose joined length stays within width, emit it, continue with the rest.

\noindent\texttt{justify(words: list, width: int) -> list[str]}
\begin{lstlisting}[style=ex]
>>> justify(['a', 'bb', 'ccc'], 6)
['a bb', 'ccc']
>>> justify(['one', 'two', 'three'], 7)
['one two', 'three']
>>> justify([], 5)
[]
>>> justify(['hello'], 5)
['hello']
\end{lstlisting}
\noindent{\small\itshape Rolled later: \texttt{words} in \S\ref{def:words} (p.~\pageref{def:words}).}

\clearpage
\chapter{Order and Text}

``Does order unlock it?'' is the recognition ladder's first question because
sorting linearises problems: tuple keys buy multi-level orderings in one
call, \texttt{minOn}/\texttt{maxOn} answer best-by-key in $O(n)$, and the
bisect pair reads a sorted list's True/False boundary three different ways
-- membership, insertion point, counting. The string quartet turns text
pipelines into list pipelines and back.
\section{Concepts}
\subsection{\texttt{sortOn}}\label{def:sortOn}
sort by a key function, evaluated once per element. TUPLE KEYS are the power move: (primary, secondary) sorts two levels at once, and negating a numeric component flips just that level -- (-count, word) is "count descending, ties alphabetical", which reverse=True cannot express.
\begin{lstlisting}[style=ex]
>>> sortOn(len, ["bbb", "a", "cc"])
['a', 'cc', 'bbb']
>>> sortOn(lambda kv: (-kv[1], kv[0]), [('b', 2), ('a', 2), ('c', 3)])
[('c', 3), ('a', 2), ('b', 2)]
\end{lstlisting}

Roll your own --- the definition:

\begin{lstlisting}[style=ex]
sortOn = lambda f, xs: sorted(xs, key=f)  # sortOn (schwartzian, f called once per element)
\end{lstlisting}
\subsection{\texttt{minOn}}\label{def:minOn}
the element with the smallest key, in O(n). sortOn + head gives the same answer for O(n log n) -- pay the sort only when you need ALL of them ordered.
\begin{lstlisting}[style=ex]
>>> minOn(abs, [-3, 2, 5])
2
>>> minOn(snd, [('a', 9), ('b', 1)])
('b', 1)
\end{lstlisting}

Roll your own --- the definition:

\begin{lstlisting}[style=ex]
minOn = lambda f, xs: min(xs, key=f)  # minimumBy (comparing f): sortOn + head without the sort
\end{lstlisting}
\subsection{\texttt{maxOn}}\label{def:maxOn}
the element with the largest key, O(n); the FIRST winner on ties, so pre-sorting the input controls the tie-break.
\begin{lstlisting}[style=ex]
>>> maxOn(len, ["aa", "b", "cc"])
'aa'
\end{lstlisting}

Roll your own --- the definition:

\begin{lstlisting}[style=ex]
maxOn = lambda f, xs: max(xs, key=f)  # maximumBy (comparing f) in O(n); first wins ties
\end{lstlisting}
\subsection{\texttt{cmp\_to\_key}}\label{def:cmp-to-key}
adapts an old-style comparator (returning negative, zero or positive) into the key= protocol. Reach for it only when a rule genuinely compares two elements against each other and no per- element key exists.
\begin{lstlisting}[style=ex]
>>> sorted(["bb", "a"], key=cmp_to_key(lambda a, b: len(a) - len(b)))
['a', 'bb']
\end{lstlisting}

Roll your own: a tiny wrapper class implementing ONE dunder -- \_\_lt\_\_ defers to cmp(a, b) < 0. Python's sort only ever asks "is a less than b?", so that single method is enough to smuggle any comparator through the key= interface.

\begin{lstlisting}[style=ex]
def cmp_to_key(cmp):  # resurrect Python 2's cmp for sortBy
    class K:
        def __init__(self, x): self.x = x
        def __lt__(self, other): return cmp(self.x, other.x) < 0
    return K
\end{lstlisting}
\noindent{\small\itshape Used before it is rolled: \texttt{sortBy} in \S\ref{def:sortBy} (p.~\pageref{def:sortBy}).}

\subsection{\texttt{sortBy}}\label{def:sortBy}
comparator sort, Python 2 style, built on cmp\_to\_key.
\begin{lstlisting}[style=ex]
>>> sortBy(lambda a, b: b - a, [1, 3, 2])
[3, 2, 1]
\end{lstlisting}

Roll your own --- the definition:

\begin{lstlisting}[style=ex]
sortBy = lambda cmp, xs: sorted(xs, key=cmp_to_key(cmp))  # comparator sort, Python 2 style
\end{lstlisting}
\subsection{\texttt{merge}}\label{def:merge}
join two ALREADY-SORTED lists into one sorted list, stably, in O(n+m): the two-pointer loop at the heart of merge sort.
\begin{lstlisting}[style=ex]
>>> merge([1, 4, 6], [2, 3, 7])
[1, 2, 3, 4, 6, 7]
>>> merge([], [1, 2])
[1, 2]
\end{lstlisting}

Roll your own: two cursors, always copying the smaller head forward; the <= (rather than <) is what makes the merge STABLE, preferring the left list on ties. When either side runs dry, glue both tails on -- one of them is empty, so the + is harmless.

\begin{lstlisting}[style=ex]
def merge(a, b):  # merge two sorted lists, stable; the heart of merge sort
    out, i, j = [], 0, 0  # (section 13's merge_lists is this same loop on ListNodes)
    while i < len(a) and j < len(b):
        if a[i] <= b[j]:
            out.append(a[i]); i += 1
        else:
            out.append(b[j]); j += 1
    return out + a[i:] + b[j:]
\end{lstlisting}
\noindent{\small\itshape Used before it is rolled: \texttt{merge\_lists} in \S\ref{def:merge-lists} (p.~\pageref{def:merge-lists}), \texttt{on} in \S\ref{def:on} (p.~\pageref{def:on}).}

\subsection{\texttt{bisect\_left}}\label{def:bisect-left}
binary search over a SORTED list: the first index whose element is >= x; equivalently, where x would insert to the left of its equals. If a[i] != x there, x is absent -- that check is binary search.
\begin{lstlisting}[style=ex]
>>> bisect_left([10, 20, 20, 30], 20)
1
>>> bisect_left([10, 20, 30], 25)       # insertion point for a missing value
2
\end{lstlisting}

How it works: on a sorted list the predicate "element < x" reads True for a prefix and False for the rest -- TTTTFFFF. Binary search does not hunt for x; it finds that BOUNDARY, halving the interval per probe. bisect\_left returns the first index >= x, bisect\_right the first > x, and from those two facts three idioms fall out.

\begin{lstlisting}[style=ex]
>>> xs = [10, 20, 20, 30]
>>> bisect_left(xs, 20), bisect_right(xs, 20)     # the run of 20s occupies [1, 3)
(1, 3)
\end{lstlisting}

Idiom one, membership: x is present iff the slot bisect\_left names actually holds it. Idiom two, insertion point: that same index is where x belongs to keep the list sorted. Idiom three, counting: right minus left counts occurrences -- and with two different probes, the values inside any range.

\begin{lstlisting}[style=ex]
>>> i = bisect_left(xs, 25)
>>> i, (i < len(xs) and xs[i] == 25)              # 25 would insert at 3; absent
(3, False)
\end{lstlisting}

Why it earns its place: every threshold ladder is secretly a sorted list. Grading bands, tax brackets, rate tiers -- one bisect replaces the if/elif staircase:

\begin{lstlisting}[style=ex]
>>> grade = lambda score: "FDCBA"[bisect_right([60, 70, 80, 90], score)]
>>> grade(59), grade(60), grade(95)
('F', 'D', 'A')
\end{lstlisting}

Remember it as: bisect finds the boundary in a sea of Trues and Falses; everything else is reading that boundary three ways.


Roll your own: the lo/hi halving loop over an invariant -- everything left of lo is < x, everything from hi rightwards is >= x. One comparison, a[mid] < x, decides which half survives; when lo meets hi, that meeting point IS the boundary. bisect\_right is the identical loop with <= in place of < -- one character moves the boundary to the far side of the equals.

\begin{lstlisting}[style=ex]
def bisect_left(a, x):  # first index where a[i] >= x
    lo, hi = 0, len(a)
    while lo < hi:
        mid = (lo + hi) // 2
        if a[mid] < x:
            lo = mid + 1
        else:
            hi = mid
    return lo
\end{lstlisting}
\noindent{\small\itshape Used before it is rolled: \texttt{bisect\_right} in \S\ref{def:bisect-right} (p.~\pageref{def:bisect-right}).}

\subsection{\texttt{bisect\_right}}\label{def:bisect-right}
the first index whose element is > x: where x would insert to the RIGHT of its equals. right minus left counts the occurrences.
\begin{lstlisting}[style=ex]
>>> bisect_right([10, 20, 20, 30], 20)
3
>>> xs = [10, 20, 20, 30]
>>> bisect_right(xs, 20) - bisect_left(xs, 20)
2
\end{lstlisting}

Roll your own --- the definition:

\begin{lstlisting}[style=ex]
def bisect_right(a, x):  # first index where a[i] > x
    lo, hi = 0, len(a)
    while lo < hi:
        mid = (lo + hi) // 2
        if a[mid] <= x:
            lo = mid + 1
        else:
            hi = mid
    return lo
\end{lstlisting}
\subsection{\texttt{words}}\label{def:words}
split on any whitespace; runs collapse, edges vanish. The inverse direction is unwords, and the round trip normalises spacing.
\begin{lstlisting}[style=ex]
>>> words("  such   spacing  ")
['such', 'spacing']
>>> unwords(words("  a   b "))
'a b'
\end{lstlisting}

Roll your own --- the definition:

\begin{lstlisting}[style=ex]
words = str.split
\end{lstlisting}
\noindent{\small\itshape Used before it is rolled: \texttt{unwords} in \S\ref{def:unwords} (p.~\pageref{def:unwords}).}

\subsection{\texttt{unwords}}\label{def:unwords}
join words with single spaces.
\begin{lstlisting}[style=ex]
>>> unwords(['a', 'b', 'c'])
'a b c'
\end{lstlisting}

Roll your own --- the definition:

\begin{lstlisting}[style=ex]
unwords = " ".join
\end{lstlisting}
\subsection{\texttt{lines}}\label{def:lines}
split a string into its lines (no trailing newlines kept).
\begin{lstlisting}[style=ex]
>>> lines("one\ntwo")
['one', 'two']
\end{lstlisting}

Roll your own --- the definition:

\begin{lstlisting}[style=ex]
lines = str.splitlines
\end{lstlisting}
\subsection{\texttt{unlines}}\label{def:unlines}
join lines, TERMINATING each with a newline (not separating -- note the trailing one).
\begin{lstlisting}[style=ex]
>>> unlines(['a', 'b'])
'a\nb\n'
\end{lstlisting}

Roll your own --- the definition:

\begin{lstlisting}[style=ex]
unlines = lambda ls: "".join(l + "\n" for l in ls)
\end{lstlisting}
\section{Homework}
Work in order: drills cement each tool, applies combine them,
challenges are interview-grade. Write each solution in a scratch file with
the given doctests pasted in, and let \texttt{python3 -m doctest} referee.
\subsection*{9.1\quad Sort by size, then spelling\hfill\textnormal{\textit{drill}}}
Sort a list of words shortest-first; words of equal length appear in alphabetical order. One sortOn with a tuple key does both levels at once.

\noindent\texttt{sort\_by\_size(ws: list[str]) -> list[str]}
\begin{lstlisting}[style=ex]
>>> sort_by_size(["bb", "a", "ccc", "aa"])
['a', 'aa', 'bb', 'ccc']
>>> sort_by_size(["same", "size", "abcd"])
['abcd', 'same', 'size']
>>> sort_by_size([])
[]
\end{lstlisting}
\subsection*{9.2\quad Closest to zero\hfill\textnormal{\textit{drill}}}
Return the element of a non-empty list closest to zero. If two elements tie in distance (like -2 and 2), the one appearing FIRST in the list wins.

\noindent\texttt{closest\_to\_zero(xs: list[int]) -> int}
\begin{lstlisting}[style=ex]
>>> closest_to_zero([3, -1, 4])
-1
>>> closest_to_zero([-2, 2])
-2
>>> closest_to_zero([7])
7
\end{lstlisting}
\subsection*{9.3\quad Count occurrences in sorted data\hfill\textnormal{\textit{drill}}}
Given an ALREADY-SORTED list and a value, count how many times the value occurs -- in O(log n), using the two bisects. No scanning allowed.

\noindent\texttt{count\_val(sorted\_xs: list[int], x: int) -> int}
\begin{lstlisting}[style=ex]
>>> count_val([1, 2, 2, 2, 3], 2)
3
>>> count_val([1, 3], 2)
0
>>> count_val([], 5)
0
\end{lstlisting}
\subsection*{9.4\quad Normalize spacing\hfill\textnormal{\textit{drill}}}
Collapse any runs of whitespace in a sentence to single spaces and strip the ends, using the words/unwords round trip.

\noindent\texttt{normalize\_spaces(s: str) -> str}
\begin{lstlisting}[style=ex]
>>> normalize_spaces("  such   spacing  ")
'such spacing'
>>> normalize_spaces("one")
'one'
>>> normalize_spaces("")
''
\end{lstlisting}
\subsection*{9.5\quad Three-way merge\hfill\textnormal{\textit{drill}}}
Merge three already-sorted lists into one sorted list by chaining the two-list merge. Stability must hold: on equal values, earlier-argument elements come first.

\noindent\texttt{merge3(a: list, b: list, c: list) -> list}
\begin{lstlisting}[style=ex]
>>> merge3([1, 4], [2, 5], [3, 6])
[1, 2, 3, 4, 5, 6]
>>> merge3([], [7], [])
[7]
>>> merge3([1, 1], [1], [0])
[0, 1, 1, 1]
\end{lstlisting}
\subsection*{9.6\quad Tournament leaderboard\hfill\textnormal{\textit{apply}}}
A tournament reports (name, score) pairs. Produce the leaderboard: names only, highest score first, and players with equal scores listed alphabetically. One key, two levels.

\noindent\texttt{leaderboard(entries: list[tuple[str, int]]) -> list[str]}
\begin{lstlisting}[style=ex]
>>> leaderboard([("ann", 50), ("bob", 70), ("cy", 50)])
['bob', 'ann', 'cy']
>>> leaderboard([("zed", 10)])
['zed']
>>> leaderboard([])
[]
\end{lstlisting}
\subsection*{9.7\quad Shipping classes\hfill\textnormal{\textit{apply}}}
Parcels are classed by weight with cutoffs at 1, 5 and 20 kilograms: under 1 is an 'envelope', then 'small', then 'medium', and 20 or over is 'freight' -- a parcel exactly at a cutoff belongs to the higher class. Replace the if/elif ladder with one bisect over the cutoffs.

\noindent\texttt{shipping\_class(w: float) -> str}
\begin{lstlisting}[style=ex]
>>> shipping_class(0.5)
'envelope'
>>> shipping_class(1)
'small'
>>> shipping_class(4.9)
'small'
>>> shipping_class(20)
'freight'
\end{lstlisting}
\subsection*{9.8\quad Version sort\hfill\textnormal{\textit{apply}}}
Sort dotted version strings like '1.10' and '1.2' in true numeric order (so '1.2' precedes '1.10'). Shorter versions that are a prefix of longer ones come first.

\noindent\texttt{sort\_versions(vs: list[str]) -> list[str]}
\begin{lstlisting}[style=ex]
>>> sort_versions(["1.10", "1.2", "2.0"])
['1.2', '1.10', '2.0']
>>> sort_versions(["0.9", "0.10.1", "0.10"])
['0.9', '0.10', '0.10.1']
>>> sort_versions(["3"])
['3']
\end{lstlisting}
\subsection*{9.9\quad Sort the score report\hfill\textnormal{\textit{apply}}}
A report arrives as newline-separated lines of 'name score'. Return the same report with lines ordered by score descending, ties broken by name ascending -- rebuilt with unlines, so the result ends with a newline.

\noindent\texttt{sort\_report(text: str) -> str}
\begin{lstlisting}[style=ex]
>>> sort_report("ann 50\nbob 70")
'bob 70\nann 50\n'
>>> sort_report("cy 70\nbob 70\nann 50")
'bob 70\ncy 70\nann 50\n'
>>> sort_report("solo 1")
'solo 1\n'
\end{lstlisting}
\subsection*{9.10\quad Combined guest list\hfill\textnormal{\textit{apply}}}
Two venues export their guest ids as sorted lists which may share entries. Produce one sorted list with each id appearing once. Merge first, then deduplicate -- both passes linear.

\noindent\texttt{merge\_unique(a: list[int], b: list[int]) -> list[int]}
\begin{lstlisting}[style=ex]
>>> merge_unique([1, 3, 5], [1, 2, 5])
[1, 2, 3, 5]
>>> merge_unique([], [4, 4])
[4]
>>> merge_unique([2], [])
[2]
\end{lstlisting}
\subsection*{9.11\quad Median of two sorted lists\hfill\textnormal{\textit{challenge}}}
Two sensors deliver their readings as sorted lists. Return the median of ALL readings as a float: the middle element of the combined order, or the mean of the two middles when the total count is even. At least one reading is guaranteed. A linear merge is accepted here; name the O(log(n+m)) partition approach in your complexity remarks.

\noindent\texttt{median\_sorted(a: list[int], b: list[int]) -> float}
\begin{lstlisting}[style=ex]
>>> median_sorted([1, 3], [2])
2.0
>>> median_sorted([1, 2], [3, 4])
2.5
>>> median_sorted([], [5])
5.0
>>> median_sorted([1, 1], [1])
1.0
\end{lstlisting}
\subsection*{9.12\quad Longest common prefix\hfill\textnormal{\textit{challenge}}}
Return the longest prefix shared by EVERY string in a list (empty string when the list is empty or shares nothing). Insight to find: after sorting, only the FIRST and LAST strings need comparing -- anything they share, every string between them shares too.

\noindent\texttt{common\_prefix(ws: list[str]) -> str}
\begin{lstlisting}[style=ex]
>>> common_prefix(["flower", "flow", "flight"])
'fl'
>>> common_prefix(["dog", "cat"])
''
>>> common_prefix([])
''
>>> common_prefix(["same"])
'same'
\end{lstlisting}
\clearpage
\chapter{Dict Folds and Bags}

Aggregation by key -- the fold aimed at a mapping. One family, one idea:
\texttt{insertWith} resolves a single collision, \texttt{fromListWith} a
list of them, \texttt{unionWith} two dicts' worth; and CHOOSING the combiner
is choosing a monoid -- addition counts, \texttt{old + new} groups,
\texttt{max} is bag union. Mind the Haskell order: the combiner hears the
NEW value first. \texttt{group}/\texttt{groupBy} handle the positional
cousin: runs, not piles.
\section{Concepts}
\subsection{\texttt{defaultdict}}\label{def:defaultdict}
a dict whose missing keys create themselves via a factory: touch d[k] and factory() is installed there. This "autovivification" makes grouping and counting one-liners. Gotcha: READING a missing key also plants it -- use 'in' or .get when you only want to look.
\begin{lstlisting}[style=ex]
>>> d = defaultdict(list)
>>> d['x'].append(1)
>>> d
{'x': [1]}
>>> d2 = defaultdict(int)
>>> d2['ghost']                         # a mere read...
0
>>> 'ghost' in d2                       # ...planted the key
True
\end{lstlisting}

How it works: one dunder does everything. When d[k] misses, Python calls d.\_\_missing\_\_(k); this subclass responds by calling the stored factory, installing the result, and returning it. The factory runs only on a miss -- and it takes NO arguments, which is why the leaf recipes in Tree/ITree are zero-argument lambdas.


Why it earns its place: it deletes the check-then-create dance from every grouping and counting loop. The write path becomes a bare append or +=, and the structure assembles itself on touch -- autovivification, which section 12's trie tools then apply recursively.


The discipline it demands: a mere READ also plants the key, so split your paths -- write with d[k], but ask questions with 'in' or .get. The prelude's bag operations read Counters with .get for exactly this reason.

\begin{lstlisting}[style=ex]
>>> d = defaultdict(int)
>>> d['hits'] += 1                      # write path: indexing is the point
>>> d.get('misses', 0), 'misses' in d   # read path: no ghost key planted
(0, False)
\end{lstlisting}

Remember it as: index to build, .get to ask.


Roll your own: subclass dict and implement one dunder -- \_\_missing\_\_, which Python calls when [k] fails. Call the factory, install the result, hand it back; iteration, repr and everything else is inherited. Note the factory takes no arguments, which is why leaf recipes are zero-argument lambdas.

\begin{lstlisting}[style=ex]
class defaultdict(dict):
    def __init__(self, factory=None, *args, **kw):
        super().__init__(*args, **kw)
        self.factory = factory
    def __missing__(self, key):
        if self.factory is None:
            raise KeyError(key)
        self[key] = self.factory()
        return self[key]
\end{lstlisting}
\subsection{\texttt{Counter}}\label{def:Counter}
count occurrences: a defaultdict(int) fold over the input. Missing keys read as 0 (mind the autovivification gotcha above); rank the .items() with sortOn.
\begin{lstlisting}[style=ex]
>>> c = Counter("abracadabra")
>>> c['a'], c['z']
(5, 0)
>>> sortOn(lambda kv: -kv[1], list(Counter("aab").items()))
[('a', 2), ('b', 1)]
\end{lstlisting}

Roll your own: a defaultdict(int) plus one fold -- d[x] += 1 per element. The += on a missing key is autovivification doing all the work.

\begin{lstlisting}[style=ex]
def Counter(xs):  # count-by-value; a defaultdict(int) fold
    d = defaultdict(int)
    for x in xs:
        d[x] += 1
    return d
\end{lstlisting}
\subsection{\texttt{insertWith}}\label{def:insertWith}
Data.Map's insertWith f k v m: insert (k, v), combining with f(new, old) on collision -- Haskell's argument order, new value first. PERSISTENT, also like Haskell: it returns a fresh dict and leaves the argument untouched.
\begin{lstlisting}[style=ex]
>>> insertWith(add, 'x', 4, {'x': 1})
{'x': 5}
>>> d = {'y': 2}
>>> insertWith(add, 'x', 4, d)
{'y': 2, 'x': 4}
>>> d                                    # the original survives
{'y': 2}
\end{lstlisting}

Roll your own: one dict display -- {**d, k: ...}: the splat copies d (that IS the persistence), and the new key's value is f(v, d[k]) on collision (new value first, Haskell's order) or plain v. Compare fromListWith, which applies the same rule but mutates, because there the dict is its own private accumulator.

\begin{lstlisting}[style=ex]
insertWith = lambda f, k, v, d: {**d, k: f(v, d[k]) if k in d else v}  # Data.Map insertWith: PERSISTENT
  # (fresh dict); f(new, old)
\end{lstlisting}
\subsection{\texttt{fromListWith}}\label{def:fromListWith}
THE dict-building fold: pour (key, value) pairs into a dict, combining collisions with f(new, old) -- Haskell's order, like insertWith. Keys keep first-seen order.
\begin{lstlisting}[style=ex]
>>> fromListWith(add, [(c, 1) for c in "aab"])
{'a': 2, 'b': 1}
>>> fromListWith(lambda new, old: old + new, [('a', [1]), ('b', [2]), ('a', [3])])
{'a': [1, 3], 'b': [2]}
\end{lstlisting}

The order gotcha, faithfully Haskell's: because f receives (new, old), grouping with a bare concat REVERSES each group -- the classic Data.Map surprise. Combine as old + new to keep arrival order, or use plain addition where order cannot matter (counts).

\begin{lstlisting}[style=ex]
>>> fromListWith(add, [(len(w), [w]) for w in ["hi", "ox", "sun"]])
{2: ['ox', 'hi'], 3: ['sun']}
>>> fromListWith(lambda new, old: old + new, [(len(w), [w]) for w in ["hi", "ox", "sun"]])
{2: ['hi', 'ox'], 3: ['sun']}
\end{lstlisting}

Why it earns its place -- one function, three monoids: choosing f is choosing what the dict MEANS. Addition over (k, 1) pairs counts; old + new over (k, [v]) pairs groups; max merges bags. That is the whole family: insertWith is one collision, fromListWith a list of them, unionWith two dicts' worth (where f receives (a's value, b's value)).


The recognition cue: whenever a loop builds a dict with an if-else on key existence, you are hand-rolling this fold. Name the pairs, name the monoid, and the loop disappears.


Remember it as: grouping, counting, and merging are one fold -- whose combiner hears the NEW value first.


Roll your own: an empty dict and one loop applying the collision rule per pair -- d[k] = f(v, d[k]) if the key exists, else v. It is insertWith inlined and IN-PLACE: mutation is fine here because the dict being built is the fold's own accumulator, invisible until returned.

\begin{lstlisting}[style=ex]
def fromListWith(f, pairs):  # Data.Map fromListWith: THE dict-building fold; f(new, old)
    d = {}  # like insertWith -- so grouping with concat REVERSES each
    for k, v in pairs:  # group (the classic Haskell gotcha); use add for counts
        d[k] = f(v, d[k]) if k in d else v
    return d  # Counter == fromListWith(add) over (x, 1);
  # grouping == fromListWith(++) over (k, [v])
\end{lstlisting}
\noindent{\small\itshape Used before it is rolled: \texttt{group} in \S\ref{def:group} (p.~\pageref{def:group}).}

\subsection{\texttt{unionWith}}\label{def:unionWith}
merge two dicts, combining SHARED keys with f; a's keys keep their order, b's new keys follow. Choosing f is choosing a monoid: max gives bag union, + gives a summing merge.
\begin{lstlisting}[style=ex]
>>> unionWith(max, {'a': 1, 'b': 5}, {'b': 2, 'c': 3})
{'a': 1, 'b': 5, 'c': 3}
>>> unionWith(add, {'a': 1, 'b': 2}, {'b': 3, 'c': 4})
{'a': 1, 'b': 5, 'c': 4}
\end{lstlisting}

Roll your own: two splats -- {**a, **{...}}: a copied whole, then overlaid with b's items after each has been through the collision rule, f(a[k], v) when k is shared, plain v when it is new. The asymmetry is deliberate: a's keys keep their positions, and only keys new to b append at the end.

\begin{lstlisting}[style=ex]
unionWith = lambda f, a, b: {**a, **{k: f(a[k], v) if k in a else v for k, v in b.items()}}
  # Data.Map unionWith: f(a's value, b's value) on shared keys;
  # bag_union == unionWith(max), merge-sum == unionWith(add)
\end{lstlisting}
\noindent{\small\itshape Used before it is rolled: \texttt{bag\_union} in \S\ref{def:bag-union} (p.~\pageref{def:bag-union}), \texttt{on} in \S\ref{def:on} (p.~\pageref{def:on}).}

\subsection{\texttt{groupBy}}\label{def:groupBy}
runs by an equivalence predicate; each newcomer is compared with eq against the run's FIRST element (span-based, exactly Haskell's groupBy), not against its neighbour -- a subtle and deliberate fidelity.
\begin{lstlisting}[style=ex]
>>> groupBy(lambda a, b: a == b, [1, 1, 2])
[[1, 1], [2]]
>>> groupBy(lambda a, b: b - a <= 1, [1, 2, 3, 10, 11])
[[1, 2], [3], [10, 11]]
\end{lstlisting}

The word that matters is CONSECUTIVE: a new run starts whenever eq against the run's first element fails, so 'abab' yields four runs. Not a weakness -- a different question: runs, streaks and compression care about position, and these are the only tools asking about it.

\begin{lstlisting}[style=ex]
>>> [(run[0], len(run)) for run in group("aaabbbbcc")]     # run lengths
[('a', 3), ('b', 4), ('c', 2)]
\end{lstlisting}

For CATEGORIES -- position irrelevant -- either sort first, so equals become adjacent, or use fromListWith, which never cared about position at all. Choosing between them is stating what your key means.

\begin{lstlisting}[style=ex]
>>> [(run[0], len(run)) for run in group(sorted("abab"))]
[('a', 2), ('b', 2)]
\end{lstlisting}

Remember it as: group answers "how long are the runs?"; fromListWith answers "what are the piles?".


Roll your own: one pass in which each element either joins the current run or starts a fresh one. The test is eq(out[-1][0], x) -- against the run's FIRST element, not its last -- which is exactly how Haskell's span-based groupBy behaves, and the detail that makes non-transitive predicates (like "within 1 of") agree with it.

\begin{lstlisting}[style=ex]
def groupBy(eq, xs):  # Data.List groupBy: runs of CONSECUTIVE elements; each new
    out = []  # element is compared via eq against the run's FIRST element,
    for x in xs:  # exactly like Haskell (span-based) -- not its neighbor
        if out and eq(out[-1][0], x):
            out[-1].append(x)
        else:
            out.append([x])
    return out
\end{lstlisting}
\noindent{\small\itshape Used before it is rolled: \texttt{group} in \S\ref{def:group} (p.~\pageref{def:group}).}

\subsection{\texttt{group}}\label{def:group}
runs of CONSECUTIVE equal elements, as bare lists ([[a]], no keys) -- Data.List's group.
\begin{lstlisting}[style=ex]
>>> group("aaabbc")
[['a', 'a', 'a'], ['b', 'b'], ['c']]
>>> group("abab")
[['a'], ['b'], ['a'], ['b']]
\end{lstlisting}

Roll your own --- the definition:

\begin{lstlisting}[style=ex]
group = lambda xs: groupBy(lambda a, b: a == b, xs)  # Data.List group: runs of equals, [[a]] -- no keys
\end{lstlisting}
\subsection{\texttt{bag\_union}}\label{def:bag-union}
multisets (bags) are Counters; union takes the LARGER count per key. Key order follows set union, so compare with == rather than printing.
\begin{lstlisting}[style=ex]
>>> bag_union(Counter("aab"), Counter("abb")) == {'a': 2, 'b': 2}
True
\end{lstlisting}

Roll your own --- the definition:

\begin{lstlisting}[style=ex]
bag_union = lambda a, b: {k: max(a.get(k, 0), b.get(k, 0)) for k in a.keys() | b.keys()}
\end{lstlisting}
\subsection{\texttt{bag\_inter}}\label{def:bag-inter}
intersection: the SMALLER count per shared key, zero counts dropped. Note the implementation reads with .get, never [] -- indexing a Counter would autovivify zeros into it.
\begin{lstlisting}[style=ex]
>>> bag_inter(Counter("aab"), Counter("ab")) == {'a': 1, 'b': 1}
True
>>> bag_inter(Counter("aa"), Counter("bb")) == {}
True
\end{lstlisting}

Roll your own --- the definition:

\begin{lstlisting}[style=ex]
bag_inter = lambda a, b: {k: v for k in a.keys() & b.keys() if (v := min(a.get(k, 0), b.get(k, 0)))}
\end{lstlisting}
\subsection{\texttt{bag\_diff}}\label{def:bag-diff}
difference, clipped at zero: what remains of a after removing b.
\begin{lstlisting}[style=ex]
>>> bag_diff(Counter("aab"), Counter("ab")) == {'a': 1}
True
\end{lstlisting}

Roll your own --- the definition:

\begin{lstlisting}[style=ex]
bag_diff = lambda a, b: {k: a[k] - b[k] for k in a if a[k] > b.get(k, 0)}  # clipped at 0
\end{lstlisting}
\subsection{\texttt{bag\_sub}}\label{def:bag-sub}
inclusion: is bag a contained in bag b, multiplicities included? The ransom-note test.
\begin{lstlisting}[style=ex]
>>> bag_sub(Counter("ab"), Counter("aabb"))
True
>>> bag_sub(Counter("aab"), Counter("ab"))
False
\end{lstlisting}

Roll your own --- the definition:

\begin{lstlisting}[style=ex]
bag_sub = lambda a, b: all(b.get(k, 0) >= n for k, n in a.items())  # a <= b (inclusion)
\end{lstlisting}
\section{Homework}
Work in order: drills cement each tool, applies combine them,
challenges are interview-grade. Write each solution in a scratch file with
the given doctests pasted in, and let \texttt{python3 -m doctest} referee.
\subsection*{10.1\quad Vowel census\hfill\textnormal{\textit{drill}}}
Count the frequency of each vowel (a, e, i, o, u) in a string, case-insensitive. Return the counts as (vowel, count) pairs sorted alphabetically; vowels that never occur are omitted.

\noindent\texttt{vowel\_census(s: str) -> list[tuple[str, int]]}
\begin{lstlisting}[style=ex]
>>> vowel_census("Banana")
[('a', 3)]
>>> vowel_census("queue")
[('e', 2), ('u', 2)]
>>> vowel_census("xyz")
[]
\end{lstlisting}
\subsection*{10.2\quad Order totals\hfill\textnormal{\textit{drill}}}
An order log is a list of (product, quantity) pairs; products repeat. Total the quantity per product with one fromListWith fold. Keys keep first-seen order.

\noindent\texttt{totals(pairs: list[tuple[str, int]]) -> dict[str, int]}
\begin{lstlisting}[style=ex]
>>> totals([("apple", 2), ("pear", 1), ("apple", 3)])
{'apple': 5, 'pear': 1}
>>> totals([])
{}
>>> totals([("kiwi", 4)])
{'kiwi': 4}
\end{lstlisting}
\subsection*{10.3\quad Longest run\hfill\textnormal{\textit{drill}}}
Return the length of the longest run of consecutive equal characters in a string (0 for the empty string). group hands you the runs; you measure them.

\noindent\texttt{longest\_run(s: str) -> int}
\begin{lstlisting}[style=ex]
>>> longest_run("aaabb")
3
>>> longest_run("abab")
1
>>> longest_run("")
0
\end{lstlisting}
\subsection*{10.4\quad Restock, persistently\hfill\textnormal{\textit{drill}}}
Given a warehouse inventory dict and one delivery (item, quantity), return a NEW inventory with the delivery added to that item's count -- and leave the original dict untouched, which is exactly insertWith's contract. New items simply appear.

\noindent\texttt{restock(inv: dict[str, int], item: str, qty: int) -> dict[str, int]}
\begin{lstlisting}[style=ex]
>>> inv = {'nail': 10}
>>> restock(inv, 'nail', 5)
{'nail': 15}
>>> inv
{'nail': 10}
>>> restock(inv, 'screw', 3)
{'nail': 10, 'screw': 3}
\end{lstlisting}
\subsection*{10.5\quad Tile check\hfill\textnormal{\textit{drill}}}
Scrabble-style: can a word be assembled from a hand of letter tiles, respecting multiplicity (two o's in the word need two o tiles)? This is multiset inclusion -- one bag\_sub.

\noindent\texttt{can\_build(word: str, tiles: str) -> bool}
\begin{lstlisting}[style=ex]
>>> can_build("cat", "tacos")
True
>>> can_build("book", "bok")
False
>>> can_build("", "xyz")
True
\end{lstlisting}
\subsection*{10.6\quad High-water marks\hfill\textnormal{\textit{apply}}}
Two flood sensors each report per-zone water maxima as a dict. Combine them into one report keeping the HIGHER reading for zones both saw; zones only one sensor saw pass through unchanged. Choose the right combining function for unionWith -- the monoid is the design decision.

\noindent\texttt{high\_water(a: dict[str, int], b: dict[str, int]) -> dict[str, int]}
\begin{lstlisting}[style=ex]
>>> high_water({'north': 3, 'south': 7}, {'south': 9, 'east': 2})
{'north': 3, 'south': 9, 'east': 2}
>>> high_water({}, {'west': 1})
{'west': 1}
>>> high_water({'north': 4}, {})
{'north': 4}
\end{lstlisting}
\subsection*{10.7\quad Click sessions\hfill\textnormal{\textit{apply}}}
A user's click timestamps arrive sorted. Clicks belong to one session while they fall within 30 seconds OF THE SESSION'S FIRST CLICK; a later click starts a new session. Split the timestamps into sessions -- note this is groupBy's compare-against-the-run's-FIRST semantics, verbatim, not a neighbor-gap rule.

\noindent\texttt{sessions(ts: list[int]) -> list[list[int]]}
\begin{lstlisting}[style=ex]
>>> sessions([0, 10, 25, 50, 60, 100])
[[0, 10, 25], [50, 60], [100]]
>>> sessions([])
[]
>>> sessions([5])
[[5]]
\end{lstlisting}
\subsection*{10.8\quad Department roster\hfill\textnormal{\textit{apply}}}
HR exports (department, employee) rows in hiring order. Build a dict mapping each department to its employees IN HIRING ORDER. Mind fromListWith's f(new, old) argument order: a bare concat would reverse each roster.

\noindent\texttt{roster(rows: list[tuple[str, str]]) -> dict[str, list[str]]}
\begin{lstlisting}[style=ex]
>>> roster([("eng", "ann"), ("ops", "bob"), ("eng", "cy")])
{'eng': ['ann', 'cy'], 'ops': ['bob']}
>>> roster([])
{}
>>> roster([("qa", "dee")])
{'qa': ['dee']}
\end{lstlisting}
\subsection*{10.9\quad Shopping list\hfill\textnormal{\textit{apply}}}
A recipe lists its ingredients with repeats ('egg' twice means two eggs); your pantry is a list of what you have. Produce the shopping list -- what is still missing and how many of each -- as (item, count) pairs sorted alphabetically. Multiset difference, clipped at zero.

\noindent\texttt{shopping\_list(need: list[str], have: list[str]) -> list[tuple[str, int]]}
\begin{lstlisting}[style=ex]
>>> shopping_list(["egg", "egg", "milk"], ["egg"])
[('egg', 1), ('milk', 1)]
>>> shopping_list(["jam"], ["jam", "jam"])
[]
>>> shopping_list([], ["salt"])
[]
\end{lstlisting}
\subsection*{10.10\quad Traffic rollup\hfill\textnormal{\textit{apply}}}
Each server in a fleet reports a dict of per-page visit counts. Roll the whole fleet up into one dict of grand totals. This is a FOLD whose combining step is itself a unionWith -- dict-merging as a monoid.

\noindent\texttt{rollup(reports: list[dict[str, int]]) -> dict[str, int]}
\begin{lstlisting}[style=ex]
>>> rollup([{'home': 1}, {'home': 2, 'faq': 5}, {'faq': 1}])
{'home': 3, 'faq': 6}
>>> rollup([])
{}
>>> rollup([{'a': 1}])
{'a': 1}
\end{lstlisting}
\subsection*{10.11\quad Majority element\hfill\textnormal{\textit{challenge}}}
An election log is a list of votes. The most common candidate wins only with a STRICT majority (more than half of all votes): return that candidate, or None when nobody qualifies. Ties for the top count must not invent a majority. Count with a bag; interrogate it deterministically.

\noindent\texttt{majority(xs: list) -> value | None}
\begin{lstlisting}[style=ex]
>>> majority([1, 2, 1, 1])
1
>>> majority([1, 2]) is None
True
>>> majority([]) is None
True
>>> majority([5])
5
\end{lstlisting}
\subsection*{10.12\quad Top spender per city\hfill\textnormal{\textit{challenge}}}
A payment feed yields (city, customer, amount) rows; customers repeat. For every city, find the customer with the highest TOTAL spend there; break ties by taking the alphabetically first name. Return {city: customer} with cities in first-appearance order. Two stacked dict folds: totals per (city, customer), then a per-city best.

\noindent\texttt{top\_spender(rows: list[tuple[str, str, int]]) -> dict[str, str]}
\begin{lstlisting}[style=ex]
>>> top_spender([("nyc", "ann", 50), ("nyc", "bob", 70), ("la", "cy", 20),
...              ("nyc", "ann", 30)])
{'nyc': 'ann', 'la': 'cy'}
>>> top_spender([("sf", "a", 10), ("sf", "b", 10)])
{'sf': 'a'}
>>> top_spender([])
{}
\end{lstlisting}
\clearpage
\chapter{Trees and Tries}

Hierarchy as nested dicts that build themselves. \texttt{ITree} autovivifies
a branch the moment you touch it -- a trie in one line -- and
\texttt{Tree(depth, leaf)} bottoms the same idea out at fixed depth.
Writing is indexing or \texttt{setpath}; asking is ALWAYS \texttt{getpath}
(reads plant ghost branches otherwise); and \texttt{paths} folds the whole
tree back into rows, which is where the value gets harvested.
\section{Concepts}
\subsection{\texttt{Tree}}\label{def:Tree}
nested defaultdicts to a FIXED depth, with a leaf factory at the bottom. Tree(1, list) is the grouping dict; Tree(2, int) is a sparse 2-D counting table. Indexing autovivifies the whole path.
\begin{lstlisting}[style=ex]
>>> t = Tree(1, list)
>>> t['evens'].append(2)
>>> t
{'evens': [2]}
>>> m = Tree(2, int)
>>> m['a']['b'] += 1
>>> m
{'a': {'b': 1}}
\end{lstlisting}

Roll your own --- the definition:

\begin{lstlisting}[style=ex]
Tree = lambda depth, leaf: (defaultdict(leaf) if depth == 1
                             else defaultdict(lambda: Tree(depth - 1, leaf)))
\end{lstlisting}
\subsection{\texttt{ITree}}\label{def:ITree}
the INFINITE tree: every branch is another ITree, so any depth autovivifies on touch. This is the trie: index letter by letter and the branch builds itself.
\begin{lstlisting}[style=ex]
>>> t = ITree()
>>> t['x']['y']['z'] = 9
>>> t
{'x': {'y': {'z': 9}}}
\end{lstlisting}

How it works: the definition is a knot -- ITree = lambda: defaultdict(ITree). Every missing key manufactures another ITree, which will do the same, without end. Indexing therefore GROWS the tree: touching t['c']['a']['t'] conjures the whole branch. A trie in one line, built from nothing but \_\_missing\_\_.

\begin{lstlisting}[style=ex]
>>> t = ITree()
>>> for w in ["cat", "car"]: setpath(t, list(w) + ['$'], w)
>>> sorted([leaf for p, leaf in paths(t) if p[-1] == '$'])
['car', 'cat']
\end{lstlisting}

Why it earns its place: hierarchies -- tries, directory trees, nested groupings -- stop needing setup code. Writing is indexing (or setpath, which also plants a leaf); harvesting is paths. The one edged rule: reads autovivify too, so QUESTIONS go through getpath, never bare indexing -- getpath's entry demonstrates the ghost branch this avoids.


Versus Tree(depth, leaf): Tree is the finite cousin -- fixed depth, with a real leaf factory (list, int) at the bottom instead of more tree. Tree(1, list) is the grouping dict; ITree is for depth you cannot know in advance, like words.


Remember it as: a dict that dreams up more of itself on demand.


Roll your own: ITree = lambda: defaultdict(ITree) mentions itself -- legal, because the lambda body only runs on a MISS, by which time the name is bound. Tree bottoms the same recursion out at a fixed depth by handing the last level a real leaf factory instead of more tree.

\begin{lstlisting}[style=ex]
ITree = lambda: defaultdict(ITree)
\end{lstlisting}
\noindent{\small\itshape Used before it is rolled: \texttt{paths} in \S\ref{def:paths} (p.~\pageref{def:paths}), \texttt{setpath} in \S\ref{def:setpath} (p.~\pageref{def:setpath}).}

\subsection{\texttt{paths}}\label{def:paths}
fold a whole Tree/ITree into a flat list of (keypath, leaf) pairs, in insertion order. The read- everything companion to setpath's write-one-path.
\begin{lstlisting}[style=ex]
>>> t = ITree()
>>> setpath(t, ['a', 'b'], 1); setpath(t, ['a', 'c'], 2)
>>> paths(t)
[(['a', 'b'], 1), (['a', 'c'], 2)]
\end{lstlisting}

How it works: structural recursion with two arms, exactly the case-arms-as-fold story. A non-dict (or empty node) is a LEAF: report the single row ([], leaf). A dict is a BRANCH: recurse into every subtree and prepend the branching key to each keypath that comes back. The whole tree flattens into (keypath, leaf) rows, insertion-ordered, losslessly.


Why it earns its place: building a trie is the easy half; the value is HARVESTING it. autocomplete filters rows under a prefix, deepest\_directory takes maxOn of keypath length, leaves reads the second column. One cata, many harvests -- you never write tree-walking code again.

\begin{lstlisting}[style=ex]
>>> t = ITree()
>>> setpath(t, ['a', 'b'], 1); setpath(t, ['c'], 2)
>>> paths(t)
[(['a', 'b'], 1), (['c'], 2)]
>>> maxOn(lambda row: len(row[0]), paths(t))
(['a', 'b'], 1)
\end{lstlisting}

Remember it as: paths turns a tree back into rows; setpath turns rows into a tree.


Roll your own: recursion on one question -- is this still a dict? A leaf reports a single row with an empty path; a branch recurses into every child and conses its key onto every path that comes back. Depth equals tree height, comfortably inside Python's limit for sane tries.

\begin{lstlisting}[style=ex]
def paths(t):  # cata for Tree/ITree: [(keypath, leaf)]
  # recursion depth = tree height (word length / len(nums)) -- well under the ~1000 limit
    if not isinstance(t, dict) or not t:  # leaf = non-dict value, or empty node
        return [([], t)]
    return [([k] + p, leaf) for k, sub in t.items() for p, leaf in paths(sub)]
\end{lstlisting}
\noindent{\small\itshape Used before it is rolled: \texttt{setpath} in \S\ref{def:setpath} (p.~\pageref{def:setpath}).}

\subsection{\texttt{leaves}}\label{def:leaves}
just the leaf values of paths.
\begin{lstlisting}[style=ex]
>>> t = ITree()
>>> setpath(t, ['a'], 1); setpath(t, ['b'], 2)
>>> leaves(t)
[1, 2]
\end{lstlisting}

Roll your own --- the definition:

\begin{lstlisting}[style=ex]
leaves = lambda t: [l for _, l in paths(t)]
\end{lstlisting}
\noindent{\small\itshape Used before it is rolled: \texttt{setpath} in \S\ref{def:setpath} (p.~\pageref{def:setpath}).}

\subsection{\texttt{setpath}}\label{def:setpath}
write a leaf value at the end of a key path, autovivifying the branch on the way. Caution: writing at an interior key CLOBBERS the subtree below it.
\begin{lstlisting}[style=ex]
>>> t = ITree()
>>> setpath(t, ['a', 'b'], 7)
>>> t
{'a': {'b': 7}}
>>> setpath(t, ['a'], 0)                # clobbers {'b': 7}
>>> t
{'a': 0}
\end{lstlisting}

Roll your own: walk all but the last key by PLAIN INDEXING -- on an ITree that autovivifies the branch as you go -- then assign at the final key, the single write. getpath is the mirror image with the safety catch on: isinstance and 'in' checks instead of indexing, so reading plants nothing.

\begin{lstlisting}[style=ex]
def setpath(t, ks, v):  # write leaf v at key path ks (autovivifies interior)
    for k in ks[:-1]:
        t = t[k]
    t[ks[-1]] = v  # NB: clobbers any subtree already at ks
\end{lstlisting}
\subsection{\texttt{getpath}}\label{def:getpath}
read along a key path WITHOUT autovivifying, returning a default when the path is absent. This is why it exists: reading a trie with [] would plant ghost branches.
\begin{lstlisting}[style=ex]
>>> t = ITree()
>>> setpath(t, ['a', 'b'], 7)
>>> getpath(t, ['a', 'b']), getpath(t, ['a', 'z'], default=0)
(7, 0)
>>> 'z' in t['a']                       # no ghost branch was planted
False
\end{lstlisting}

Roll your own --- the definition:

\begin{lstlisting}[style=ex]
def getpath(t, ks, default=None):  # read along ks WITHOUT autovivifying
    for k in ks:
        if not isinstance(t, dict) or k not in t:
            return default
        t = t[k]
    return t
\end{lstlisting}
\section{Homework}
Work in order: drills cement each tool, applies combine them,
challenges are interview-grade. Write each solution in a scratch file with
the given doctests pasted in, and let \texttt{python3 -m doctest} referee.
\subsection*{11.1\quad Plant a settings tree\hfill\textnormal{\textit{drill}}}
Write set\_all(pairs), where each pair is (key\_path, value) with key\_path a list of keys. Plant every value in a fresh ITree with setpath and return the tree. Later pairs overwrite earlier ones at the same path.

\noindent\texttt{set\_all(pairs: list[tuple[list, object]]) -> dict}
\begin{lstlisting}[style=ex]
>>> set_all([(['a', 'b'], 1), (['c'], 2)])
{'a': {'b': 1}, 'c': 2}
>>> set_all([(['x'], 1), (['x'], 9)])
{'x': 9}
>>> set_all([])
{}
\end{lstlisting}
\subsection*{11.2\quad Dotted config lookup\hfill\textnormal{\textit{drill}}}
Write get\_setting(cfg, dotted, default): read a nested dict with a dotted key string like 'server.tls.on'. Missing keys -- or a path that runs THROUGH a non-dict leaf -- answer the default, and the read must never plant ghost keys.

\noindent\texttt{get\_setting(cfg: dict, dotted: str, default) -> object}
\begin{lstlisting}[style=ex]
>>> cfg = {'server': {'port': 8080, 'tls': {'on': True}}}
>>> get_setting(cfg, 'server.port', 0)
8080
>>> get_setting(cfg, 'server.tls.on', False)
True
>>> get_setting(cfg, 'server.host', 'localhost')
'localhost'
>>> get_setting(cfg, 'server.port.max', -1)
-1
\end{lstlisting}
\noindent{\small\itshape Rolled later: \texttt{on} in \S\ref{def:on} (p.~\pageref{def:on}).}

\subsection*{11.3\quad Every leaf value\hfill\textnormal{\textit{drill}}}
Write leaf\_values(t): all leaf values of a nested dict, in insertion order. Recognise the prelude tool before writing any recursion.

\noindent\texttt{leaf\_values(t: dict) -> list}
\begin{lstlisting}[style=ex]
>>> leaf_values({'a': {'b': 1, 'c': 2}, 'd': 3})
[1, 2, 3]
>>> leaf_values({'x': 7})
[7]
>>> leaf_values({'a': {'b': {'c': 'deep'}}})
['deep']
\end{lstlisting}
\subsection*{11.4\quad Two-level tally\hfill\textnormal{\textit{drill}}}
Write sales\_table(pairs): count (city, product) sale events into a two-level table -- table[city][product] is the number of sales. Use Tree(2, int) so no key ever needs checking.

\noindent\texttt{sales\_table(pairs: list[tuple[str, str]]) -> dict}
\begin{lstlisting}[style=ex]
>>> sales_table([('nyc', 'ai'), ('nyc', 'ai'), ('la', 'uav')])
{'nyc': {'ai': 2}, 'la': {'uav': 1}}
>>> sales_table([])
{}
>>> sales_table([('x', 'y')])['x']['y']
1
\end{lstlisting}
\subsection*{11.5\quad Slash-joined key paths\hfill\textnormal{\textit{drill}}}
Write key\_paths(t): the key path of every leaf in a nested dict (string keys), joined with '/', in insertion order.

\noindent\texttt{key\_paths(t: dict) -> list[str]}
\begin{lstlisting}[style=ex]
>>> key_paths({'a': {'b': 1}, 'c': 2})
['a/b', 'c']
>>> key_paths({'solo': 9})
['solo']
>>> key_paths({'a': {'b': {'c': 0}}})
['a/b/c']
\end{lstlisting}
\subsection*{11.6\quad Phone prefix counter\hfill\textnormal{\textit{apply}}}
A call-routing box holds a set of DISTINCT phone numbers (digit strings). Write count\_with\_prefix(numbers, prefix): how many numbers start with the given prefix. Build a digit trie once ('\$'-terminated so a full number survives being a prefix of another), descend to the prefix subtree with getpath, and count the '\$' leaves below it.

\noindent\texttt{count\_with\_prefix(numbers: list[str], prefix: str) -> int}
\begin{lstlisting}[style=ex]
>>> nums = ['5551234', '5555678', '5551299']
>>> count_with_prefix(nums, '555')
3
>>> count_with_prefix(nums, '5551')
2
>>> count_with_prefix(nums, '9')
0
>>> count_with_prefix(['12', '123'], '12')
2
\end{lstlisting}
\subsection*{11.7\quad Deep config merge\hfill\textnormal{\textit{apply}}}
Deployment configs are nested dicts. Write deep\_merge(a, b): b's settings override a's, except where BOTH sides hold dicts -- then merge recursively. Return a new dict; neither argument may be mutated at any level you touch. Key order: a's keys first, then b's new ones.

\noindent\texttt{deep\_merge(a: dict, b: dict) -> dict}
\begin{lstlisting}[style=ex]
>>> deep_merge({'a': {'x': 1}, 'b': 2}, {'a': {'y': 9}, 'c': 3})
{'a': {'x': 1, 'y': 9}, 'b': 2, 'c': 3}
>>> deep_merge({'a': {'x': 1}}, {'a': 5})
{'a': 5}
>>> deep_merge({}, {'k': 1})
{'k': 1}
>>> orig = {'a': {'x': 1}}
>>> _ = deep_merge(orig, {'a': {'y': 2}})
>>> orig
{'a': {'x': 1}}
\end{lstlisting}
\subsection*{11.8\quad Folder rollup\hfill\textnormal{\textit{apply}}}
Given (path, size) pairs like ('a/b/f.txt', 10), write folder\_size(files, folder): the total size under a folder such as 'a/b' (or the size of a single file if the path names one). Unknown paths answer 0. Build the tree with setpath, descend with getpath, sum the leaves.

\noindent\texttt{folder\_size(files: list[tuple[str, int]], folder: str) -> int}
\begin{lstlisting}[style=ex]
>>> fs = [('a/b/f.txt', 10), ('a/b/g.txt', 5), ('a/c/h.txt', 7), ('d.txt', 1)]
>>> folder_size(fs, 'a/b')
15
>>> folder_size(fs, 'a')
22
>>> folder_size(fs, 'zzz')
0
>>> folder_size(fs, 'a/b/f.txt')
10
\end{lstlisting}
\subsection*{11.9\quad Word index, two keys deep\hfill\textnormal{\textit{apply}}}
Write word\_index(words): index words first by initial letter, then by length -- index[letter][length] is the list of matching words in arrival order. Words are non-empty and lowercase.

\noindent\texttt{word\_index(words: list[str]) -> dict}
\begin{lstlisting}[style=ex]
>>> word_index(['ant', 'ape', 'bee', 'at'])
{'a': {3: ['ant', 'ape'], 2: ['at']}, 'b': {3: ['bee']}}
>>> word_index([])
{}
>>> word_index(['x'])['x'][1]
['x']
\end{lstlisting}
\subsection*{11.10\quad Flatten to dotted lines\hfill\textnormal{\textit{apply}}}
The inverse of a config tree: write dotted(t) producing sorted 'path.to.key=value' strings for every leaf (str() the value). An empty tree yields []. Keys are strings without dots.

\noindent\texttt{dotted(t: dict) -> list[str]}
\begin{lstlisting}[style=ex]
>>> dotted({'srv': {'port': 8080, 'tls': {'on': True}}, 'debug': False})
['debug=False', 'srv.port=8080', 'srv.tls.on=True']
>>> dotted({})
[]
>>> dotted({'a': 1})
['a=1']
\end{lstlisting}
\noindent{\small\itshape Rolled later: \texttt{on} in \S\ref{def:on} (p.~\pageref{def:on}).}

\subsection*{11.11\quad Longest common prefix via trie\hfill\textnormal{\textit{challenge}}}
Write common\_prefix(words): the longest string that every word starts with. Build a '\$'-terminated trie, then walk from the root while the current node has EXACTLY one child and that child is not the terminator. [] answers ''. Classic interview question -- the trie walk is the O(total chars) answer that scales to millions of words.

\noindent\texttt{common\_prefix(words: list[str]) -> str}
\begin{lstlisting}[style=ex]
>>> common_prefix(['flower', 'flow', 'flight'])
'fl'
>>> common_prefix(['car', 'card'])
'car'
>>> common_prefix(['abc'])
'abc'
>>> common_prefix(['dog', 'race'])
''
>>> common_prefix([])
''
\end{lstlisting}
\subsection*{11.12\quad Render the directory listing\hfill\textnormal{\textit{challenge}}}
Given directory paths like 'a/b' (parents may or may not be listed separately), write tree\_lines(dirs): the indented listing -- each name on its own line, two spaces per depth, children in first-mention order. Build by folding indexing over each path's parts (walking IS building on an ITree), then emit with your own recursion over items().

\noindent\texttt{tree\_lines(dirs: list[str]) -> list[str]}
\begin{lstlisting}[style=ex]
>>> tree_lines(['a/b', 'a/c', 'd'])
['a', '  b', '  c', 'd']
>>> tree_lines(['x/y/z'])
['x', '  y', '    z']
>>> tree_lines([])
[]
>>> tree_lines(['m', 'm/n'])
['m', '  n']
\end{lstlisting}
\clearpage
\chapter{Machines: Deque, DSU, Heap}

Three data structures rebuilt from nothing, each carrying one great argument:
the two-stack \texttt{deque}'s amortised $O(1)$ ends (each element migrates
once); union--find's path halving making connectivity effectively constant
-- with \texttt{union} returning False AS the cycle detector; and the heap's
single invariant (parent beats children) keeping the minimum readable at
\texttt{h[0]} forever. Owning these means never being helpless without
imports.
\section{Concepts}
\subsection{\texttt{deque}}\label{def:deque}
a double-ended queue built from two stacks, all four end operations amortised O(1) (elements migrate between stacks only when one runs dry). This is the BFS frontier; a plain list's pop(0) is O(n).
\begin{lstlisting}[style=ex]
>>> d = deque([1, 2, 3])
>>> d.appendleft(0); d.append(4)
>>> d.popleft(), d.pop()
(0, 4)
>>> len(d), list(d)
(3, [1, 2, 3])
\end{lstlisting}

How it works: two plain lists, nose to nose -- \_out holds the front (reversed), \_in holds the back. append pushes onto \_in, appendleft onto \_out, both O(1). The trick is popping an EMPTY side: the other list is reversed across in one O(n) migration, and popping is O(1) again.


Why that still counts as O(1): amortization. Each element migrates at most once between being pushed and popped, so n operations cost O(n) in total -- constant per operation on average, which is what the workload's big-O cares about. A bare Python list cannot say this: pop(0) shifts every element, every single time.

\begin{lstlisting}[style=ex]
>>> q = deque()
>>> for x in [1, 2, 3]: q.append(x)
>>> q.popleft(), q.popleft()            # first pop pays the migration; the next is free
(1, 2)
\end{lstlisting}

Why it earns its place: BFS. The frontier must yield its OLDEST member (popleft) while newcomers join at the back (append) -- levelorder and grid\_path are this class doing its one job.


Remember it as: two stacks facing each other make a queue; the reversal bill is paid once per element.


Roll your own: two plain lists nose to nose, each end appending and popping its own list at O(1). All the cleverness is in the pop guard: when your side is empty, reverse the OTHER list across in one go. Each element makes that crossing at most once between entering and leaving, so the occasional O(n) reversal amortises to O(1) per operation.

\begin{lstlisting}[style=ex]
class deque:  # two-stack; all four ends amortized O(1)
    def __init__(self, xs=()):
        self._in, self._out = [], list(reversed(list(xs)))
    def append(self, x): self._in.append(x)
    def appendleft(self, x): self._out.append(x)
    def pop(self):
        if not self._in:
            self._in, self._out = list(reversed(self._out)), []
        return self._in.pop()
    def popleft(self):
        if not self._out:
            self._out, self._in = list(reversed(self._in)), []
        return self._out.pop()
    def __len__(self): return len(self._in) + len(self._out)
    def __iter__(self): return iter(self._out[::-1] + self._in)
\end{lstlisting}
\subsection{\texttt{dsu}}\label{def:dsu}
disjoint-set union (union-find) over n elements, with path halving. find gives a set's representative; union merges and returns False when the two were ALREADY connected -- which is exactly a cycle detector (Kruskal's test).
\begin{lstlisting}[style=ex]
>>> find, union = dsu(4)
>>> union(0, 1), union(1, 2)
(True, True)
>>> union(0, 2)                         # already connected: a cycle
False
>>> find(0) == find(2), find(0) == find(3)
(True, False)
\end{lstlisting}

How it works: parent[] links every element toward its set's ROOT, and find follows the chain up. The subtle line is parent[x] = parent[parent[x]] -- path HALVING: every find re-points each visited node at its grandparent, so chains shrink as a side effect of being walked. That one line makes operations effectively constant time (inverse Ackermann, for the curious).


Why it earns its place: dynamic connectivity. "Are these two in the same group yet?" under a stream of merges is awkward for every other structure and trivial here. And union returning False -- already connected -- IS cycle detection: Kruskal's spanning tree is built by skipping every False.

\begin{lstlisting}[style=ex]
>>> find, union = dsu(5)
>>> [union(a, b) for a, b in [(0, 1), (1, 2), (3, 4), (0, 2)]]
[True, True, True, False]
>>> len({find(i) for i in range(5)})   # two components remain
2
\end{lstlisting}

Remember it as: find walks to the root, halving as it goes; a False union is a cycle caught.


Roll your own: parent[] as the forest and two closures instead of a class. find's loop carries the entire trick in one line -- parent[x] = parent[parent[x]] re-points every visited node at its grandparent (path halving), flattening chains as a side effect of walking them. union is find twice plus one root re-point, answering False when the roots already agree.

\begin{lstlisting}[style=ex]
def dsu(n):  # union-find, path halving; union -> False if already joined
    parent = list(range(n))
    def find(x):
        while parent[x] != x:
            parent[x] = parent[parent[x]]
            x = parent[x]
        return x
    def union(a, b):
        ra, rb = find(a), find(b)
        if ra == rb:
            return False
        parent[ra] = rb
        return True
    return find, union
\end{lstlisting}
\subsection{\texttt{heappush}}\label{def:heappush}
push onto a min-heap kept in a plain list (sift-up). The smallest element is always h[0]. For a MAX-heap, push negated values.
\begin{lstlisting}[style=ex]
>>> h = []
>>> for x in [5, 1, 8]: heappush(h, x)
>>> h[0]
1
\end{lstlisting}

How it works: the flat list IS a binary tree -- h[i]'s children live at 2i+1 and 2i+2 -- kept under one invariant: every parent <= its children. heappush appends at the first free leaf and sifts UP, swapping with its parent while smaller; heappop moves the last leaf to the root and sifts DOWN toward the smaller child. Each walks one root-to-leaf path: O(log n).


Why the invariant is enough: nothing is fully sorted -- siblings sit in any order -- yet h[0] is ALWAYS the minimum, because the root beats its children, who beat theirs. You pay log n per operation for a permanently readable minimum; draining the heap is heapsort; pushing negatives turns min into max.

\begin{lstlisting}[style=ex]
>>> h = []
>>> for x in [7, 2, 9, 4]: heappush(h, x)
>>> h[0]
2
>>> [heappop(h) for _ in range(4)]
[2, 4, 7, 9]
\end{lstlisting}

Why it earns its place: "repeatedly take the best next thing" -- Dijkstra's frontier, running\_median's two halves, any streaming top-k. When next-best is needed many times while the data keeps changing, the heap beats re-sorting every time.


Remember it as: a flat list, a parent-beats-children promise, and log-n repairs.


Roll your own: append at the first free leaf, then sift UP -- while the parent at (i - 1) // 2 is larger, swap and climb. The index arithmetic IS the tree; watch the array after four pushes:

\begin{lstlisting}[style=ex]
>>> h = []
>>> for x in [7, 2, 9, 4]: heappush(h, x)
>>> h                       # root 2; children 4 and 9; 7 pushed down under 4
[2, 4, 9, 7]
\end{lstlisting}
\begin{lstlisting}[style=ex]
def heappush(h, x):  # min-heap on a plain list: sift-up
    h.append(x)
    i = len(h) - 1
    while i and h[(i - 1) // 2] > h[i]:
        h[(i - 1) // 2], h[i] = h[i], h[(i - 1) // 2]
        i = (i - 1) // 2
\end{lstlisting}
\noindent{\small\itshape Used before it is rolled: \texttt{heappop} in \S\ref{def:heappop} (p.~\pageref{def:heappop}), \texttt{on} in \S\ref{def:on} (p.~\pageref{def:on}).}

\subsection{\texttt{heappop}}\label{def:heappop}
remove and return the minimum (sift-down). Popping repeatedly drains in sorted order -- that is heapsort.
\begin{lstlisting}[style=ex]
>>> h = []
>>> for x in [5, 1, 8]: heappush(h, x)
>>> [heappop(h) for _ in range(3)]
[1, 5, 8]
>>> g = []
>>> for x in [3, 9, 5]: heappush(g, -x)
>>> -heappop(g)                         # max-heap via negation
9
\end{lstlisting}

Roll your own: swap the root with the last leaf, pop it off the end, then sift DOWN -- repeatedly swap with the SMALLER of the two children (choosing the smaller is the classic bug site) until neither child is smaller. One root-to-leaf path each way: O(log n).

\begin{lstlisting}[style=ex]
def heappop(h):  # min-heap on a plain list: sift-down
    h[0], h[-1] = h[-1], h[0]
    top = h.pop()
    i, n = 0, len(h)
    while True:
        s, l, r = i, 2*i + 1, 2*i + 2
        if l < n and h[l] < h[s]: s = l
        if r < n and h[r] < h[s]: s = r
        if s == i:
            return top
        h[i], h[s] = h[s], h[i]
        i = s
\end{lstlisting}
\noindent{\small\itshape Used before it is rolled: \texttt{on} in \S\ref{def:on} (p.~\pageref{def:on}).}

\section{Homework}
Work in order: drills cement each tool, applies combine them,
challenges are interview-grade. Write each solution in a scratch file with
the given doctests pasted in, and let \texttt{python3 -m doctest} referee.
\subsection*{12.1\quad Rotate left with a deque\hfill\textnormal{\textit{drill}}}
Write rotate\_left(xs, k): rotate a list k positions left using a deque -- popleft feeds append, k times. k may be 0, huge, or applied to an empty list.

\noindent\texttt{rotate\_left(xs: list, k: int) -> list}
\begin{lstlisting}[style=ex]
>>> rotate_left([1, 2, 3, 4], 1)
[2, 3, 4, 1]
>>> rotate_left([1, 2, 3, 4], 6)
[3, 4, 1, 2]
>>> rotate_left([], 3)
[]
>>> rotate_left([7], 99)
[7]
\end{lstlisting}
\subsection*{12.2\quad Keep only the last k\hfill\textnormal{\textit{drill}}}
A dashboard shows only the k most recent readings. Write last\_k(xs, k): feed the stream through a deque, evicting from the front whenever the buffer exceeds k, and return the final buffer as a list.

\noindent\texttt{last\_k(xs: list, k: int) -> list}
\begin{lstlisting}[style=ex]
>>> last_k([1, 2, 3, 4, 5], 2)
[4, 5]
>>> last_k([1, 2], 5)
[1, 2]
>>> last_k([], 3)
[]
\end{lstlisting}
\subsection*{12.3\quad Same group yet?\hfill\textnormal{\textit{drill}}}
Write same\_group(n, pairs, queries) for elements 0..n-1: apply every merge in pairs, then answer each query (a, b) with whether a and b ended up connected. Compare find results -- never raw parent values, which are arbitrary.

\noindent\texttt{same\_group(n: int, pairs: list[tuple], queries: list[tuple]) -> list[bool]}
\begin{lstlisting}[style=ex]
>>> same_group(5, [(0, 1), (1, 2)], [(0, 2), (0, 3), (3, 3)])
[True, False, True]
>>> same_group(3, [], [(0, 1)])
[False]
>>> same_group(2, [(0, 1)], [(1, 0)])
[True]
\end{lstlisting}
\subsection*{12.4\quad Heapsort by draining\hfill\textnormal{\textit{drill}}}
Write heap\_sorted(xs): push every element onto a min-heap, then pop until empty. The drain order IS sorted order -- that is heapsort.

\noindent\texttt{heap\_sorted(xs: list) -> list}
\begin{lstlisting}[style=ex]
>>> heap_sorted([5, 1, 4, 1, 3])
[1, 1, 3, 4, 5]
>>> heap_sorted([])
[]
>>> heap_sorted([2])
[2]
\end{lstlisting}
\subsection*{12.5\quad The k smallest\hfill\textnormal{\textit{drill}}}
Write k\_smallest(xs, k): the k smallest values in ascending order (all of them if k exceeds the length). Push everything, pop k times.

\noindent\texttt{k\_smallest(xs: list, k: int) -> list}
\begin{lstlisting}[style=ex]
>>> k_smallest([9, 4, 7, 1, 0], 3)
[0, 1, 4]
>>> k_smallest([5, 5, 5], 2)
[5, 5]
>>> k_smallest([1, 2], 10)
[1, 2]
\end{lstlisting}
\subsection*{12.6\quad Hot potato\hfill\textnormal{\textit{apply}}}
Players stand in a circle and count around: every m-th player is eliminated until one remains (the Josephus game). Write survivor(names, m): rotate the deque m-1 times (popleft feeding append), eliminate with a popleft, repeat. Return the last name standing.

\noindent\texttt{survivor(names: list[str], m: int) -> str}
\begin{lstlisting}[style=ex]
>>> survivor(['a', 'b', 'c', 'd', 'e'], 2)
'c'
>>> survivor(['solo'], 7)
'solo'
>>> survivor(['a', 'b', 'c'], 1)
'c'
>>> survivor(['a', 'b'], 3)
'b'
\end{lstlisting}
\subsection*{12.7\quad Merge k sorted feeds\hfill\textnormal{\textit{apply}}}
k sensors each emit an already-sorted list of readings. Write merge\_all(lists): one sorted stream, heap-powered: seed the heap with each list's head as (value, list\_index, position), then repeatedly pop the smallest and push that list's next reading. The index in the tuple breaks value ties so comparison never reaches unorderable data.

\noindent\texttt{merge\_all(lists: list[list[int]]) -> list[int]}
\begin{lstlisting}[style=ex]
>>> merge_all([[1, 4, 6], [2, 3], [5]])
[1, 2, 3, 4, 5, 6]
>>> merge_all([[], [1], []])
[1]
>>> merge_all([])
[]
>>> merge_all([[1, 1], [1]])
[1, 1, 1]
\end{lstlisting}
\subsection*{12.8\quad How many friend circles?\hfill\textnormal{\textit{apply}}}
n people, a list of friendships. Friendship is transitive through the merges. Write components(n, pairs): how many separate circles remain -- union every pair, then count DISTINCT roots via find over everyone.

\noindent\texttt{components(n: int, pairs: list[tuple]) -> int}
\begin{lstlisting}[style=ex]
>>> components(5, [(0, 1), (3, 4)])
3
>>> components(4, [])
4
>>> components(3, [(0, 1), (1, 2)])
1
>>> components(0, [])
0
\end{lstlisting}
\subsection*{12.9\quad The edge that closes a loop\hfill\textnormal{\textit{apply}}}
Network links arrive one at a time. Write first\_cycle\_edge(n, edges): the FIRST link that connects two already-connected nodes (union answering False), or None if every link merges something new. This is exactly how Kruskal's algorithm rejects edges.

\noindent\texttt{first\_cycle\_edge(n: int, edges: list[tuple]) -> tuple | None}
\begin{lstlisting}[style=ex]
>>> first_cycle_edge(5, [(0, 1), (1, 2), (0, 2), (3, 4)])
(0, 2)
>>> first_cycle_edge(3, [(0, 1), (1, 2)]) is None
True
>>> first_cycle_edge(2, [(0, 0)])
(0, 0)
\end{lstlisting}
\subsection*{12.10\quad k-th largest, live\hfill\textnormal{\textit{apply}}}
A leaderboard shows the k-th largest score seen so far, updating per submission. Write kth\_largest\_stream(xs, k): after each element, report the k-th largest so far, or None while fewer than k have arrived. Keep a min-heap of the k best -- its root IS the answer.

\noindent\texttt{kth\_largest\_stream(xs: list[int], k: int) -> list}
\begin{lstlisting}[style=ex]
>>> kth_largest_stream([3, 1, 5, 12, 2, 11], 3)
[None, None, 1, 3, 3, 5]
>>> kth_largest_stream([4, 4], 1)
[4, 4]
>>> kth_largest_stream([], 2)
[]
\end{lstlisting}
\subsection*{12.11\quad Sliding window maximum\hfill\textnormal{\textit{challenge}}}
Write window\_max(xs, k): the maximum of every length-k window, in order (classic hard interview question -- O(n) required, so no per-window max()). Keep a deque of INDICES with values decreasing: each newcomer pops smaller values off the back; the front is the window's max, dropped when it slides out of range. The prelude deque has no peek: peek by pop-then-push-back on the side you need.

\noindent\texttt{window\_max(xs: list[int], k: int) -> list[int]}
\begin{lstlisting}[style=ex]
>>> window_max([1, 3, -1, -3, 5, 3, 6, 7], 3)
[3, 3, 5, 5, 6, 7]
>>> window_max([2, 1], 1)
[2, 1]
>>> window_max([4, 2, 3], 3)
[4]
>>> window_max([], 2)
[]
\end{lstlisting}
\subsection*{12.12\quad Cheapest network (Kruskal)\hfill\textnormal{\textit{challenge}}}
Weighted links (w, u, v) can connect n offices. Write kruskal(n, edges): the minimum total cost to connect ALL offices, or None if impossible. Sort edges by weight, take every edge whose union succeeds, stop counting when n-1 merges have landed. sortOn + dsu is the whole algorithm.

\noindent\texttt{kruskal(n: int, edges: list[tuple]) -> int | None}
\begin{lstlisting}[style=ex]
>>> kruskal(4, [(1, 0, 1), (2, 1, 2), (10, 0, 2), (3, 2, 3)])
6
>>> kruskal(3, [(5, 0, 1)]) is None
True
>>> kruskal(1, [])
0
>>> kruskal(4, [(4, 0, 1), (1, 2, 3), (2, 1, 2), (7, 0, 3)])
7
\end{lstlisting}
\clearpage
\chapter{Nodes: Linked Lists and Binary Trees}

The LeetCode givens. For trees, one catamorphism -- \texttt{tfold} -- makes
every traversal and measure a one-liner, because \texttt{z} is the
empty-tree arm and \texttt{f} the node arm; \texttt{levelorder} adds the
deque for breadth-first. For linked lists, the pointer idioms worth knowing
cold: three-pointer reversal, fast/slow middles, Floyd's cycle detection,
dummy-head merging. Convert node results with \texttt{to\_list} so
everything stays testable.
\section{Concepts}
\subsection{\texttt{TreeNode}}\label{def:TreeNode}
the LeetCode binary tree: val, left, right, with None for absent children.
\begin{lstlisting}[style=ex]
>>> t = TreeNode(2, TreeNode(1), TreeNode(3))
>>> t.left.val, t.val, t.right.val
(1, 2, 3)
\end{lstlisting}

Roll your own --- the definition:

\begin{lstlisting}[style=ex]
class TreeNode:  # the LeetCode binary tree
    def __init__(self, val=0, left=None, right=None):
        self.val, self.left, self.right = val, left, right
\end{lstlisting}
\subsection{\texttt{tfold}}\label{def:tfold}
the tree catamorphism: recursively combines each node's value with the already-folded left and right results, z standing in for empty subtrees. EVERY structural tree computation is one choice of f and z -- the five traversal/measure functions below are all tfold one-liners.
\begin{lstlisting}[style=ex]
>>> t = TreeNode(2, TreeNode(1), TreeNode(3))
>>> tfold(lambda v, l, r: v + l + r, 0, t)      # sum
6
>>> tfold(lambda v, l, r: max(v, l, r), 0, t)   # max
3
\end{lstlisting}

How it works: the tree catamorphism -- structural recursion with a contract. Two arms, exactly like a case expression: the empty tree answers z; a node answers f(value, left\_result, right\_result), both results computed by the same recursion. Hand it (f, z) and any question about a tree becomes one line -- the five traversals and both measures in section 13 are all tfold with different arms.


Why it earns its place: it separates WALKING from DECIDING. The recursion is written once, correct forever; you only ever author the arms. Even tree TRANSFORMATION fits -- mirror a tree by swapping the child results while rebuilding:

\begin{lstlisting}[style=ex]
>>> t = TreeNode(2, TreeNode(1), TreeNode(3))
>>> mirror = tfold(lambda v, l, r: TreeNode(v, r, l), None, t)
>>> inorder(mirror)
[3, 2, 1]
\end{lstlisting}

The one caveat: recursion depth equals tree height -- fine for balanced and interview-sized trees; say a sentence out loud about adversarially skewed ones.


Remember it as: z is the empty-tree arm, f is the node arm -- the fold walks, you decide.


Roll your own: one line of structural recursion -- z when the tree is None, else f over the node's value and the two recursive results.


Roll your own: one line of structural recursion -- z when the tree is None, else f over the node's value and the two recursive results. Every traversal and measure in this section is that line with a different f plugged in.

\begin{lstlisting}[style=ex]
tfold = lambda f, z, t: z if t is None else f(t.val, tfold(f, z, t.left), tfold(f, z, t.right))
\end{lstlisting}
\noindent{\small\itshape Used before it is rolled: \texttt{inorder} in \S\ref{def:inorder} (p.~\pageref{def:inorder}).}

\subsection{\texttt{inorder}}\label{def:inorder}
left, node, right; on a binary SEARCH tree this is sorted order.
\begin{lstlisting}[style=ex]
>>> inorder(TreeNode(2, TreeNode(1), TreeNode(3)))
[1, 2, 3]
\end{lstlisting}

Roll your own --- the definition:

\begin{lstlisting}[style=ex]
inorder = lambda t: tfold(lambda v, l, r: l + [v] + r, [], t)
\end{lstlisting}
\subsection{\texttt{preorder}}\label{def:preorder}
node first, then children: the copy/serialise order.
\begin{lstlisting}[style=ex]
>>> preorder(TreeNode(2, TreeNode(1), TreeNode(3)))
[2, 1, 3]
\end{lstlisting}

Roll your own --- the definition:

\begin{lstlisting}[style=ex]
preorder = lambda t: tfold(lambda v, l, r: [v] + l + r, [], t)
\end{lstlisting}
\subsection{\texttt{postorder}}\label{def:postorder}
children first, then node: the delete/evaluate order.
\begin{lstlisting}[style=ex]
>>> postorder(TreeNode(2, TreeNode(1), TreeNode(3)))
[1, 3, 2]
\end{lstlisting}

Roll your own --- the definition:

\begin{lstlisting}[style=ex]
postorder = lambda t: tfold(lambda v, l, r: l + r + [v], [], t)
\end{lstlisting}
\subsection{\texttt{tdepth}}\label{def:tdepth}
the height of the tree.
\begin{lstlisting}[style=ex]
>>> tdepth(TreeNode(2, TreeNode(1, TreeNode(0)), TreeNode(3)))
3
\end{lstlisting}

Roll your own --- the definition:

\begin{lstlisting}[style=ex]
tdepth = lambda t: tfold(lambda _, l, r: 1 + max(l, r), 0, t)
\end{lstlisting}
\subsection{\texttt{tsize}}\label{def:tsize}
the number of nodes.
\begin{lstlisting}[style=ex]
>>> tsize(TreeNode(2, TreeNode(1), TreeNode(3)))
3
\end{lstlisting}

Roll your own --- the definition:

\begin{lstlisting}[style=ex]
tsize = lambda t: tfold(lambda _, l, r: 1 + l + r, 0, t)
\end{lstlisting}
\subsection{\texttt{levelorder}}\label{def:levelorder}
breadth-first values, level by level: the deque earning its keep. Compare preorder on the same tree to see DFS vs BFS.
\begin{lstlisting}[style=ex]
>>> t = TreeNode(1, TreeNode(2, TreeNode(4)), TreeNode(3))
>>> levelorder(t)
[1, 2, 3, 4]
>>> preorder(t)
[1, 2, 4, 3]
\end{lstlisting}

Roll your own: seed a deque with the root, then loop -- popleft, record the value, append the children that exist. FIFO order IS breadth-first; swap the deque for a list used as a stack and the identical loop turns into depth-first.

\begin{lstlisting}[style=ex]
def levelorder(t):  # BFS: the deque earning its keep
    out, q = [], deque([t] if t else [])
    while len(q):
        n = q.popleft()
        out.append(n.val)
        if n.left:  q.append(n.left)
        if n.right: q.append(n.right)
    return out
\end{lstlisting}
\subsection{\texttt{ListNode}}\label{def:ListNode}
the LeetCode singly-linked list: val and next.
\begin{lstlisting}[style=ex]
>>> n = ListNode(1, ListNode(2))
>>> n.val, n.next.val
(1, 2)
\end{lstlisting}

Roll your own --- the definition:

\begin{lstlisting}[style=ex]
class ListNode:  # the LeetCode linked list
    def __init__(self, val=0, next=None):
        self.val, self.next = val, next
\end{lstlisting}
\subsection{\texttt{from\_list}}\label{def:from-list}
build a linked list from a Python list: literally foldr(ListNode, xs, None). foldr's f takes (element, accumulator), and ListNode(val, next) has exactly that shape, so the last element is wrapped first and each earlier one points at the list built so far.
\begin{lstlisting}[style=ex]
>>> to_list(from_list([1, 2, 3]))
[1, 2, 3]
\end{lstlisting}

Roll your own --- the definition:

\begin{lstlisting}[style=ex]
from_list = lambda xs: foldr(ListNode, xs, None)  # build; foldr of ListNode
\end{lstlisting}
\noindent{\small\itshape Used before it is rolled: \texttt{to\_list} in \S\ref{def:to-list} (p.~\pageref{def:to-list}).}

\subsection{\texttt{to\_list}}\label{def:to-list}
walk a linked list back into a Python list; the round-trip partner of from\_list, and the way to make node results printable. An unfoldr: the seed is the node, each step yields (node.val, node.next), and None ends it.
\begin{lstlisting}[style=ex]
>>> to_list(from_list("ab"))
['a', 'b']
>>> to_list(None)
[]
\end{lstlisting}

Roll your own --- the definition:

\begin{lstlisting}[style=ex]
to_list = lambda node: list(unfoldr(lambda n: None if n is None else (n.val, n.next), node))
\end{lstlisting}
\subsection{\texttt{reverse\_list}}\label{def:reverse-list}
reverse in place with the three-pointer walk; prev is a fold accumulator wearing pointers.
\begin{lstlisting}[style=ex]
>>> to_list(reverse_list(from_list([1, 2, 3])))
[3, 2, 1]
\end{lstlisting}

Roll your own: the three-pointer walk compressed into one simultaneous assignment -- node.next, prev, node = prev, node, node.next. The right-hand side is evaluated in full before any assignment lands, which is what saves the pointer you are about to overwrite.

\begin{lstlisting}[style=ex]
def reverse_list(node):  # three-pointer; prev is a fold accumulator
    prev = None
    while node:
        node.next, prev, node = prev, node, node.next
    return prev
\end{lstlisting}
\subsection{\texttt{middle}}\label{def:middle}
the middle node by fast/slow pointers (fast moves twice per step); on even lengths, the SECOND middle.
\begin{lstlisting}[style=ex]
>>> middle(from_list([1, 2, 3])).val
2
>>> middle(from_list([1, 2, 3, 4])).val
3
\end{lstlisting}

Roll your own: fast advances two hops for slow's one, so when fast falls off the end, slow stands at the middle. The loop condition (fast and fast.next) is precisely what selects the SECOND middle on even lengths.

\begin{lstlisting}[style=ex]
def middle(node):  # fast/slow: second middle for even lengths
    slow = fast = node
    while fast and fast.next:
        slow, fast = slow.next, fast.next.next
    return slow
\end{lstlisting}
\subsection{\texttt{has\_cycle}}\label{def:has-cycle}
Floyd's tortoise and hare: the fast pointer laps the slow one if and only if a cycle exists; O(1) space.
\begin{lstlisting}[style=ex]
>>> has_cycle(from_list([1, 2, 3]))
False
>>> n = ListNode(1); n.next = n         # a self-loop
>>> has_cycle(n)
True
\end{lstlisting}

Roll your own: the same fast/slow walk, but watching for the pointers to MEET -- on a cycle the fast runner laps the slow one; on a straight list it simply falls off. O(1) space, no visited set.

\begin{lstlisting}[style=ex]
def has_cycle(node):  # Floyd: fast laps slow iff a cycle exists
    slow = fast = node
    while fast and fast.next:
        slow, fast = slow.next, fast.next.next
        if slow is fast:
            return True
    return False
\end{lstlisting}
\subsection{\texttt{merge\_lists}}\label{def:merge-lists}
merge two SORTED linked lists, using the dummy-head idiom to avoid special-casing the first node.
\begin{lstlisting}[style=ex]
>>> to_list(merge_lists(from_list([1, 3]), from_list([2, 4])))
[1, 2, 3, 4]
\end{lstlisting}

Roll your own: the dummy-head idiom -- a throwaway node to hang the result from, so appending the first real node needs no special case; walk both lists copying the smaller head, glue the survivor, return dummy.next.

\begin{lstlisting}[style=ex]
def merge_lists(a, b):  # merge two sorted lists; dummy-head idiom
    dummy = tail_ = ListNode()
    while a and b:
        if a.val <= b.val:
            tail_.next, a = a, a.next
        else:
            tail_.next, b = b, b.next
        tail_ = tail_.next
    tail_.next = a or b
    return dummy.next
\end{lstlisting}
\section{Homework}
Work in order: drills cement each tool, applies combine them,
challenges are interview-grade. Write each solution in a scratch file with
the given doctests pasted in, and let \texttt{python3 -m doctest} referee.
\subsection*{13.1\quad Sum of the tree\hfill\textnormal{\textit{drill}}}
Return the sum of every value in a binary tree. The empty tree sums to 0. Use tfold: pick the combining arm and the empty-tree seed.

\noindent\texttt{tree\_sum(t: TreeNode | None) -> int}
\begin{lstlisting}[style=ex]
>>> tree_sum(TreeNode(2, TreeNode(1), TreeNode(3)))
6
>>> tree_sum(TreeNode(5, TreeNode(3, TreeNode(1)), TreeNode(2)))
11
>>> tree_sum(None)
0
\end{lstlisting}
\subsection*{13.2\quad Largest value in the tree\hfill\textnormal{\textit{drill}}}
Return the largest value in a NON-EMPTY binary tree of integers, again as one tfold. Think about what seed is safe when every real value must be able to beat it.

\noindent\texttt{tree\_max(t: TreeNode) -> int   \# t is non-empty}
\begin{lstlisting}[style=ex]
>>> tree_max(TreeNode(2, TreeNode(9), TreeNode(3)))
9
>>> tree_max(TreeNode(-5, TreeNode(-2, TreeNode(-8)), None))
-2
>>> tree_max(TreeNode(7))
7
\end{lstlisting}
\subsection*{13.3\quad Count the leaves\hfill\textnormal{\textit{drill}}}
Count the LEAVES (nodes with no children) of a binary tree. One tfold suffices -- notice that a node is a leaf exactly when both child results are zero.

\noindent\texttt{count\_leaves(t: TreeNode | None) -> int}
\begin{lstlisting}[style=ex]
>>> count_leaves(TreeNode(2, TreeNode(1), TreeNode(3)))
2
>>> count_leaves(TreeNode(1, TreeNode(2, TreeNode(3))))
1
>>> count_leaves(TreeNode(7))
1
>>> count_leaves(None)
0
\end{lstlisting}
\subsection*{13.4\quad Linked length\hfill\textnormal{\textit{drill}}}
Return how many nodes a linked list holds. Convert once and measure -- to\_list is the bridge between pointer-land and list-land.

\noindent\texttt{llen(node: ListNode | None) -> int}
\begin{lstlisting}[style=ex]
>>> llen(from_list([1, 2, 3]))
3
>>> llen(from_list([]))
0
>>> llen(ListNode(9))
1
\end{lstlisting}
\subsection*{13.5\quad The traversal triple\hfill\textnormal{\textit{drill}}}
Return a tree's (inorder, preorder, postorder) value lists as one tuple. All three are tfold one-liners already in the prelude; this drill is about KNOWING their orders cold.

\noindent\texttt{traversals(t: TreeNode | None) -> tuple[list, list, list]}
\begin{lstlisting}[style=ex]
>>> traversals(TreeNode(2, TreeNode(1), TreeNode(3)))
([1, 2, 3], [2, 1, 3], [1, 3, 2])
>>> traversals(TreeNode(3, TreeNode(2, TreeNode(1))))
([1, 2, 3], [3, 2, 1], [1, 2, 3])
>>> traversals(None)
([], [], [])
\end{lstlisting}
\subsection*{13.6\quad Kth newest audit entry\hfill\textnormal{\textit{apply}}}
A server keeps an append-only audit trail as a linked list, oldest entry first. Ops asks for the k-th NEWEST entry (k = 1 is the last one). 1 <= k <= length is guaranteed.

\noindent\texttt{kth\_from\_end(node: ListNode, k: int) -> value}
\begin{lstlisting}[style=ex]
>>> kth_from_end(from_list([10, 20, 30, 40]), 1)
40
>>> kth_from_end(from_list([10, 20, 30, 40]), 4)
10
>>> kth_from_end(from_list([7]), 1)
7
\end{lstlisting}
\subsection*{13.7\quad Palindrome ticket tape\hfill\textnormal{\textit{apply}}}
A ticket printer spools IDs onto a tape held as a linked list. QA wants to know whether the tape reads the same in both directions. The empty tape counts as a palindrome.

\noindent\texttt{is\_palindrome\_list(node: ListNode | None) -> bool}
\begin{lstlisting}[style=ex]
>>> is_palindrome_list(from_list([1, 2, 1]))
True
>>> is_palindrome_list(from_list([1, 2]))
False
>>> is_palindrome_list(from_list([]))
True
>>> is_palindrome_list(from_list([7]))
True
\end{lstlisting}
\subsection*{13.8\quad Deduplicate a sorted list\hfill\textnormal{\textit{apply}}}
Subscription events arrive as a SORTED linked list of customer ids, possibly with repeats. Return a linked list with each id once, order preserved.

\noindent\texttt{dedup(node: ListNode | None) -> ListNode | None}
\begin{lstlisting}[style=ex]
>>> to_list(dedup(from_list([1, 1, 2, 3, 3])))
[1, 2, 3]
>>> to_list(dedup(from_list([4, 4, 4])))
[4]
>>> to_list(dedup(from_list([])))
[]
>>> to_list(dedup(from_list([1, 2])))
[1, 2]
\end{lstlisting}
\subsection*{13.9\quad Is it a search tree?\hfill\textnormal{\textit{apply}}}
A binary tree claims to be a binary SEARCH tree with strictly increasing values. Verify the claim. Key insight: a BST is exactly a tree whose inorder walk is strictly increasing -- so check the walk, not the tree.

\noindent\texttt{is\_bst(t: TreeNode | None) -> bool}
\begin{lstlisting}[style=ex]
>>> is_bst(TreeNode(2, TreeNode(1), TreeNode(3)))
True
>>> is_bst(TreeNode(2, TreeNode(3), TreeNode(1)))
False
>>> is_bst(TreeNode(2, TreeNode(2)))
False
>>> is_bst(None)
True
\end{lstlisting}
\subsection*{13.10\quad Symmetric tree\hfill\textnormal{\textit{apply}}}
An org chart is 'symmetric' if the whole tree equals its own mirror image -- values AND structure. Check it by folding the tree into a nested-tuple shape twice: once as-is, once with the child results swapped, and comparing.

\noindent\texttt{is\_symmetric(t: TreeNode | None) -> bool}
\begin{lstlisting}[style=ex]
>>> is_symmetric(TreeNode(1, TreeNode(2), TreeNode(2)))
True
>>> is_symmetric(TreeNode(1, TreeNode(2), TreeNode(3)))
False
>>> is_symmetric(TreeNode(1, TreeNode(2, TreeNode(3)), TreeNode(2, TreeNode(3))))
False
>>> is_symmetric(TreeNode(1, TreeNode(2, None, TreeNode(3)), TreeNode(2, TreeNode(3))))
True
>>> is_symmetric(None)
True
\end{lstlisting}
\subsection*{13.11\quad Merge k sorted feeds\hfill\textnormal{\textit{challenge}}}
k sensor feeds each deliver readings as a sorted linked list. Produce one sorted linked list containing everything. An empty collection of feeds yields the empty list. Hint: you already own a two-list merge; what fold turns two-way into k-way?

\noindent\texttt{merge\_k(feeds: list[ListNode | None]) -> ListNode | None}
\begin{lstlisting}[style=ex]
>>> to_list(merge_k([from_list([1, 4]), from_list([2, 3]), from_list([5])]))
[1, 2, 3, 4, 5]
>>> to_list(merge_k([from_list([]), from_list([2]), from_list([1, 3])]))
[1, 2, 3]
>>> to_list(merge_k([]))
[]
>>> to_list(merge_k([from_list([9])]))
[9]
\end{lstlisting}
\subsection*{13.12\quad Height-balanced?\hfill\textnormal{\textit{challenge}}}
A tree is height-balanced when EVERY node's two subtrees differ in height by at most 1. Decide it in one pass. Hint: tfold can carry a PAIR upward -- (height, balanced\_so\_far) -- so each node judges itself from its children's pairs.

\noindent\texttt{is\_balanced(t: TreeNode | None) -> bool}
\begin{lstlisting}[style=ex]
>>> is_balanced(TreeNode(1, TreeNode(2), TreeNode(3)))
True
>>> is_balanced(TreeNode(1, TreeNode(2, TreeNode(3))))
False
>>> is_balanced(None)
True
>>> deep = TreeNode(2, TreeNode(3, TreeNode(4)))
>>> is_balanced(TreeNode(1, deep, TreeNode(9)))
False
\end{lstlisting}
\clearpage
\chapter{Grids and Windows}

Two small tools with outsized reach. The neighbour generators put the
bounds check in ONE place so every flood-fill and BFS body stays clean.
\texttt{longest\_window} inverts control for sliding windows: the skeleton
owns the two pointers, you own the state as three closures --
\texttt{valid}, \texttt{add}, \texttt{rem} -- and designing them is just
answering ``what makes a window illegal?''.
\section{Concepts}
\subsection{\texttt{neighbors4}}\label{def:neighbors4}
the in-bounds orthogonal neighbours of cell (r, c) in an R-by-C grid. Bounds checking lives HERE, so BFS bodies stay clean.
\begin{lstlisting}[style=ex]
>>> sorted(neighbors4(0, 0, 2, 2))      # corner: two neighbours
[(0, 1), (1, 0)]
>>> sorted(neighbors4(1, 1, 3, 3))      # centre: four
[(0, 1), (1, 0), (1, 2), (2, 1)]
\end{lstlisting}

Roll your own: yield each in-bounds delta of (r, c) -- the bounds check lives HERE so every caller's BFS body stays clean. neighbors8 is the same generator over a 3x3 delta grid, with (dr or dc) excluding the centre cell.

\begin{lstlisting}[style=ex]
def neighbors4(r, c, R, C):  # in-bounds orthogonal neighbors
    for dr, dc in ((1, 0), (-1, 0), (0, 1), (0, -1)):
        if 0 <= r + dr < R and 0 <= c + dc < C:
            yield r + dr, c + dc
\end{lstlisting}
\subsection{\texttt{neighbors8}}\label{def:neighbors8}
the same plus diagonals.
\begin{lstlisting}[style=ex]
>>> len(list(neighbors8(1, 1, 3, 3)))
8
>>> sorted(neighbors8(0, 0, 2, 2))
[(0, 1), (1, 0), (1, 1)]
\end{lstlisting}

Roll your own --- the definition:

\begin{lstlisting}[style=ex]
def neighbors8(r, c, R, C):  # ... plus diagonals
    for dr in (-1, 0, 1):
        for dc in (-1, 0, 1):
            if (dr or dc) and 0 <= r + dr < R and 0 <= c + dc < C:
                yield r + dr, c + dc
\end{lstlisting}
\subsection{\texttt{longest\_window}}\label{def:longest-window}
the sliding-window skeleton: it grows the right edge one element at a time (calling add), shrinks the left edge while your valid() says the window is illegal (calling rem), and tracks the best length. YOU supply the state as closures; the skeleton does the two-pointer bookkeeping. Longest- substring-without-repeats and longest-run-under-a- budget are both three closures away.
\begin{lstlisting}[style=ex]
>>> seen = []
>>> longest_window("abcabb", lambda: len(seen) == len(set(seen)),
...                seen.append, lambda c: seen.pop(0))
3
>>> w = []
>>> longest_window([2, 1, 3, 4], lambda: sum(w) <= 6,
...                w.append, lambda x: w.pop(0))
3
\end{lstlisting}

How it works: inversion of control. The skeleton owns the two pointers -- it grows hi one element per step (calling add), then shrinks lo while your valid() reports the window broken (calling rem per evicted element) -- and it tracks the best length ever seen legal. You own only the STATE, held in closures over variables sitting next to the call.


Designing the closures is answering one question: "what makes a window ILLEGAL, and what must I track to know?" Distinct characters -> a duplicate count. A budget -> a running sum. At most k of something -> a counter. add and rem are that state's increment and decrement; valid is the legality test.

\begin{lstlisting}[style=ex]
>>> w = []
>>> longest_window("aabcb", lambda: len(w) == len(set(w)), w.append, lambda c: w.pop(0))
3
\end{lstlisting}

Why closures and not a class: the state lives exactly as long as the call, sits in scope where you can read it, and needs no ceremony. Functions carrying state -- the same muscle as curry's captured argument, here working a two-pointer job.


Remember it as: the skeleton slides the window; your three closures say what legal means.


Roll your own: hi sweeps forward exactly once, add() greeting each newcomer; an inner while shrinks lo, rem() by rem(), until valid() approves again. Every element is added once and removed at most once -- the entire two-pointer O(n) argument, visible in two loops.

\begin{lstlisting}[style=ex]
def longest_window(xs, valid, add, rem):  # sliding-window skeleton; state lives in the closures
    lo = best = 0
    for hi in range(len(xs)):
        add(xs[hi])
        while not valid():  # shrink until legal again
            rem(xs[lo])
            lo += 1
        best = max(best, hi - lo + 1)
    return best
\end{lstlisting}
\noindent{\small\itshape Used before it is rolled: \texttt{until} in \S\ref{def:until} (p.~\pageref{def:until}).}

\section{Homework}
Work in order: drills cement each tool, applies combine them,
challenges are interview-grade. Write each solution in a scratch file with
the given doctests pasted in, and let \texttt{python3 -m doctest} referee.
\subsection*{14.1\quad How many neighbours?\hfill\textnormal{\textit{drill}}}
Return how many in-bounds orthogonal neighbours cell (r, c) has in an R-by-C grid. neighbors4 already does the bounds checking -- you just count what it yields.

\noindent\texttt{degree(r: int, c: int, R: int, C: int) -> int}
\begin{lstlisting}[style=ex]
>>> degree(0, 0, 2, 2)
2
>>> degree(0, 1, 3, 3)
3
>>> degree(1, 1, 3, 3)
4
>>> degree(0, 0, 1, 1)
0
\end{lstlisting}
\subsection*{14.2\quad Diagonal neighbours only\hfill\textnormal{\textit{drill}}}
Return the DIAGONAL in-bounds neighbours of (r, c), sorted. You own two neighbour generators; the diagonals are exactly what one yields and the other does not.

\noindent\texttt{diag\_neighbors(r: int, c: int, R: int, C: int) -> list[tuple[int, int]]}
\begin{lstlisting}[style=ex]
>>> diag_neighbors(0, 0, 3, 3)
[(1, 1)]
>>> diag_neighbors(1, 1, 3, 3)
[(0, 0), (0, 2), (2, 0), (2, 2)]
>>> diag_neighbors(0, 1, 1, 3)
[]
\end{lstlisting}
\subsection*{14.3\quad Live neighbours\hfill\textnormal{\textit{drill}}}
A Game-of-Life board is a list of '0'/'1' strings. Count the live cells among the 8-neighbourhood of (r, c).

\noindent\texttt{live\_around(grid: list[str], r: int, c: int) -> int}
\begin{lstlisting}[style=ex]
>>> board = ['111', '101', '111']
>>> live_around(board, 1, 1)
8
>>> live_around(board, 0, 0)
2
>>> live_around(['1'], 0, 0)
0
\end{lstlisting}
\subsection*{14.4\quad Longest positive stretch\hfill\textnormal{\textit{drill}}}
Given daily profit/loss numbers, find the length of the longest run of strictly positive days. Use longest\_window: the state you must track is how many non-positive values sit inside the current window.

\noindent\texttt{longest\_positive(xs: list[int]) -> int}
\begin{lstlisting}[style=ex]
>>> longest_positive([1, 2, -1, 3, 4, 5])
3
>>> longest_positive([-1, -2])
0
>>> longest_positive([])
0
>>> longest_positive([5, 6])
2
\end{lstlisting}
\subsection*{14.5\quad Cells on the border\hfill\textnormal{\textit{drill}}}
How many cells of an R-by-C grid lie on its border? Characterise the border through neighbours: a border cell is one with fewer than 4 orthogonal neighbours.

\noindent\texttt{border\_count(R: int, C: int) -> int}
\begin{lstlisting}[style=ex]
>>> border_count(1, 1)
1
>>> border_count(2, 3)
6
>>> border_count(3, 3)
8
>>> border_count(1, 5)
5
\end{lstlisting}
\subsection*{14.6\quad Count the islands\hfill\textnormal{\textit{apply}}}
A satellite tile is a list of '1'/'0' strings, land and water. Count the connected regions of land (orthogonal adjacency). Scan every cell; each unseen land cell starts a new island and a BFS flood that claims the rest of it.

\noindent\texttt{count\_islands(grid: list[str]) -> int}
\begin{lstlisting}[style=ex]
>>> count_islands(['11000', '11000', '00100'])
2
>>> count_islands(['101', '010', '101'])
5
>>> count_islands(['000'])
0
>>> count_islands(['1'])
1
\end{lstlisting}
\subsection*{14.7\quad At most k distinct products\hfill\textnormal{\textit{apply}}}
A shelf camera streams the product code seen each second, as a string of letters. Find the longest stretch showing at most k DISTINCT products. Window state: per-code counts plus how many codes are currently present.

\noindent\texttt{longest\_at\_most\_k(s: str, k: int) -> int}
\begin{lstlisting}[style=ex]
>>> longest_at_most_k('eceba', 2)
3
>>> longest_at_most_k('aa', 1)
2
>>> longest_at_most_k('aabbcc', 1)
2
>>> longest_at_most_k('abc', 0)
0
\end{lstlisting}
\subsection*{14.8\quad Longest run within budget\hfill\textnormal{\textit{apply}}}
Task durations stream in (non-negative). What is the longest run of CONSECUTIVE tasks whose total fits in a time budget? Window state: the running sum.

\noindent\texttt{longest\_within\_budget(xs: list[int], budget: int) -> int}
\begin{lstlisting}[style=ex]
>>> longest_within_budget([2, 1, 3, 4], 6)
3
>>> longest_within_budget([5], 3)
0
>>> longest_within_budget([], 10)
0
>>> longest_within_budget([1, 1, 1, 1], 2)
2
\end{lstlisting}
\subsection*{14.9\quad Island perimeter\hfill\textnormal{\textit{apply}}}
One island of '1' cells sits in a '0' sea. Its perimeter is the number of unit edges touching water or the map border. Count per land cell: 4 minus its LAND neighbours (off-grid sides count as water automatically, since neighbors4 never yields them).

\noindent\texttt{island\_perimeter(grid: list[str]) -> int}
\begin{lstlisting}[style=ex]
>>> island_perimeter(['0100', '1110', '0100'])
12
>>> island_perimeter(['1'])
4
>>> island_perimeter(['11'])
6
>>> island_perimeter(['00'])
0
\end{lstlisting}
\subsection*{14.10\quad Flood clock\hfill\textnormal{\textit{apply}}}
Water enters an open grid ('.' floor, '\#' wall) at the top-left and spreads one orthogonal step per minute. Return the minute the LAST reachable floor cell gets wet (0 if only the entrance), or -1 if the entrance itself is a wall.

\noindent\texttt{flood\_minutes(grid: list[str]) -> int}
\begin{lstlisting}[style=ex]
>>> flood_minutes(['..', '..'])
2
>>> flood_minutes(['.'])
0
>>> flood_minutes(['#.'])
-1
>>> flood_minutes(['.#.', '...', '.#.'])
4
\end{lstlisting}
\subsection*{14.11\quad Best uptime with k repairs\hfill\textnormal{\textit{challenge}}}
A service log is a list of 1s (up) and 0s (down). You may retroactively excuse at most k downtimes. What is the longest stretch that then reads fully up? Classic interview form: longest window containing at most k zeros.

\noindent\texttt{max\_up\_with\_repairs(xs: list[int], k: int) -> int}
\begin{lstlisting}[style=ex]
>>> max_up_with_repairs([1, 1, 0, 0, 1], 1)
3
>>> max_up_with_repairs([1, 1, 1, 0, 0, 0, 1, 1, 1, 1, 0], 2)
6
>>> max_up_with_repairs([1, 1, 0, 1], 0)
2
>>> max_up_with_repairs([], 3)
0
\end{lstlisting}
\subsection*{14.12\quad Rotting oranges\hfill\textnormal{\textit{challenge}}}
A crate map holds 0 (empty), 1 (fresh), 2 (rotten). Every minute, fresh oranges next to a rotten one (orthogonally) rot. Return the minutes until nothing fresh remains, or -1 if some fresh orange can never be reached. Multi-source BFS: seed the queue with EVERY rotten orange at minute 0.

\noindent\texttt{orange\_minutes(grid: list[list[int]]) -> int}
\begin{lstlisting}[style=ex]
>>> orange_minutes([[2, 1, 1], [1, 1, 0], [0, 1, 1]])
4
>>> orange_minutes([[2, 1, 1], [0, 1, 1], [1, 0, 1]])
-1
>>> orange_minutes([[0, 2]])
0
>>> orange_minutes([[1]])
-1
\end{lstlisting}
\clearpage
\chapter{Control and Dynamic Programming}

The capstone machinery. \texttt{memo} turns an honest recurrence into
dynamic programming -- state the truth, decorate it, and the cache is your
DP table, inspectable at \texttt{.cache}; the craft is naming the state so
``the future only needs $X$'' comes out in one sentence. \texttt{until}
handles the loops with no list to fold: iterate a step to a goal or a fixed
point, Bellman--Ford being the textbook case.
\section{Concepts}
\subsection{\texttt{memo}}\label{def:memo}
an unbounded memoizer as a plain decorator (no parentheses): results are cached by the argument tuple, so each distinct call computes once. Keyword arguments fold into the key; the cache is exposed as .cache.
\begin{lstlisting}[style=ex]
>>> @memo
... def fib(n): return n if n < 2 else fib(n - 1) + fib(n - 2)
>>> fib(30)
832040
>>> fib.cache[(10,)]
55
>>> len(fib.cache)                       # one entry per distinct argument
31
\end{lstlisting}

How it works: the decorator wraps f with a dict keyed on the arguments. First call computes and stores; every repeat is a lookup. That is the entire mechanism -- and why arguments must be HASHABLE: lists become tuples before they can be keys.


Why it earns its place: memoized recursion IS dynamic programming. Write the honest recurrence -- the answer for state (i, j) in terms of smaller states -- decorate it, and the cache is your DP table, filled in dependency order with no index bookkeeping. The exposed .cache lets you SEE the table: its size is the state count, exactly the number complexity analysis wants said out loud.

\begin{lstlisting}[style=ex]
>>> @memo
... def ways(r, c):                       # lattice paths to (r, c), moving right/down
...     return 1 if r == 0 or c == 0 else ways(r - 1, c) + ways(r, c - 1)
>>> ways(3, 3)
20
>>> len(ways.cache)                       # the 4x4 grid minus (0,0), never needed
15
\end{lstlisting}

Versus functools.lru\_cache: the standard library's version adds eviction (maxsize); memo is deliberately unbounded, because DP wants every state kept -- and it is nine lines you can retype in any interview that bans imports.


Remember it as: recursion states the truth; the cache makes it affordable; .cache is the table.


Roll your own: a dict captured in a closure, keyed by the argument tuple (keyword arguments folded in as a frozenset when present -- tuples and frozensets because keys must be hashable). Compute on miss, look up on hit, and hang the dict on the wrapper as .cache so the DP table stays inspectable.

\begin{lstlisting}[style=ex]
def memo(f):  # unbounded memoizer; enough for DP (no eviction by design)
    cache = {}
    def wrapped(*args, **kw):
        key = (args, frozenset(kw.items())) if kw else args
        if key not in cache:
            cache[key] = f(*args, **kw)
        return cache[key]
    wrapped.cache = cache  # peek at the DP table if curious
    return wrapped
\end{lstlisting}
\subsection{\texttt{until}}\label{def:until}
iterate f from x until the predicate holds, returning the first satisfying value: a fixed-point loop with no recursion-depth limit. The state can be a tuple carrying whatever the loop must remember.
\begin{lstlisting}[style=ex]
>>> until(lambda x: x > 100, lambda x: x * 2, 1)
128
>>> collatz = lambda s: (s[0] // 2 if even(s[0]) else 3 * s[0] + 1, s[1] + 1)
>>> until(lambda s: s[0] == 1, collatz, (6, 0))     # 8 steps to reach 1
(1, 8)
\end{lstlisting}

How it works: keep applying f until p approves -- iteration to a goal or a FIXED POINT, as a loop rather than recursion, so depth is unbounded. The state can be anything; a tuple lets one loop carry several facts at once, like the collatz example's (value, steps).


Why it earns its place: some processes have no list to fold over -- they run "until stable". Newton's isqrt converges; bellman\_ford relaxes every edge until the distance map stops changing (the textbook fixed point; g7's hint calls it by name). until states such loops declaratively: here is the step, here is done.

\begin{lstlisting}[style=ex]
>>> until(lambda n: n * n >= 30, succ, 0)      # smallest n with n squared >= 30
6
>>> until(lambda xs: xs == sorted(xs), sorted, [3, 1, 2])
[1, 2, 3]
\end{lstlisting}

Remember it as: while-not-done, promoted to an expression that returns its answer.


Roll your own: while not p(x): x = f(x); return x. Iteration promoted to a value-returning expression -- and a loop rather than recursion on purpose, because a fixed point may take unboundedly many steps.

\begin{lstlisting}[style=ex]
def until(p, f, x):  # until p f x -- loop, not recursion (unbounded depth)
    while not p(x):
        x = f(x)
    return x
\end{lstlisting}
\section{Homework}
Work in order: drills cement each tool, applies combine them,
challenges are interview-grade. Write each solution in a scratch file with
the given doctests pasted in, and let \texttt{python3 -m doctest} referee.
\subsection*{15.1\quad Tribonacci\hfill\textnormal{\textit{drill}}}
The tribonacci sequence starts 0, 1, 1 and every later term is the sum of the previous THREE. Write trib(n) with @memo so large n is instant. State the recurrence before coding.

\noindent\texttt{trib(n: int) -> int}
\begin{lstlisting}[style=ex]
>>> trib(0), trib(1), trib(2)
(0, 1, 1)
>>> trib(4)
4
>>> trib(10)
149
>>> trib(25)
1389537
\end{lstlisting}
\subsection*{15.2\quad First power of two\hfill\textnormal{\textit{drill}}}
Using until, return the smallest power of two that is >= n (n >= 1). No loops of your own: state the stop predicate and the step, let until iterate.

\noindent\texttt{first\_pow2(n: int) -> int}
\begin{lstlisting}[style=ex]
>>> first_pow2(1)
1
>>> first_pow2(100)
128
>>> first_pow2(128)
128
>>> first_pow2(1025)
2048
\end{lstlisting}
\subsection*{15.3\quad Digital root\hfill\textnormal{\textit{drill}}}
The digital root of n repeatedly replaces n with the sum of its digits until a single digit remains. Write it with until -- the fixed point is 'already one digit'.

\noindent\texttt{digital\_root(n: int) -> int}
\begin{lstlisting}[style=ex]
>>> digital_root(0)
0
>>> digital_root(7)
7
>>> digital_root(9875)
2
>>> digital_root(99999)
9
\end{lstlisting}
\subsection*{15.4\quad Staircase, three strides\hfill\textnormal{\textit{drill}}}
You climb a staircase of n steps taking 1, 2 or 3 steps at a time. Count the ways to reach the top (distinct ordered stride sequences) with a memoized recursion; ways(0) is 1 (stand still), negative overshoot is 0.

\noindent\texttt{steps3(n: int) -> int}
\begin{lstlisting}[style=ex]
>>> steps3(0)
1
>>> steps3(1)
1
>>> steps3(4)
7
>>> steps3(10)
274
\end{lstlisting}
\subsection*{15.5\quad Reading the DP table\hfill\textnormal{\textit{drill}}}
Write fib with @memo inside a helper, compute fib(n), and return the pair (fib(n), number\_of\_cached\_states). The cache IS the DP table; its size is the state count that complexity analysis asks you to say out loud.

\noindent\texttt{fib\_with\_table(n: int) -> tuple[int, int]}
\begin{lstlisting}[style=ex]
>>> fib_with_table(0)
(0, 1)
>>> fib_with_table(10)
(55, 11)
>>> fib_with_table(20)
(6765, 21)
\end{lstlisting}
\subsection*{15.6\quad House robber\hfill\textnormal{\textit{apply}}}
A row of houses holds these amounts of cash. You cannot rob two ADJACENT houses. Return the maximum total you can take. Classic interview DP: at each house, rob it and skip a neighbour, or walk past.

\noindent\texttt{rob(xs: list[int]) -> int}
\begin{lstlisting}[style=ex]
>>> rob([])
0
>>> rob([5])
5
>>> rob([1, 2, 3, 1])
4
>>> rob([2, 7, 9, 3, 1])
12
\end{lstlisting}
\subsection*{15.7\quad Decode ways\hfill\textnormal{\textit{apply}}}
A message of letters A-Z was encoded as digits (A=1 .. Z=26) and concatenated. Count the ways to decode a digit string: '12' is 2 ('AB' or 'L'). A leading zero at any decision point kills that branch; the empty string counts as one way.

\noindent\texttt{decode\_ways(s: str) -> int}
\begin{lstlisting}[style=ex]
>>> decode_ways('')
1
>>> decode_ways('12')
2
>>> decode_ways('226')
3
>>> decode_ways('06')
0
>>> decode_ways('11106')
2
\end{lstlisting}
\subsection*{15.8\quad Jump game\hfill\textnormal{\textit{apply}}}
xs[i] is the maximum jump length from index i. Starting at index 0, can you reach the last index? Solve it as one fold over enum(xs) carrying the furthest reachable index -- poison the accumulator when an index is unreachable.

\noindent\texttt{can\_jump(xs: list[int]) -> bool}
\begin{lstlisting}[style=ex]
>>> can_jump([0])
True
>>> can_jump([2, 3, 1, 1, 4])
True
>>> can_jump([3, 2, 1, 0, 4])
False
>>> can_jump([1, 0, 1])
False
\end{lstlisting}
\subsection*{15.9\quad Perfect squares\hfill\textnormal{\textit{apply}}}
Return the fewest perfect squares (1, 4, 9, 16, ...) summing exactly to n. 12 -> 3 (4+4+4), 13 -> 2 (4+9). Memoize on the remaining amount; isqrt bounds the candidates.

\noindent\texttt{num\_squares(n: int) -> int}
\begin{lstlisting}[style=ex]
>>> num_squares(0)
0
>>> num_squares(1)
1
>>> num_squares(12)
3
>>> num_squares(13)
2
>>> num_squares(7)
4
\end{lstlisting}
\subsection*{15.10\quad Growth to target\hfill\textnormal{\textit{apply}}}
A town of pop people grows rate percent per year (integer floor) and then gains influx newcomers. Using until with a (population, years) state tuple, return how many whole years until the population reaches target (0 if already there).

\noindent\texttt{years\_to\_target(pop: int, rate: int, influx: int, target: int) -> int}
\begin{lstlisting}[style=ex]
>>> years_to_target(500, 5, 10, 400)
0
>>> years_to_target(100, 10, 0, 200)
8
>>> years_to_target(1, 0, 10, 51)
5
\end{lstlisting}
\subsection*{15.11\quad Longest palindromic subsequence\hfill\textnormal{\textit{challenge}}}
Return the length of the longest subsequence of s (not necessarily contiguous) that reads the same forwards and backwards. 'bbbab' -> 4 ('bbbb'). Memoize on the (i, j) window: matching ends extend the palindrome by two; otherwise drop one end.

\noindent\texttt{lps(s: str) -> int}
\begin{lstlisting}[style=ex]
>>> lps('')
0
>>> lps('a')
1
>>> lps('bbbab')
4
>>> lps('cbbd')
2
>>> lps('abcdef')
1
\end{lstlisting}
\subsection*{15.12\quad Word break\hfill\textnormal{\textit{challenge}}}
Can s be segmented into a sequence of words from the given list (words reusable)? 'leetcode' with ['leet', 'code'] -> True. Memoize on the start index; isPrefixOf tests each candidate word against the remaining suffix.

\noindent\texttt{word\_break(s: str, words: list[str]) -> bool}
\begin{lstlisting}[style=ex]
>>> word_break('', ['a'])
True
>>> word_break('leetcode', ['leet', 'code'])
True
>>> word_break('applepenapple', ['apple', 'pen'])
True
>>> word_break('catsandog', ['cats', 'dog', 'sand', 'and', 'cat'])
False
>>> word_break('aaaaaaa', ['aaa', 'aaaa'])
True
\end{lstlisting}
\clearpage
\chapter{Combinators Capstone}

The functions about functions come FIRST in the file and LAST in a person:
you have been using them all course, and now they get named as a
discipline. \texttt{curry} and \texttt{partial} stage arguments;
\texttt{flip} and \texttt{on} adapt argument order and route comparisons
through keys; \texttt{compose} makes a pipeline a value. The capstone skill
is judgement -- writing point-free when it clarifies, and refusing to when
it does not.
\section{Concepts}
\subsection{\texttt{id}}\label{def:id}
the identity function: returns its argument unchanged. Sounds useless; is everywhere. It is the "do nothing" you pass where a transformation is required -- a sort key meaning "the value itself", the seed of compose.
\begin{lstlisting}[style=ex]
>>> id(42)
42
>>> sortOn(id, [3, 1, 2])
[1, 2, 3]
\end{lstlisting}

Roll your own --- the definition:

\begin{lstlisting}[style=ex]
id = lambda x: x  # shadows builtin id() by design
\end{lstlisting}
\subsection{\texttt{const}}\label{def:const}
freezes a value into a function: const(x) returns a function that ignores its argument and always answers x. Use it to fill a callback slot with a constant.
\begin{lstlisting}[style=ex]
>>> five = const(5)
>>> five("anything"), five(99)
(5, 5)
>>> map_(const('*'), [1, 2, 3])
['*', '*', '*']
\end{lstlisting}

Roll your own --- the definition:

\begin{lstlisting}[style=ex]
const = lambda x: lambda _: x
\end{lstlisting}
\subsection{\texttt{flip}}\label{def:flip}
swaps the two arguments of a function. An adapter: the function you have takes (a, b), the slot you are filling supplies (b, a). This is how the prelude derives foldr from foldl.
\begin{lstlisting}[style=ex]
>>> sub = lambda a, b: a - b
>>> sub(10, 1), flip(sub)(10, 1)
(9, -9)
>>> flip(pow)(3, 2)                     # pow(2, 3)
8
\end{lstlisting}

How it works: flip(f) returns a new function that hands its two arguments to f in the opposite order. Nothing is computed at flip time -- it is pure plumbing, applied when the data arrives in the "wrong" order for the function you already have.

\begin{lstlisting}[style=ex]
>>> pair = lambda a, b: (a, b)
>>> flip(pair)(1, 2)
(2, 1)
\end{lstlisting}

Why it earns its place: the prelude's own foldr is the showcase. foldl's combiner takes (accumulator, element); a right fold's combiner takes (element, accumulator). Rather than build a second fold machine, foldr = foldl of the FLIPPED combiner over the reversed list -- one adapter turned an existing engine into its mirror image.

\begin{lstlisting}[style=ex]
>>> f = lambda x, acc: '(' + x + acc + ')'
>>> foldr(f, "abc", "*") == foldl(flip(f), list(reversed("abc")), "*")
True
\end{lstlisting}

Remember it as: when argument order is the only thing wrong, adapt -- don't rewrite.


Roll your own --- the definition:

\begin{lstlisting}[style=ex]
flip = lambda f: (lambda x, y: f(y, x))  # flip f x y = f y x
\end{lstlisting}
\subsection{\texttt{curry}}\label{def:curry}
splits a two-argument function into two one-argument stages, so you can supply the first argument now and the second later.
\begin{lstlisting}[style=ex]
>>> add = lambda a, b: a + b
>>> curry(add)(1)(2)
3
>>> add_ten = curry(add)(10)
>>> add_ten(5)
15
\end{lstlisting}

How it works: read the three lambdas inside-out, then watch a call peel one layer at a time. Each call binds ONE argument by closure -- after curry(add)(10), the 10 lives in the returned function's captured environment, and that half-applied function is a value you can name, store, and pass.

\begin{lstlisting}[style=ex]
>>> curry(add)(10)(5)
15
>>> add_ten = curry(add)(10)
>>> add_ten(5), add_ten(90)
(15, 100)
\end{lstlisting}

Why it earns its place: the prelude's pipelines -- map\_, filter\_, find, compose -- all want ONE-argument functions in their function slot, while real logic usually takes two ("multiply by k", "is x greater than n?"). Currying is the adapter: close over the configuration argument now, hand the pipeline the unary stage it wants.

\begin{lstlisting}[style=ex]
>>> gt = curry(lambda n, x: x > n)
>>> filter_(gt(10), [4, 11, 8, 40])
[11, 40]
>>> scale = curry(lambda k, x: k * x)
>>> map_(scale(3), [1, 2, 3])
[3, 6, 9]
\end{lstlisting}

The Haskell connection: in Haskell EVERY function is curried automatically -- add :: Int -> Int -> Int really means Int -> (Int -> Int), so add 10 is ordinary code and map (add 10) xs falls out for free. Every arrow in a Haskell type is another layer of this doll.


Versus partial: both bind early. partial binds any number of arguments, all at once, when you have them; curry restructures the function so binding happens one call at a time, decided later by whoever holds it. In Python, partial is the cheaper everyday tool; curry is the pipeline-shaper and the reading key for Haskell signatures.


Remember it as: curry turns one two-argument function into two one-argument functions -- and one-argument functions are the only thing pipelines eat.


Roll your own: three lambdas, each returning the next -- every layer is a closure capturing one more argument until the innermost body finally has both. No state, no storage: the captured environment IS the data structure.

\begin{lstlisting}[style=ex]
curry = lambda f: lambda x: lambda y: f(x, y)
\end{lstlisting}
\subsection{\texttt{uncurry}}\label{def:uncurry}
the reverse adaptation: a two-argument function becomes a one-argument function of a PAIR. Reach for it when your data is already (a, b) tuples and your function is not.
\begin{lstlisting}[style=ex]
>>> uncurry(add)((1, 2))
3
>>> map_(uncurry(lambda a, b: a * b), [(2, 3), (4, 5)])
[6, 20]
\end{lstlisting}

Roll your own --- the definition:

\begin{lstlisting}[style=ex]
uncurry = lambda f: lambda p: f(*p)
\end{lstlisting}
\subsection{\texttt{partial}}\label{def:partial}
freezes any number of leading arguments, returning a function that takes the rest. Like curry but takes the arguments now, all at once.
\begin{lstlisting}[style=ex]
>>> partial(pow, 2)(10)                 # pow(2, 10)
1024
>>> clamp_floor = partial(max, 0)       # "never below zero"
>>> clamp_floor(-5), clamp_floor(7)
(0, 7)
\end{lstlisting}

Roll your own --- the definition:

\begin{lstlisting}[style=ex]
partial = lambda f, *bound: lambda *args: f(*bound, *args)
\end{lstlisting}
\subsection{\texttt{on}}\label{def:on}
builds a two-argument function that first sends BOTH arguments through a key function: on(f, key)(x, y) == f(key(x), key(y)). Its natural habitat is comparators: compare two records by one field.
\begin{lstlisting}[style=ex]
>>> on(add, len)("ab", "c")             # len("ab") + len("c")
3
>>> same_length = on(lambda a, b: a == b, len)
>>> same_length("abc", "xyz"), same_length("ab", "abc")
(True, False)
>>> sortBy(on(lambda a, b: a - b, snd), [(1, 9), (2, 3)])
[(2, 3), (1, 9)]
\end{lstlisting}

How it works: on(f, key) builds a two-argument function that routes BOTH inputs through key before combining them with f: on(f, key)(x, y) == f(key(x), key(y)). Combination logic and projection logic stay separate, so each is reusable on its own.

\begin{lstlisting}[style=ex]
>>> shorter = on(lambda a, b: a < b, len)
>>> shorter("ab", "abc"), shorter("abc", "ab")
(True, False)
\end{lstlisting}

Why it earns its place: comparator factories. Any "compare records by one field" collapses to on(compare, field) -- Haskell's sortBy (compare 'on' snd) is exactly sortBy(on(cmp, snd)) here, and reads as its own documentation.

\begin{lstlisting}[style=ex]
>>> sortBy(on(lambda a, b: a - b, snd), [('a', 9), ('b', 1), ('c', 5)])
[('b', 1), ('c', 5), ('a', 9)]
\end{lstlisting}

Versus sortOn: when a per-element key exists, sortOn(key, xs) is simpler and faster (the key runs once per element). on shines when an API demands a genuine two-argument comparator, or when the combining function is something other than comparison.


Remember it as: compare the keys, not the things.


Roll your own --- the definition:

\begin{lstlisting}[style=ex]
on = lambda f, g: lambda x, y: f(g(x), g(y))  # Data.Function: combine via a key -- cmp `on` fst
\end{lstlisting}
\subsection{\texttt{compose}}\label{def:compose}
fold any number of functions into one; the RIGHTMOST runs first, like mathematical composition f(g(x)).
\begin{lstlisting}[style=ex]
>>> compose(str.strip, str.upper)("  hi  ")
'HI'
>>> f = compose(succ, succ, succ)
>>> f(0)
3
\end{lstlisting}

How it works: a fold over the functions themselves -- neighbours merge into lambda x: f(g(x)), seeded with the identity, so compose() with no arguments is id. The RIGHTMOST function runs first, matching mathematical (f . g)(x) = f(g(x)).

\begin{lstlisting}[style=ex]
>>> compose()(42) == id(42)
True
>>> compose(len, str.strip)("  hi  ")     # strip first, then len
2
\end{lstlisting}

Why it earns its place: it turns a pipeline into a VALUE. A composed function can be named, mapped, stored in a rules table -- processing steps become data you can shuffle. The constraint is the one curry answers: every stage must be unary, which is why compose-heavy code leans on curry and partial to shape its stages.

\begin{lstlisting}[style=ex]
>>> normalize = compose(str.lower, str.strip)
>>> map_(normalize, ["  Hi ", " YO "])
['hi', 'yo']
\end{lstlisting}

Remember it as: foldr (.) id -- a pipeline you can hold in your hand.


Roll your own: it IS foldr (.) id -- foldr over the functions with x as the seed, each stage applied to the accumulated result: foldr(lambda f, acc: f(acc), fns, x). Because the prelude's foldr is an iterative loop underneath, stack depth stays constant no matter how many stages you compose.

\begin{lstlisting}[style=ex]
compose = lambda *fns: lambda x: foldr(lambda f, acc: f(acc), fns, x)  # foldr (.) id: RIGHTMOST first
\end{lstlisting}
\section{Homework}
Work in order: drills cement each tool, applies combine them,
challenges are interview-grade. Write each solution in a scratch file with
the given doctests pasted in, and let \texttt{python3 -m doctest} referee.
\subsection*{16.1\quad Compose three\hfill\textnormal{\textit{drill}}}
Use compose to build a single function that strips a string, lowercases it, then reverses it -- in that order of effect. Remember compose runs its RIGHTMOST function first, so the last stage you write is the first that happens.

\noindent\texttt{clean\_reverse: str -> str}
\begin{lstlisting}[style=ex]
>>> clean_reverse('  Hello  ')
'olleh'
>>> clean_reverse('ABC')
'cba'
>>> clean_reverse('')
''
\end{lstlisting}
\subsection*{16.2\quad Flip subtract\hfill\textnormal{\textit{drill}}}
Given the two-argument function sub(a, b) = a - b, use flip to make a function that subtracts the FIRST argument from the second. flip swaps argument order and nothing else.

\noindent\texttt{rsub: (int, int) -> int   (rsub(a, b) == b - a)}
\begin{lstlisting}[style=ex]
>>> rsub(1, 10)
9
>>> rsub(10, 1)
-9
>>> rsub(5, 5)
0
\end{lstlisting}
\subsection*{16.3\quad Add via curry\hfill\textnormal{\textit{drill}}}
Curry the function add(a, b) = a + b, then use the curried form to build add\_ten, a one-argument function that adds 10. Show that binding the first argument yields a reusable function.

\noindent\texttt{add\_ten: int -> int}
\begin{lstlisting}[style=ex]
>>> add_ten(5)
15
>>> add_ten(-10)
0
>>> add_ten(0)
10
\end{lstlisting}
\subsection*{16.4\quad Constant column\hfill\textnormal{\textit{drill}}}
Use const and map\_ to replace every element of a list with the string 'x', regardless of its value. const(v) is the function that ignores its argument and always answers v.

\noindent\texttt{mask(xs: list) -> list[str]}
\begin{lstlisting}[style=ex]
>>> mask([1, 2, 3])
['x', 'x', 'x']
>>> mask([])
[]
>>> mask(['a', None, 7])
['x', 'x', 'x']
\end{lstlisting}
\subsection*{16.5\quad Partial power\hfill\textnormal{\textit{drill}}}
Use partial to freeze the base of pow, producing powers\_of\_two(exp) == 2 ** exp. partial(f, *leading) fixes the first arguments now and takes the rest later.

\noindent\texttt{powers\_of\_two: int -> int}
\begin{lstlisting}[style=ex]
>>> powers_of_two(0)
1
>>> powers_of_two(10)
1024
>>> powers_of_two(5)
32
\end{lstlisting}
\subsection*{16.6\quad Sort records by a field\hfill\textnormal{\textit{apply}}}
You have (name, age) tuples. Sort them oldest-first, using on to build the comparator from a plain numeric compare and the snd accessor, fed through sortBy. on(f, key)(x, y) routes BOTH arguments through key before f.

\noindent\texttt{by\_age\_desc(people: list[tuple]) -> list[tuple]}
\begin{lstlisting}[style=ex]
>>> by_age_desc([('ada', 36), ('bob', 41), ('cy', 19)])
[('bob', 41), ('ada', 36), ('cy', 19)]
>>> by_age_desc([])
[]
>>> by_age_desc([('x', 5)])
[('x', 5)]
\end{lstlisting}
\subsection*{16.7\quad Traffic-light rules table\hfill\textnormal{\textit{apply}}}
Classify an integer as 'neg', 'zero', or 'pos' using a RULES table: a list of (predicate, result) pairs tried in order, with find picking the first match. The catch-all arm is const(otherwise) -- since table guards are APPLIED to the value, the always-true guard must itself be a function of the value.

\noindent\texttt{classify(n: int) -> str}
\begin{lstlisting}[style=ex]
>>> classify(-4)
'neg'
>>> classify(0)
'zero'
>>> classify(9)
'pos'
>>> map_(classify, [-1, 0, 1])
['neg', 'zero', 'pos']
\end{lstlisting}
\subsection*{16.8\quad Pipeline from a list of steps\hfill\textnormal{\textit{apply}}}
Given a list of one-argument functions, build a single function that runs them left-to-right (first in the list runs first) and apply it to x. Note this is the OPPOSITE order from compose, so reverse before composing.

\noindent\texttt{pipe(steps: list, x)}
\begin{lstlisting}[style=ex]
>>> pipe([lambda a: a + 1, lambda a: a * 2], 5)
12
>>> pipe([], 7)
7
>>> pipe([str.strip, str.upper], '  hi ')
'HI'
\end{lstlisting}
\subsection*{16.9\quad Group-and-count, point-free key\hfill\textnormal{\textit{apply}}}
Count how many words fall into each length bucket. Build the (key, value) pairs with the length as key and 1 as value, then fold them into a dict with fromListWith and addition. Use partial or a small helper to make the pair-builder read cleanly.

\noindent\texttt{length\_histogram(words: list[str]) -> dict[int, int]}
\begin{lstlisting}[style=ex]
>>> length_histogram(['a', 'bb', 'cc', 'ddd'])
{1: 1, 2: 2, 3: 1}
>>> length_histogram([])
{}
>>> length_histogram(['x', 'y', 'z'])
{1: 3}
\end{lstlisting}
\subsection*{16.10\quad Equality up to a key\hfill\textnormal{\textit{apply}}}
Write same\_by(key, a, b): true when a and b are equal AFTER passing through key. Build it with on so the projection lives in one place. Then a doctest shows it grouping case-insensitively.

\noindent\texttt{same\_by(key, a, b) -> bool}
\begin{lstlisting}[style=ex]
>>> same_by(str.lower, 'ABC', 'abc')
True
>>> same_by(len, 'ab', 'cd')
True
>>> same_by(len, 'ab', 'c')
False
>>> groupBy(partial(same_by, str.lower), ['a', 'A', 'b'])
[['a', 'A'], ['b']]
\end{lstlisting}
\subsection*{16.11\quad Compose a validation pipeline\hfill\textnormal{\textit{challenge}}}
Each validator is a function that returns None on success or an error string on failure. Run a list of validators over a value and return the FIRST error, or None if all pass. Build the check with map\_ and find; a value passes when find turns up no error.

\noindent\texttt{first\_error(validators: list, x)}
\begin{lstlisting}[style=ex]
>>> pos = lambda n: None if n > 0 else 'not positive'
>>> even_v = lambda n: None if n % 2 == 0 else 'not even'
>>> first_error([pos, even_v], 4) is None
True
>>> first_error([pos, even_v], -3)
'not positive'
>>> first_error([pos, even_v], 3)
'not even'
>>> first_error([], 0) is None
True
\end{lstlisting}
\subsection*{16.12\quad Curried lookup chain\hfill\textnormal{\textit{challenge}}}
Build get\_in(keys) -- a function that, given nested dicts, walks the key path and returns the leaf, or None if any key is missing. get\_in is CURRIED: get\_in(['a', 'b']) returns a function you can reuse on many dicts. Fold the keys with a safe step that stops at None.

\noindent\texttt{get\_in(keys: list) -> (dict -> value | None)}
\begin{lstlisting}[style=ex]
>>> deep = {'a': {'b': {'c': 42}}}
>>> get_in(['a', 'b', 'c'])(deep)
42
>>> get_in(['a', 'x'])(deep) is None
True
>>> get_in([])(deep) == deep
True
>>> pluck = get_in(['a', 'b'])
>>> pluck(deep)
{'c': 42}
\end{lstlisting}
\clearpage
\part{Applied Labs: The Nine Problem Classes}
\chapter*{About the labs}
\addcontentsline{toc}{chapter}{About the labs}
Seventy-four problems, each a self-contained practice file in the repository
(\texttt{fpython/g*/}) carrying its needed snippets inline. Statements,
contracts, hints and doctests are printed here; solve them in the repository files, where the doctest runner referees.
\chapter{Arith And Unfolds}

\section{sign: Sign of a number}
\begin{lstlisting}[style=ex]
Return -1 for negative, 0 for zero, 1 for positive.

Contract:
    sign(x: int | float) -> int

Hint:
    The prelude ships this exact function: signum, one comparison
    subtraction, no branches. Recognize it, don't reinvent it.
\end{lstlisting}
\begin{lstlisting}[style=ex]
>>> sign(-7)
-1
>>> sign(0)
0
>>> sign(3)
1
>>> sign(-0.5)
-1
\end{lstlisting}

\section{rect\_area: Rectangle area}
\begin{lstlisting}[style=ex]
Contract:
    rect_area(w, h) -> w * h  (result type follows the inputs)

Hint:
    Deliberate overkill: area is the numeric fold, product([w, h]).
    See the smallest possible fold once; recognize every fold after.
\end{lstlisting}
\begin{lstlisting}[style=ex]
>>> rect_area(3, 4)
12
>>> rect_area(2.5, 4)
10.0
>>> rect_area(0, 99)
0
\end{lstlisting}

\section{reverse\_digits: Reverse an integer's digits}
\begin{lstlisting}[style=ex]
Negatives keep their sign; trailing zeros vanish.

Contract:
    reverse_digits(n: int) -> int

Hint:
    unfoldr peels the digits with (m % 10, m // 10) -- least
    significant first, which IS reversed order. Fold them back into a
    number with acc * 10 + d; signum reapplies the sign.
\end{lstlisting}
\begin{lstlisting}[style=ex]
>>> reverse_digits(123)
321
>>> reverse_digits(-45)
-54
>>> reverse_digits(120)
21
>>> reverse_digits(7)
7
\end{lstlisting}

\section{to\_base: Render n in base b}
\begin{lstlisting}[style=ex]
Non-negative n, base 2..16, uppercase digits.

Contract:
    to_base(n: int, b: int) -> str

Hint:
    unfoldr emits digits low-to-high -- "0123456789ABCDEF"[m % b],
    seed m // b, stop at 0. Reverse and join; `or "0"` covers zero
    (its unfold is empty).
\end{lstlisting}
\begin{lstlisting}[style=ex]
>>> to_base(255, 16)
'FF'
>>> to_base(5, 2)
'101'
>>> to_base(0, 8)
'0'
>>> to_base(100, 10)
'100'
>>> to_base(31, 16)
'1F'
\end{lstlisting}

\section{to\_roman: Integer to Roman numeral, 1 <= n <= 3999}
\begin{lstlisting}[style=ex]
Contract:
    to_roman(n: int) -> str

Hint:
    A left fold over the value table (subtractive pairs included:
    900 'CM', 40 'XL', ...) carrying (remaining, out): each step
    appends sym * (remaining // val) and keeps remaining % val.
\end{lstlisting}
\begin{lstlisting}[style=ex]
>>> to_roman(1994)
'MCMXCIV'
>>> to_roman(9)
'IX'
>>> to_roman(58)
'LVIII'
>>> to_roman(3000)
'MMM'
>>> to_roman(1)
'I'
\end{lstlisting}

\section{from\_roman: Parse a valid Roman numeral}
\begin{lstlisting}[style=ex]
Contract:
    from_roman(s: str) -> int

Hint:
    pairwise gives each symbol its right neighbor; a value counts
    NEGATIVE when its neighbor is bigger ('IV' = -1 + 5). Pad the
    value list with a trailing 0 so the last symbol gets a neighbor
    too, then sum.
\end{lstlisting}
\begin{lstlisting}[style=ex]
>>> from_roman('MCMXCIV')
1994
>>> from_roman('IX')
9
>>> from_roman('LVIII')
58
>>> from_roman('III')
3
>>> from_roman('MMM')
3000
\end{lstlisting}
\chapter{Lists And Strings}

\section{greet: Greeting string}
\begin{lstlisting}[style=ex]
Contract:
    greet(name: str) -> 'Hello, <name>!'

Hint:
    unwords joins words with single spaces; the bang concatenates.
    Trivial on purpose -- the warmup proves the pipeline.
\end{lstlisting}
\begin{lstlisting}[style=ex]
>>> greet("Ada")
'Hello, Ada!'
>>> greet("world")
'Hello, world!'
\end{lstlisting}

\section{count\_vowels: Count vowels, case-insensitive}
\begin{lstlisting}[style=ex]
Contract:
    count_vowels(s: str) -> int

Hint:
    len . filter_ (in "aeiou") . lower -- a filter consumed by len is
    the counting idiom.
\end{lstlisting}
\begin{lstlisting}[style=ex]
>>> count_vowels("Programming")
3
>>> count_vowels("AEIOU")
5
>>> count_vowels("xyz")
0
>>> count_vowels("")
0
\end{lstlisting}

\section{fizzbuzz: FizzBuzz as data}
\begin{lstlisting}[style=ex]
For 1..n: multiples of 3 -> "Fizz", of 5 -> "Buzz", of both ->
"FizzBuzz", else the number as a string.

Contract:
    fizzbuzz(n: int) -> list[str]

Hint:
    Solve it BOTH ways and compare -- the contrast is the lesson:
    1. COMPOSITION OF ARMS: map_ one per-number rule over
       range(1, n+1), building the word by concatenating guards --
       "Fizz" * (i % 3 == 0) + "Buzz" * (i % 5 == 0) -- with
       `or str(i)` catching the empty case. Overlaps compose:
       15 earns "FizzBuzz" from two rules for free.
    2. FIRST-MATCH SELECTION: a RULES table of (predicate, result)
       pairs, find picking the first hit; const(otherwise) is the
       catch-all arm. Overlaps need their own row (15 goes first),
       but the ladder is now DATA -- appendable, testable rule by
       rule. Which generalizes better depends on whether cases
       overlap.
\end{lstlisting}
\begin{lstlisting}[style=ex]
>>> fizzbuzz(5)
['1', '2', 'Fizz', '4', 'Buzz']
>>> fizzbuzz(15)[-1]
'FizzBuzz'
>>> fizzbuzz(0)
[]
\end{lstlisting}

\section{word\_lengths: Map each word to its length}
\begin{lstlisting}[style=ex]
Contract:
    word_lengths(sentence: str) -> dict[str, int]
    First-appearance key order; repeats collapse harmlessly.

Hint:
    words splits; zip_ the words against map_ len of them; dict() the
    pairs. Duplicate keys overwrite with the same value -- free dedup.
\end{lstlisting}
\begin{lstlisting}[style=ex]
>>> word_lengths("the quick fox")
{'the': 3, 'quick': 5, 'fox': 3}
>>> word_lengths("go go go")
{'go': 2}
>>> word_lengths("")
{}
\end{lstlisting}

\section{flatten: Flatten one level of nesting}
\begin{lstlisting}[style=ex]
Contract:
    flatten(lists: list[list]) -> list

Hint:
    This IS the prelude's concat. One name, zero work.
\end{lstlisting}
\begin{lstlisting}[style=ex]
>>> flatten([[1, 2], [3], [4, 5]])
[1, 2, 3, 4, 5]
>>> flatten([[], [1], []])
[1]
>>> flatten([])
[]
\end{lstlisting}

\section{caesar: Caesar cipher, letters only, case preserved}
\begin{lstlisting}[style=ex]
Contract:
    caesar(s: str, k: int) -> str

Hint:
    map_ a per-char shift over the string; the wrap is mod 26 on
    ord() offsets; join the pieces. One helper, one map.
\end{lstlisting}
\begin{lstlisting}[style=ex]
>>> caesar("abc", 2)
'cde'
>>> caesar("XYZ", 3)
'ABC'
>>> caesar("a-b", 1)
'b-c'
>>> caesar("Hello, World!", 13)
'Uryyb, Jbeyq!'
>>> caesar("abc", 0)
'abc'
\end{lstlisting}

\section{contains\_duplicate: NeetCode: Contains Duplicate}
\begin{lstlisting}[style=ex]
Contract:
    contains_duplicate(nums: list[int]) -> bool
    True iff any value appears at least twice.

Hint:
    nub keeps first occurrences; duplicates exist iff nub shrinks the
    list.
\end{lstlisting}
\begin{lstlisting}[style=ex]
>>> contains_duplicate([1, 2, 3, 1])
True
>>> contains_duplicate([1, 2, 3, 4])
False
>>> contains_duplicate([])
False
\end{lstlisting}

\section{chunks: Split a list into pieces of length k}
\begin{lstlisting}[style=ex]
Last piece may be shorter.

Contract:
    chunks(xs: list, k: int) -> list[list]

Hint:
    splitAt already returns (front, rest) == (value, next_seed), so
    chunks IS unfoldr of splitAt; null(rest) is the stop. One line.
    (You are deriving the prelude's chunksOf -- deriving a tool once
    is why it sticks.)
\end{lstlisting}
\begin{lstlisting}[style=ex]
>>> chunks([1, 2, 3, 4, 5], 2)
[[1, 2], [3, 4], [5]]
>>> chunks([1, 2, 3, 4], 2)
[[1, 2], [3, 4]]
>>> chunks([], 3)
[]
>>> chunks([1], 5)
[[1]]
\end{lstlisting}

\section{permutations: All distinct orderings, sorted}
\begin{lstlisting}[style=ex]
Contract:
    permutations(s: str) -> list[str]

Hint:
    concatMap over the choice of first character: pick s[i], prepend
    it to every permutation of the rest. nub kills duplicates from
    repeated letters; sort at the end.
\end{lstlisting}
\begin{lstlisting}[style=ex]
>>> permutations("abc")
['abc', 'acb', 'bac', 'bca', 'cab', 'cba']
>>> permutations("aab")
['aab', 'aba', 'baa']
>>> permutations("x")
['x']
>>> permutations("")
['']
\end{lstlisting}

\section{reverse\_words: Reverse the word order}
\begin{lstlisting}[style=ex]
Runs of whitespace collapse (that's what a bare split does).

Contract:
    reverse_words(s: str) -> str

Hint:
    An unwords . reverse . words sandwich -- three names, no loop.
\end{lstlisting}
\begin{lstlisting}[style=ex]
>>> reverse_words("the quick brown fox")
'fox brown quick the'
>>> reverse_words("  hello   world  ")
'world hello'
>>> reverse_words("solo")
'solo'
\end{lstlisting}

\section{transpose\_matrix: Flip rows and columns}
\begin{lstlisting}[style=ex]
Contract:
    transpose_matrix(rows: list[list]) -> list[list]

Hint:
    This IS the prelude's transpose -- like flatten was concat.
    One name, zero work; know WHY zip(*rows) does it.
\end{lstlisting}
\begin{lstlisting}[style=ex]
>>> transpose_matrix([[1, 2, 3], [4, 5, 6]])
[[1, 4], [2, 5], [3, 6]]
>>> transpose_matrix([[5]])
[[5]]
>>> transpose_matrix([[7], [8], [9]])
[[7, 8, 9]]
\end{lstlisting}

\section{rotate\_right: Rotate a list right by k}
\begin{lstlisting}[style=ex]
k may be 0 or larger than the list.

Contract:
    rotate_right(xs: list, k: int) -> list

Hint:
    One splitAt at len - k%len, then glue the halves swapped.
    The k=0 trap: xs[-0:] thinking sneaks in if you slice by -k.
\end{lstlisting}
\begin{lstlisting}[style=ex]
>>> rotate_right([1, 2, 3, 4, 5], 2)
[4, 5, 1, 2, 3]
>>> rotate_right([1, 2], 0)
[1, 2]
>>> rotate_right([1, 2, 3], 5)
[2, 3, 1]
>>> rotate_right([], 3)
[]
\end{lstlisting}
\chapter{Folds Scans Streams}

\section{two\_sum: NeetCode: Two Sum}
\begin{lstlisting}[style=ex]
Contract:
    two_sum(nums: list[int], target: int) -> tuple[int, int]
    Return indices (i, j), i < j, with nums[i] + nums[j] == target.
    Exactly one solution exists. len(nums) >= 2.

Hint:
    A fold over enum(nums) carrying a {value: index} dict.
    The future needs to know: which values have I seen, and where?
\end{lstlisting}
\begin{lstlisting}[style=ex]
>>> two_sum([2, 7, 11, 15], 9)
(0, 1)
>>> two_sum([3, 2, 4], 6)
(1, 2)
>>> two_sum([3, 3], 6)
(0, 1)
\end{lstlisting}

\section{best\_time\_stock: NeetCode: Best Time to Buy and Sell Stock}
\begin{lstlisting}[style=ex]
Contract:
    max_profit(prices: list[int]) -> int
    One buy, one later sell. Best achievable profit, 0 if none.
    len(prices) >= 1.

Hint:
    scanl1 min gives the running minimum -- the best buy so far.
    zipWith (-) prices against it, then maximum. A scan, not a loop.
\end{lstlisting}
\begin{lstlisting}[style=ex]
>>> max_profit([7, 1, 5, 3, 6, 4])
5
>>> max_profit([7, 6, 4, 3, 1])
0
>>> max_profit([5])
0
\end{lstlisting}

\section{maximum\_subarray: NeetCode: Maximum Subarray (Kadane)}
\begin{lstlisting}[style=ex]
Contract:
    max_subarray(nums: list[int]) -> int
    Largest sum of a non-empty contiguous subarray. len(nums) >= 1.

Hint:
    Kadane IS scanl1: best-ending-here = max(x, best_prev + x).
    The scan is the DP table; maximum reads the answer off it.
\end{lstlisting}
\begin{lstlisting}[style=ex]
>>> max_subarray([-2, 1, -3, 4, -1, 2, 1, -5, 4])
6
>>> max_subarray([1])
1
>>> max_subarray([5, 4, -1, 7, 8])
23
>>> max_subarray([-3, -1, -2])
-1
\end{lstlisting}

\section{product\_except\_self: NeetCode: Product of Array Except Self}
\begin{lstlisting}[style=ex]
Contract:
    product_except_self(nums: list[int]) -> list[int]
    out[i] == product of all nums except nums[i]. No division. O(n).
    len(nums) >= 2.

Hint:
    Prefix products = scanl (*) 1, suffix products = scanr (*) 1.
    out = zipWith (*) prefixes suffixes -- offset by one on each side.
\end{lstlisting}
\begin{lstlisting}[style=ex]
>>> product_except_self([1, 2, 3, 4])
[24, 12, 8, 6]
>>> product_except_self([-1, 1, 0, -3, 3])
[0, 0, 9, 0, 0]
\end{lstlisting}

\section{moving\_average: Average of every k-window}
\begin{lstlisting}[style=ex]
Contract:
    moving_average(xs: list, k: int) -> list[float]
    len(xs) >= k >= 1.

Hint:
    Prefix sums = scanl (+) 0 -- one scan, then every window sum is
    pref[i+k] - pref[i]: zipWith (-) the two offset views of the same
    scan, divide by k.
\end{lstlisting}
\begin{lstlisting}[style=ex]
>>> moving_average([1, 2, 3, 4], 2)
[1.5, 2.5, 3.5]
>>> moving_average([1, 1, 1], 3)
[1.0]
>>> moving_average([5], 1)
[5.0]
\end{lstlisting}

\section{pascal: First n rows of Pascal's triangle}
\begin{lstlisting}[style=ex]
Contract:
    pascal(n: int) -> list[list[int]]

Hint:
    The triangle is a stream: iterate the row step from [1], where
    next = zipWith (+) (0:row) (row:0) -- the row against itself,
    shifted. take n consumes the stream.
\end{lstlisting}
\begin{lstlisting}[style=ex]
>>> pascal(4)
[[1], [1, 1], [1, 2, 1], [1, 3, 3, 1]]
>>> pascal(1)
[[1]]
>>> pascal(0)
[]
>>> pascal(6)[-1]
[1, 5, 10, 10, 5, 1]
\end{lstlisting}

\section{climbing\_stairs: NeetCode: Climbing Stairs}
\begin{lstlisting}[style=ex]
Contract:
    climb_stairs(n: int) -> int
    Ways to reach step n taking 1 or 2 steps. n >= 1.

Hint:
    Fibonacci in disguise. iterate the pair step, take the nth --
    an unfold consumed by indexing, no table needed.
\end{lstlisting}
\begin{lstlisting}[style=ex]
>>> climb_stairs(1)
1
>>> climb_stairs(2)
2
>>> climb_stairs(5)
8
>>> climb_stairs(45)
1836311903
\end{lstlisting}

\section{primes\_up\_to: All primes <= n, ascending}
\begin{lstlisting}[style=ex]
Contract:
    primes_up_to(n: int) -> list[int]

Hint:
    Streams, Haskell style: candidates = takeWhile (<= n) count(2);
    keep p when no divisor in takeWhile (<= isqrt p) count(2) divides
    it -- all over the remainders. (The array sieve is faster; this
    version is about reading streams.)
\end{lstlisting}
\begin{lstlisting}[style=ex]
>>> primes_up_to(20)
[2, 3, 5, 7, 11, 13, 17, 19]
>>> primes_up_to(2)
[2]
>>> primes_up_to(1)
[]
>>> len(primes_up_to(100))
25
\end{lstlisting}

\section{totals\_and\_deltas: A scan and its inverse}
\begin{lstlisting}[style=ex]
Contract:
    running_totals(xs: list) -> list   # cumulative sums
    deltas(xs: list) -> list           # each element minus the previous

Hint:
    running_totals IS scanl1 (+). deltas is the pairwise differences.
    They undo each other: prepend the first element to
    deltas(running_totals(xs)) and xs comes back -- say why.
\end{lstlisting}
\begin{lstlisting}[style=ex]
>>> running_totals([3, 1, 4])
[3, 4, 8]
>>> running_totals([])
[]
>>> deltas([5, 3, 8])
[-2, 5]
>>> deltas([7])
[]
>>> ts = running_totals([2, 5, 1]); [2] + deltas(ts) == [2, 5, 1]
True
\end{lstlisting}

\section{paren\_depth: Deepest paren nesting}
\begin{lstlisting}[style=ex]
Input is balanced, only '(' and ')'.

Contract:
    max_depth(s: str) -> int

Hint:
    Depth is a running total: map each paren to +1/-1, scanl (+) 0
    is the entire depth trace, maximum reads off the answer. No
    stack needed -- the scan IS the stack height's history.
\end{lstlisting}
\begin{lstlisting}[style=ex]
>>> max_depth("(()(()))")
3
>>> max_depth("()()")
1
>>> max_depth("")
0
\end{lstlisting}
\chapter{Sorting And Searching}

\section{is\_anagram\_sorted: Anagrams via canonical form}
\begin{lstlisting}[style=ex]
True when a and b use exactly the same letters, ignoring case and
non-letters. (g5 solves the same problem with bags -- do both, argue
which you'd ship.)

Contract:
    is_anagram(a: str, b: str) -> bool

Hint:
    A canonical form makes equality trivial: sort of the cleaned,
    lowered characters. filter_ str.isalpha does the cleaning.
\end{lstlisting}
\begin{lstlisting}[style=ex]
>>> is_anagram("Listen", "Silent")
True
>>> is_anagram("hello", "world")
False
>>> is_anagram("a gentleman", "elegant man")
True
>>> is_anagram("ab", "abb")
False
\end{lstlisting}

\section{second\_largest: Second largest distinct value}
\begin{lstlisting}[style=ex]
At least two distinct values guaranteed.

Contract:
    second_largest(xs: list) -> value

Hint:
    nub then sort; the answer is one from the end.
\end{lstlisting}
\begin{lstlisting}[style=ex]
>>> second_largest([1, 5, 3])
3
>>> second_largest([5, 1, 5, 3])
3
>>> second_largest([2, 1])
1
>>> second_largest([-1, -5, -3])
-3
\end{lstlisting}

\section{binary\_search: Index of target in a sorted list, else -1}
\begin{lstlisting}[style=ex]
Contract:
    binary_search(xs: list, target) -> int

Hint:
    bisect_left finds the leftmost slot where the target COULD be;
    one bounds-check plus one equality test finishes the job. Know
    why bisect_left(xs, x) == len(xs) is possible.
\end{lstlisting}
\begin{lstlisting}[style=ex]
>>> binary_search([1, 3, 5, 7], 5)
2
>>> binary_search([1, 3, 5, 7], 1)
0
>>> binary_search([1, 3, 5, 7], 7)
3
>>> binary_search([1, 3], 2)
-1
>>> binary_search([], 1)
-1
\end{lstlisting}

\section{merge\_intervals: Merge overlapping (and touching) intervals}
\begin{lstlisting}[style=ex]
Intervals are inclusive [s, e]: [1,2] and [2,3] merge to [1,3].

Contract:
    merge_intervals(intervals: list[list[int]]) -> list[list[int]]
    Result sorted by start.

Hint:
    sortOn fst, then one fold: each interval either extends the
    accumulator's last entry or starts a new one.
\end{lstlisting}
\begin{lstlisting}[style=ex]
>>> merge_intervals([[1, 3], [2, 6], [8, 10], [15, 18]])
[[1, 6], [8, 10], [15, 18]]
>>> merge_intervals([[1, 2], [2, 3]])
[[1, 3]]
>>> merge_intervals([[5, 7], [1, 2], [3, 6]])
[[1, 2], [3, 7]]
>>> merge_intervals([[4, 4]])
[[4, 4]]
\end{lstlisting}

\section{photo\_lineup: The reported Anduril prompt, prelude-sized}
\begin{lstlisting}[style=ex]
Back row and front row, same length; can the columns be arranged so
every back person is STRICTLY taller than the front person?

Contract:
    can_lineup(back: list[int], front: list[int]) -> bool

Hint:
    Sort both rows; all over zipWith (>). Be ready to argue why
    sorted-vs-sorted decides it (the exchange argument).
\end{lstlisting}
\begin{lstlisting}[style=ex]
>>> can_lineup([5, 1, 3], [2, 4, 0])
True
>>> can_lineup([1, 2, 3], [1, 2, 3])
False
>>> can_lineup([10, 9], [1, 2])
True
>>> can_lineup([2], [3])
False
\end{lstlisting}

\section{merge\_sort: A real O(n log n) sort, by hand}
\begin{lstlisting}[style=ex]
No sorted(), no .sorted().

Contract:
    merge_sorted(xs: list) -> list  (new list; stable)

Hint:
    splitAt the midpoint, recurse on both halves, and the prelude's
    merge does the two-pointer join. Also retype merge itself from
    memory once -- it's the loop interviews ask for.
\end{lstlisting}
\begin{lstlisting}[style=ex]
>>> merge_sorted([5, 2, 8, 1])
[1, 2, 5, 8]
>>> merge_sorted([])
[]
>>> merge_sorted([3, 3, 1])
[1, 3, 3]
>>> merge_sorted([1])
[1]
\end{lstlisting}

\section{lis: Longest strictly increasing subsequence (length)}
\begin{lstlisting}[style=ex]
Contract:
    lis(xs: list[int]) -> int

Hint:
    Patience sorting: fold the sequence into `tails`, where tails[i]
    is the smallest possible tail of an increasing run of length i+1.
    bisect_left says which pile each element lands on; a landing past
    the end grows the answer. O(n log n) where the obvious DP is
    O(n^2) -- say both.
\end{lstlisting}
\begin{lstlisting}[style=ex]
>>> lis([10, 9, 2, 5, 3, 7, 101, 18])
4
>>> lis([5, 4, 3])
1
>>> lis([1, 2, 3])
3
>>> lis([])
0
>>> lis([7, 7, 7])
1
\end{lstlisting}

\section{russian\_doll: LC 354: Russian Doll Envelopes}
\begin{lstlisting}[style=ex]
Envelope (w, h) fits inside another only when BOTH are strictly
smaller. No rotation. Longest nesting chain.

Contract:
    max_dolls(envelopes: list[tuple[int, int]]) -> int

Hint:
    Sort by (width asc, height DESC), then plain LIS on the heights
    -- g4_07 verbatim. The descending tie-break is the whole trick:
    equal-width envelopes arrive tallest-first, so no two of them can
    ever chain in a strictly-increasing height run. Sorting
    linearized one dimension; LIS handles the other. O(n log n).
\end{lstlisting}
\begin{lstlisting}[style=ex]
>>> max_dolls([(5, 4), (6, 4), (6, 7), (2, 3)])
3
>>> max_dolls([(1, 1), (1, 1), (1, 1)])
1
>>> max_dolls([(3, 4), (3, 5), (3, 6)])
1
>>> max_dolls([])
0
>>> max_dolls([(2, 3)])
1
\end{lstlisting}

\section{pairs\_within\_budget: Count pairs whose sum fits a budget}
\begin{lstlisting}[style=ex]
Pairs are (i, j) with i < j; values may be negative; the count does
not depend on input order -- that sentence licenses sorting.

Contract:
    pairs_within_budget(t: int, xs: list[int]) -> int

Hint:
    Sort first. Then each element asks: how many LATER elements keep
    the sum <= t? bisect_right(ys, t - x) answers in O(log n) --
    subtract the i+1 elements at or before x itself, clip at 0.
\end{lstlisting}
\begin{lstlisting}[style=ex]
>>> pairs_within_budget(5, [3, 1, 4, 2])
4
>>> pairs_within_budget(-1, [-2, -3, 5])
1
>>> pairs_within_budget(100, [1, 2, 3])
3
>>> pairs_within_budget(0, [1, 2])
0
\end{lstlisting}

\section{sensor\_pairing: The screen's domain problem, function form}
\begin{lstlisting}[style=ex]
Two sensor feeds watch the same airspace. A detection is
(id, t, bearing_degrees); ids are sortable strings. Detections a, b
are COMPATIBLE when |a.t - b.t| <= t_tol and their bearing difference
is <= b_tol -- and bearings wrap: 359.5 and 0.3 are 0.8 apart.

Contract:
    pair_detections(a, b, t_tol, b_tol)
        -> (pairs, unpaired_a_ids, unpaired_b_ids)
    Greedy: repeatedly take the lowest-cost compatible pair of still-
    unused detections, cost = |dt| + wrapped bearing diff (exact ties
    by ids). pairs as (a_id, b_id); all three outputs sorted.

Hint:
    Build the compatible candidates as one comprehension over both
    feeds (cost, a_id, b_id, i, j), sort -- the fold that follows
    takes each candidate whose indices are still unused. The wrap is
    min(d, 360 - d) in ONE helper. This is mock_phased.py distilled.
\end{lstlisting}
\begin{lstlisting}[style=ex]
>>> a = [('A0', 10.0, 90.0), ('A1', 20.0, 180.0)]
>>> b = [('B0', 10.2, 91.0), ('B1', 20.1, 179.5)]
>>> pair_detections(a, b, 0.5, 2.0)
([('A0', 'B0'), ('A1', 'B1')], [], [])
>>> pair_detections([('A0', 10.0, 359.5)], [('B0', 10.1, 0.3)], 0.5, 2.0)
([('A0', 'B0')], [], [])
>>> pair_detections([('A0', 0.0, 10.0)], [('B0', 900.0, 200.0)], 0.5, 2.0)
([], ['A0'], ['B0'])
>>> a = [('A0', 10.0, 45.0), ('A1', 30.0, 90.0)]
>>> b = [('B0', 10.5, 44.5), ('B1', 10.6, 46.0), ('B2', 50.0, 200.0)]
>>> pair_detections(a, b, 1.0, 2.0)
([('A0', 'B0')], ['A1'], ['B1', 'B2'])
\end{lstlisting}
\chapter{Bags And Grouping}

\section{valid\_anagram: NeetCode: Valid Anagram}
\begin{lstlisting}[style=ex]
Contract:
    is_anagram(s: str, t: str) -> bool
    True iff t is a permutation of s. Lowercase ASCII.

Hint:
    Two Counter folds; anagram == equal multisets. (Compare with
    g4's sorted-canonical version: same problem, different algebra.)
\end{lstlisting}
\begin{lstlisting}[style=ex]
>>> is_anagram("anagram", "nagaram")
True
>>> is_anagram("rat", "car")
False
>>> is_anagram("", "")
True
\end{lstlisting}

\section{ransom\_note: HackerRank: Ransom Note}
\begin{lstlisting}[style=ex]
Contract:
    can_make(note: list[str], magazine: list[str]) -> bool
    True iff every word in note is available in magazine with at
    least the required multiplicity. Case-sensitive whole words.

Hint:
    Multiset inclusion: Counter(note) <= Counter(magazine), spelled
    as an `all` over the note's counts.
\end{lstlisting}
\begin{lstlisting}[style=ex]
>>> can_make("give me one grand today night".split(), "give me one grand today".split())
False
>>> can_make("two times three is not four".split(), "two times three is not four".split())
True
>>> can_make(["give"], ["Give"])
False
\end{lstlisting}

\section{most\_frequent: Most frequent value; ties -> smallest}
\begin{lstlisting}[style=ex]
Contract:
    most_frequent(xs: list) -> value

Hint:
    Counter, then minOn the tuple key (-count, value). The
    two-level tie-break lives entirely in the key, and minOn is
    O(n) where sortOn + head pays O(n log n) for the same answer.
\end{lstlisting}
\begin{lstlisting}[style=ex]
>>> most_frequent([1, 2, 2, 3, 3])
2
>>> most_frequent([5, 5, 1])
5
>>> most_frequent([3, 1, 2])
1
>>> most_frequent(["b", "a", "b", "a"])
'a'
\end{lstlisting}

\section{top\_k\_words: The k most frequent words}
\begin{lstlisting}[style=ex]
Ordered by count descending, ties by word ascending. At least k
distinct words guaranteed.

Contract:
    top_k_words(k: int, text: str) -> list[tuple[str, int]]

Hint:
    Counter over the words, sortOn (-count, word), take k. The whole
    pipeline is three prelude names deep.
\end{lstlisting}
\begin{lstlisting}[style=ex]
>>> top_k_words(2, "the cat the dog the cat bird")
[('the', 3), ('cat', 2)]
>>> top_k_words(3, "b a a b c c")
[('a', 2), ('b', 2), ('c', 2)]
\end{lstlisting}

\section{group\_words: Group words by first letter}
\begin{lstlisting}[style=ex]
First-appearance order for both keys and members.

Contract:
    group_words(words: list[str]) -> dict[str, list[str]]

Hint:
    THE fromListWith showcase: pairs (first letter, [word]). The
    combiner receives f(new, old) -- Haskell's order -- so KEEP
    arrival order by combining old + new; a bare (+) would reverse
    each group (the classic Haskell gotcha). Keys keep first-seen
    order.
\end{lstlisting}
\begin{lstlisting}[style=ex]
>>> group_words(["ant", "bee", "ape"])
{'a': ['ant', 'ape'], 'b': ['bee']}
>>> group_words(["cat"])
{'c': ['cat']}
>>> group_words([])
{}
\end{lstlisting}

\section{group\_anagrams: NeetCode: Group Anagrams}
\begin{lstlisting}[style=ex]
Contract:
    group_anagrams(words: list[str]) -> list[list[str]]
    Partition words into anagram classes; classes in first-seen order,
    members in input order.

Hint:
    Key = canonical form (sorted letters). Two prelude routes, do
    both: Tree(1, list) with a bare append (autovivification aids the
    write-heavy access), or one fromListWith over (key, [w]) pairs --
    combining old + new, since f receives (new, old); a bare (+)
    would reverse each group. Compare.
\end{lstlisting}
\begin{lstlisting}[style=ex]
>>> group_anagrams(["eat", "tea", "tan", "ate", "nat", "bat"])
[['eat', 'tea', 'ate'], ['tan', 'nat'], ['bat']]
>>> group_anagrams([""])
[['']]
>>> group_anagrams(["a"])
[['a']]
\end{lstlisting}

\section{merge\_sum: Merge two dicts of numbers, summing shared keys}
\begin{lstlisting}[style=ex]
Don't mutate the inputs. Key order: a's keys, then b's new ones.

Contract:
    merge_sum(a: dict, b: dict) -> dict

Hint:
    This IS unionWith with (+): merge, combining shared keys by
    addition. The prelude's bag_union is the SAME merge with max --
    one combinator, different monoid. Spot the difference and say it.
    (fromListWith over the two items() lists concatenated gets there
    too; decide which reads better.)
\end{lstlisting}
\begin{lstlisting}[style=ex]
>>> merge_sum({'a': 1, 'b': 2}, {'b': 3, 'c': 4})
{'a': 1, 'b': 5, 'c': 4}
>>> merge_sum({}, {'x': 1})
{'x': 1}
>>> merge_sum({'x': 1}, {})
{'x': 1}
\end{lstlisting}

\section{rle\_encode: Run-length encode}
\begin{lstlisting}[style=ex]
"aaabbc" -> "a3b2c1"; every run gets a count, even runs of one.

Contract:
    rle_encode(s: str) -> str

Hint:
    The prelude's group gathers runs of CONSECUTIVE equals as bare
    lists. map_ each run to run[0] + str(len(run)), join.
\end{lstlisting}
\begin{lstlisting}[style=ex]
>>> rle_encode("aaabbc")
'a3b2c1'
>>> rle_encode("abc")
'a1b1c1'
>>> rle_encode("aaaa")
'a4'
>>> rle_encode("")
''
\end{lstlisting}

\section{lru\_class: An LRU cache class, dict-only}
\begin{lstlisting}[style=ex]
Capacity at construction; put(key, value); get(key) -> value or None.
Both get and put count as use; puts beyond capacity evict the least
recently used entry.

Contract:
    class LRU(capacity)
        .put(key, value) -> None
        .get(key) -> value | None

Hint:
    The container lesson with no snippet to lean on: dicts remember
    insertion order, so pop + re-insert IS move-to-most-recent, and
    next(iter(d)) is the least recent. That's the whole machine.
\end{lstlisting}
\begin{lstlisting}[style=ex]
>>> c = LRU(2)
>>> c.put('a', 1); c.put('b', 2)
>>> c.get('a')
1
>>> c.put('c', 3)
>>> c.get('b') is None
True
>>> c.get('c')
3
>>> c.get('a')
1
>>> c.put('a', 99); c.get('a')
99
\end{lstlisting}

\section{first\_unique: First non-repeating character}
\begin{lstlisting}[style=ex]
Contract:
    first_unique(s: str) -> str   # the char, or '_' if none

Hint:
    Three prelude names, one line: Counter the string, find the first
    char whose count is 1, fromMaybe the '_' default. Two passes,
    O(n) -- and find stops at the first hit.
\end{lstlisting}
\begin{lstlisting}[style=ex]
>>> first_unique("swiss")
'w'
>>> first_unique("aabb")
'_'
>>> first_unique("")
'_'
\end{lstlisting}

\section{log\_report: Error report from messy logs}
\begin{lstlisting}[style=ex]
A well-formed line is "<ts> <LEVEL> <service>: <message...>": LEVEL
in {INFO, WARN, ERROR}, service nonempty and glued to its colon,
message may be empty. Blank lines are ignored; any other non-blank
line is malformed.

Contract:
    error_report(loglines: list[str]) -> (report, malformed_count)
    report = [(service, error_count)] for services with >= 1 ERROR,
    sorted by count desc, then name asc.

Hint:
    THE mapMaybe shape: parse(line) -> (level, service) or None.
    The Nones ARE the malformed count; catMaybes + a filter feed
    Counter; sortOn (-count, name) ranks. Messy input, no raw loops.
\end{lstlisting}
\begin{lstlisting}[style=ex]
>>> logs = ["t1 ERROR api: down", "t2 INFO web: ok", "garbage line",
...         "t3 ERROR api: db", "t4 ERROR web: 500", "t5 WARN api: slow",
...         "t6 DEBUG api: no", "t7 ERROR db missing"]
>>> error_report(logs)
([('api', 2), ('web', 1)], 3)
>>> error_report(["t1 INFO web: ok", "half a line"])
([], 1)
>>> error_report([])
([], 0)
\end{lstlisting}

\section{invert\_dict: Swap a dict's keys and values}
\begin{lstlisting}[style=ex]
Values are unique (so the inversion is lossless).

Contract:
    invert(d: dict) -> dict

Hint:
    A dict is a list of pairs wearing a hat: items(), swap each pair,
    dict() the result. The swap showcase.
\end{lstlisting}
\begin{lstlisting}[style=ex]
>>> invert({'a': 1, 'b': 2})
{1: 'a', 2: 'b'}
>>> invert({})
{}
\end{lstlisting}
\chapter{Tries And Paths}

\section{deepest\_directory: Challenge: deepest path in a directory forest}
\begin{lstlisting}[style=ex]
Contract:
    deepest(dirs: list[str]) -> str
    Each string is a '/'-separated path ("a/b/c"). Return the deepest
    full path (unique deepest guaranteed). Paths may share prefixes.

Hint:
    Building IS walking: foldl indexing over the parts autovivifies the
    whole branch -- no setpath needed since there is no leaf value.
    Then one paths() + maxOn keypath length. Note why setpath would
    be WRONG here: writing "a/b" after "a/b/c" would clobber the
    subtree.
\end{lstlisting}
\begin{lstlisting}[style=ex]
>>> deepest(["a/b/c", "a/d", "e"])
'a/b/c'
>>> deepest(["x"])
'x'
>>> deepest(["a/b", "a/b/c/d", "a/b/c"])
'a/b/c/d'
\end{lstlisting}

\section{autocomplete: Challenge: prefix trie with completion}
\begin{lstlisting}[style=ex]
Contract:
    make_trie(words: list[str]) -> ITree
        Letter-by-letter trie; each word terminated by '$' whose leaf
        is the word itself.
    complete(trie, prefix: str) -> list[str]
        All stored words starting with prefix, sorted. A word counts
        as its own completion.

Hint:
    setpath writes each word; getpath descends to the prefix's subtree
    WITHOUT autovivifying (reading an ITree with [] would plant ghost
    branches); paths() harvests every '$' leaf below.
    (No-trie baseline: filter_ + isPrefixOf over the word list, O(n*m)
    per query -- the trie is what repeated queries buy.)
\end{lstlisting}
\begin{lstlisting}[style=ex]
>>> make_trie(["at"]) == {'a': {'t': {'$': 'at'}}}
True
>>> t = make_trie(["cat", "car", "card", "dog", "do"])
>>> complete(t, "ca")
['car', 'card', 'cat']
>>> complete(t, "do")
['do', 'dog']
>>> complete(t, "z")
[]
>>> 'z' in t
False
\end{lstlisting}

\section{subset\_sum\_paths: Challenge: all paths in the decision tree summing to a total}
\begin{lstlisting}[style=ex]
Contract:
    build_tree(nums: list[int], total: int) -> tree
        The take/skip decision tree over nums. Node = dict with keys
        ('take', x) and ('skip', x); leaf = bool (True iff the takes
        along that path sum to total).
    tree_paths(tree) -> list[list[int]]
        The taken values along every path ending in True, take-first
        DFS order.
    sum_paths(nums: list[int], total: int) -> list[list[int]]
        Every subsequence of nums summing to total. Negatives allowed,
        so a path only ends at full depth -- rem == 0 mid-path proves
        nothing.

Hint:
    An unfold builds the tree from seed (nums, remaining); a fold over
    the tree collects paths ending in True. unfold + fold = hylomorphism:
    fuse them and the tree never exists -- but here it is the exhibit.
    2^n leaves: every subset gets judged, exactly once.
    (The prelude now also ships subsequences -- filtering the powerset
    by sum is the same 2^n with the tree left implicit.)
\end{lstlisting}
\begin{lstlisting}[style=ex]
>>> build_tree([], 0)
True
>>> build_tree([5], 5) == {('take', 5): True, ('skip', 5): False}
True
>>> tree_paths(True)
[[]]
>>> tree_paths(False)
[]
>>> tree_paths(build_tree([1, 2], 3))
[[1, 2]]
>>> sum_paths([1, 2, 3], 3)
[[1, 2], [3]]
>>> sum_paths([2, 4, 6, 10], 16)
[[2, 4, 10], [6, 10]]
>>> sum_paths([1, -1, 2], 2)
[[1, -1, 2], [2]]
>>> sum_paths([], 0)
[[]]
>>> sum_paths([1, 1], 2)
[[1, 1]]
\end{lstlisting}

\section{subset\_sum\_trie: Challenge: subset-sum, winners recorded in a solution trie}
\begin{lstlisting}[style=ex]
Contract:
    solution_trie(nums: list[int], total: int) -> ITree
        Search the take/skip space; write only the WINNING paths (the
        taken values, terminated by '$') into an ITree.
    sum_paths(nums: list[int], total: int) -> list[list[int]]
        The winning paths read back off the trie via paths().
        NB: a trie stores a SET of solutions -- duplicates arising from
        different index choices (e.g. [1, 1] with total 1) collapse to one.
        That is a semantic change from 03's multiset.

Hint:
    The full trio: unfold searches, ITree records, paths() reads.
    Why the '$' terminator? One solution can be a prefix of another
    ([2] inside [2, 4, -4]) -- without an end marker the shorter one
    vanishes into the longer one's interior.
\end{lstlisting}
\begin{lstlisting}[style=ex]
>>> solution_trie([1, 2], 3) == {1: {2: {'$': True}}}
True
>>> sum_paths([1, 2, 3], 3)
[[1, 2], [3]]
>>> sum_paths([2, 4, -4], 2)
[[2, 4, -4], [2]]
>>> sum_paths([1, 1], 1)
[[1]]
>>> sum_paths([], 0)
[[]]
\end{lstlisting}
\chapter{Graphs Grids Heaps}

\section{longest\_unique: Longest substring without repeats}
\begin{lstlisting}[style=ex]
Contract:
    longest_unique(s: str) -> int

Hint:
    The longest_window skeleton does the two-pointer bookkeeping; you
    supply the state as closures: per-char counts (Tree(1, int)) plus
    a duplicate tally. valid() == "no duplicates in the window".
\end{lstlisting}
\begin{lstlisting}[style=ex]
>>> longest_unique("abcabcbb")
3
>>> longest_unique("bbbbb")
1
>>> longest_unique("")
0
>>> longest_unique("abba")
2
\end{lstlisting}

\section{grid\_path: Shortest path through a grid}
\begin{lstlisting}[style=ex]
Steps from top-left to bottom-right through '.' cells, around '#'
walls, moving orthogonally. -1 if unreachable (or an endpoint is a
wall); a 1x1 open grid is 0.

Contract:
    grid_path(grid: list[str]) -> int

Hint:
    BFS: the prelude deque is the frontier, neighbors4 the edges.
    Mark seen when you ENQUEUE, not when you pop.
\end{lstlisting}
\begin{lstlisting}[style=ex]
>>> grid_path(["..", ".."])
2
>>> grid_path([".#", "#."])
-1
>>> grid_path(["."])
0
>>> grid_path(["..#", "#.#", "#.."])
4
\end{lstlisting}

\section{rate\_limiter: Sliding-window rate limiter}
\begin{lstlisting}[style=ex]
A request at time ts is denied when its client already has `limit`
ALLOWED requests in (ts - window, ts]. Denied requests consume no
budget. Timestamps arrive non-decreasing.

Contract:
    rate_limiter(window: int, limit: int,
                 events: list[tuple[int, str]]) -> list[tuple[int, str]]
    The denied (ts, client) pairs, in stream order.

Hint:
    Tree(1, deque): each client's allowed timestamps in a deque. The
    prelude deque has no peek -- evict by popleft, and appendleft the
    survivor back. Each stamp enters and leaves once: O(1) amortized.
\end{lstlisting}
\begin{lstlisting}[style=ex]
>>> rate_limiter(10, 2, [(1, 'alice'), (2, 'alice'), (3, 'alice'),
...                      (5, 'bob'), (12, 'alice'), (13, 'alice')])
[(3, 'alice')]
>>> rate_limiter(5, 3, [(1, 'a'), (2, 'a'), (3, 'a')])
[]
>>> rate_limiter(5, 1, [(1, 'a'), (2, 'b'), (3, 'a'), (7, 'a'), (8, 'b')])
[(3, 'a')]
\end{lstlisting}

\section{running\_median: Median after each arrival}
\begin{lstlisting}[style=ex]
Even counts average the two middle values; one decimal is the
caller's problem -- return floats.

Contract:
    running_median(xs: list[int]) -> list[float]

Hint:
    Two heaps: lo holds the lower half NEGATED (the prelude heap is
    min-only), hi the upper half. Push through lo into hi, rebalance
    so lo keeps the extra; the median reads off the tops.
\end{lstlisting}
\begin{lstlisting}[style=ex]
>>> running_median([12, 4, 5, 3, 8, 7])
[12.0, 8.0, 5.0, 4.5, 5.0, 6.0]
>>> running_median([10])
[10.0]
>>> running_median([1, 2, 3, 4])
[1.0, 1.5, 2.0, 2.5]
\end{lstlisting}

\section{dijkstra: Challenge: shortest paths from weighted (from, to, dist) edges}
\begin{lstlisting}[style=ex]
Contract:
    dijkstra(edges: list[tuple], src) -> (dist: dict, parent: dict)
        Directed edges (u, v, w), w >= 0. dist[n] = shortest distance
        src -> n; unreachable nodes absent. parent[n] = predecessor on
        one shortest path (parent[src] = None).
    shortest_path(edges, src, dst) -> list | None
        The node sequence of a shortest src -> dst path, None if
        unreachable.

Hint:
    Adjacency map = Tree(1, list) fold over edges. Frontier = the heap.
    Lazy deletion: skip a popped entry unless it matches the current
    best -- the heap may hold stale distances. Path reconstruction =
    an unfold: iterate the parent chain, takeWhile not None, reverse.
    Graphs are not trees -- cycles make tree unfolds diverge, so fold
    over STATES to a fixed point instead.
\end{lstlisting}
\begin{lstlisting}[style=ex]
>>> dist, _ = dijkstra([('a','b',1), ('b','c',2), ('a','c',4)], 'a')
>>> dist == {'a': 0, 'b': 1, 'c': 3}
True
>>> dist, _ = dijkstra([('a','b',5), ('a','c',1), ('c','b',1)], 'a')
>>> dist['b']
2
>>> dijkstra([('x','y',1)], 'y')[0]
{'y': 0}
>>> shortest_path([('a','b',5), ('a','c',1), ('c','b',1)], 'a', 'b')
['a', 'c', 'b']
>>> shortest_path([('a','b',1)], 'a', 'a')
['a']
>>> shortest_path([('a','b',1)], 'b', 'a') is None
True
\end{lstlisting}

\section{dag\_longest\_path: Challenge: longest path in a weighted DAG}
\begin{lstlisting}[style=ex]
Contract:
    dag_longest(edges: list[tuple]) -> (total: int, path: list)
    Directed acyclic edges (u, v, w). The maximum-weight path anywhere
    in the DAG and its node sequence. Ties broken toward the
    lexically-first end node. Input guaranteed acyclic.

Hint:
    Longest path is NP-hard on general graphs; acyclicity is what
    makes it linear. Kahn's topo sort (indegree fold + deque), then a
    fold over topo order relaxing best[v] = max(best[v], best[u] + w).
    Same skeleton as Dijkstra's relax with max for min and topo order
    replacing the heap -- the LCS-vs-edit-distance move on graphs.
    maxOn best picks the end node; feed it the SORTED nodes so
    ties break lexically.
\end{lstlisting}
\begin{lstlisting}[style=ex]
>>> dag_longest([('a','b',3), ('b','c',4), ('a','c',5)])
(7, ['a', 'b', 'c'])
>>> dag_longest([('s','a',1), ('s','b',5), ('a','t',6), ('b','t',1)])
(7, ['s', 'a', 't'])
>>> dag_longest([('u','v',2)])
(2, ['u', 'v'])
\end{lstlisting}

\section{bellman\_ford: Challenge: shortest paths as a fixed point (negative weights ok)}
\begin{lstlisting}[style=ex]
Contract:
    bellman_ford(edges: list[tuple], src) -> dict | None
    Directed edges (u, v, w); w may be negative. dist[n] = shortest
    distance src -> n, unreachable nodes absent. Returns None if a
    negative cycle is reachable from src (distances then have no
    fixed point -- they diverge to -inf).

Hint:
    Shortest paths ARE a fixed point: dist is the least map satisfying
    dist[v] <= dist[u] + w for every edge. So: relax ALL edges as one
    step function, `until` the map stops changing. Convergence within
    |V| - 1 steps is guaranteed (a shortest path has at most |V| - 1
    edges); a change on step |V| proves a negative cycle. Dijkstra is
    the same relax with a clever ORDER; Bellman-Ford orders nothing
    and pays for it in passes.
\end{lstlisting}
\begin{lstlisting}[style=ex]
>>> bellman_ford([('a','b',1), ('b','c',2), ('a','c',4)], 'a') == {'a': 0, 'b': 1, 'c': 3}
True
>>> bellman_ford([('a','b',5), ('b','c',-4), ('a','c',2)], 'a')['c']
1
>>> bellman_ford([('a','b',1), ('b','c',-3), ('c','b',1), ('b','d',1)], 'a') is None
True
>>> bellman_ford([('x','y',7)], 'y')
{'y': 0}
>>> bellman_ford([('a','b',-2), ('b','a',3)], 'a') == {'a': 0, 'b': -2}
True
\end{lstlisting}

\section{course\_schedule: NeetCode: Course Schedule (cycle detection via Kahn)}
\begin{lstlisting}[style=ex]
Contract:
    can_finish(n: int, prereqs: list[tuple]) -> bool
        Courses 0..n-1; (a, b) means b must precede a. True iff every
        course can be taken, i.e. the prerequisite graph is acyclic.
    find_order(n: int, prereqs: list[tuple]) -> list | None
        A valid course order (smallest-numbered course first on ties),
        or None if impossible.

Hint:
    Kahn's insight inverted: topo sort doesn't FAIL on a cycle, it
    just leaves the cycle behind -- nodes on a cycle never reach
    indegree 0. So: run 06's Kahn, then len(topo) == n is the whole
    cycle test. Heap instead of deque only to make ties deterministic.
\end{lstlisting}
\begin{lstlisting}[style=ex]
>>> find_order(2, [(1, 0)])
[0, 1]
>>> find_order(4, [(1, 0), (2, 0), (3, 1), (3, 2)])
[0, 1, 2, 3]
>>> find_order(2, [(1, 0), (0, 1)]) is None
True
>>> find_order(3, [])
[0, 1, 2]
>>> can_finish(2, [(1, 0)])
True
>>> can_finish(2, [(1, 0), (0, 1)])
False
>>> can_finish(5, [(1, 4), (2, 4), (3, 1), (3, 2)])
True
\end{lstlisting}
\chapter{Dp And Control}

\section{coin\_change: NeetCode: Coin Change}
\begin{lstlisting}[style=ex]
Contract:
    coin_change(coins: list[int], amount: int) -> int
    Fewest coins summing to amount (unlimited supply), -1 if impossible.
    amount >= 0, coins positive.

Hint:
    Knapsack, dp[remaining]. Top-down with memo on the single
    index `amount` -- the future only needs what remains.
\end{lstlisting}
\begin{lstlisting}[style=ex]
>>> coin_change([1, 2, 5], 11)
3
>>> coin_change([2], 3)
-1
>>> coin_change([1], 0)
0
>>> coin_change([186, 419, 83, 408], 6249)
20
\end{lstlisting}

\section{lcs: NeetCode: Longest Common Subsequence}
\begin{lstlisting}[style=ex]
Contract:
    lcs(a: str, b: str) -> int
    Length of the longest common subsequence (not substring).

Hint:
    Alignment pattern, dp[i][j] -- same skeleton as 03's edit
    distance, but max instead of min. Index-based args for the cache.
\end{lstlisting}
\begin{lstlisting}[style=ex]
>>> lcs("abcde", "ace")
3
>>> lcs("abc", "abc")
3
>>> lcs("abc", "def")
0
>>> lcs("", "abc")
0
\end{lstlisting}

\section{edit\_distance: Levenshtein distance}
\begin{lstlisting}[style=ex]
Minimum single-character insertions, deletions, or substitutions
turning a into b.

Contract:
    edit_distance(a: str, b: str) -> int

Hint:
    The alignment DP again -- 02's skeleton with min-of-three plus
    one. Matching heads cost nothing; otherwise pay 1 and recurse on
    the three repairs. memo over the index pair.
\end{lstlisting}
\begin{lstlisting}[style=ex]
>>> edit_distance("kitten", "sitting")
3
>>> edit_distance("", "abc")
3
>>> edit_distance("same", "same")
0
>>> edit_distance("flaw", "lawn")
2
>>> edit_distance("a", "b")
1
\end{lstlisting}

\section{knapsack: 0/1 knapsack}
\begin{lstlisting}[style=ex]
Each (weight, value) item usable at most once; maximize value within
capacity.

Contract:
    knapsack(capacity: int, items: list[tuple[int, int]]) -> int

Hint:
    memo over (i, room): take item i or skip it. Two integers of
    state -- the cache keys on the args tuple, so keep them hashable.
\end{lstlisting}
\begin{lstlisting}[style=ex]
>>> knapsack(10, [(5, 10), (4, 40), (6, 30), (3, 50)])
90
>>> knapsack(0, [(1, 100)])
0
>>> knapsack(5, [])
0
>>> knapsack(3, [(4, 100), (2, 7)])
7
\end{lstlisting}

\section{calculator: Integer expression evaluator (capstone)}
\begin{lstlisting}[style=ex]
+, -, * and parentheses, optional spaces, normal precedence,
left-associative. No unary minus, no division.

Contract:
    evaluate(expr: str) -> int

Grammar (one function per rule, shared position cursor):
    expression := term (('+' | '-') term)*
    term       := factor ('*' factor)*
    factor     := NUMBER | '(' expression ')'

Hint:
    Recursive descent. Lex numbers with takeWhile str.isdigit on the
    remaining input -- span IS the lexer. And remember the trap you
    certified earlier: '' in "+-" is True; test membership against a
    tuple.
\end{lstlisting}
\begin{lstlisting}[style=ex]
>>> evaluate("2*(3+4)")
14
>>> evaluate("2*(3+4)-5")
9
>>> evaluate("10-2-3")
5
>>> evaluate("2+3*4")
14
>>> evaluate("((1+2))*3")
9
>>> evaluate(" 12 * 2 ")
24
\end{lstlisting}

\section{max\_height\_cuboids: LC 1691: Maximum Height by Stacking Cuboids}
\begin{lstlisting}[style=ex]
Any rotation allowed; cuboid A sits on B only when ALL THREE of A's
dims are <= B's, compared side to side. Each cuboid used at most
once. Maximize stack height.

Contract:
    max_height(cuboids: list[tuple[int, int, int]]) -> int

Hint:
    The rotation freedom COLLAPSES: sort each cuboid's own dims
    ascending -- the biggest dim as height is always at least as good,
    and sorted-vs-sorted is the only comparison that can succeed (the
    photo_lineup exchange argument, twice). Then sort the cuboids and
    run weighted LIS: best ending at i = h_i + best over j fitting
    under i. Fold-with-history, O(n^2) -- the g4_07 shape with sum
    instead of length. (g8_06's mask version is the ground truth
    oracle for THIS file: same >= 17-style surprises, different fit
    rule -- all three dims here, base only there.)
\end{lstlisting}
\begin{lstlisting}[style=ex]
>>> max_height([(50, 45, 20), (95, 37, 53), (45, 23, 12)])
190
>>> max_height([(38, 25, 45), (76, 35, 3)])
76
>>> max_height([(7, 11, 17), (7, 17, 11), (11, 7, 17), (11, 17, 7), (17, 7, 11), (17, 11, 7)])
102
>>> max_height([(1, 1, 1)])
1
\end{lstlisting}

\section{guillotine\_cut: Max value cut from a sheet, guillotine cuts only}
\begin{lstlisting}[style=ex]
A W x H sheet; piece types (pw, ph, value), unlimited supply, fixed
orientation (add rotated copies yourself if rotation is allowed).
Every cut runs edge to edge, so each cut splits one rectangle into
two rectangles. Unused offcuts score 0.

Contract:
    guillotine(W: int, H: int, pieces: list[tuple[int, int, int]]) -> int

Hint:
    Guillotine is what re-admits DP: free-form 2D packing has no
    small state (the frontier goes irregular), but edge-to-edge cuts
    keep every subproblem a RECTANGLE -- (w, h), two ints. Interval
    pattern: best(w,h) = max of a piece that fits the sheet EXACTLY,
    every vertical split x, every horizontal split y. A smaller piece
    never needs a direct case: the cut that isolates it is one of the
    splits. memo over (w, h); O(W*H*(W+H)) -- say it.
\end{lstlisting}
\begin{lstlisting}[style=ex]
>>> guillotine(4, 4, [(1, 1, 1)])
16
>>> guillotine(4, 4, [(3, 3, 7), (1, 1, 1)])
16
>>> guillotine(4, 4, [(3, 3, 20), (1, 1, 1)])
27
>>> guillotine(4, 4, [(4, 2, 10), (2, 2, 3)])
20
>>> guillotine(3, 3, [(2, 2, 5)])
5
>>> guillotine(2, 2, [(3, 1, 99)])
0
\end{lstlisting}

\section{partition\_k\_subsets: LC 698: Partition to K Equal Sum Subsets}
\begin{lstlisting}[style=ex]
Contract:
    can_partition(nums: list[int], k: int) -> bool
    True iff nums splits into k non-empty groups of equal sum.
    Positive ints, small n (say the bound: 2^n states).

Hint:
    The mask family, pure form: the future needs (which numbers
    remain, how much room is left in the group being filled). When a
    group closes (room hits 0) the room resets to target -- so the
    state is just (used, room): the number of CLOSED groups is
    implied by the used-sum. memo over that pair; fill greedily
    into one group at a time to avoid counting group orderings.
    (g8_06's go is this skeleton plus geometry.)
\end{lstlisting}
\begin{lstlisting}[style=ex]
>>> can_partition([4, 3, 2, 3, 5, 2, 1], 4)
True
>>> can_partition([1, 2, 3, 4], 3)
False
>>> can_partition([2, 2, 2, 2], 2)
True
>>> can_partition([1], 1)
True
>>> can_partition([1, 1], 3)
False
\end{lstlisting}

\section{min\_path\_sum: Cheapest right/down path through a grid}
\begin{lstlisting}[style=ex]
Contract:
    min_path_sum(grid: list[list[int]]) -> int
    Top-left to bottom-right, moving only right or down.

Hint:
    Grid DP: state (r, c) -- the future only cares where you stand.
    memo over the index pair; base case is the goal cell.
    Compare g7_02: uniform steps there needed BFS; weighted cells
    with only right/down moves make a DAG, so plain DP suffices.
\end{lstlisting}
\begin{lstlisting}[style=ex]
>>> min_path_sum([[1, 3, 1], [1, 5, 1], [4, 2, 1]])
7
>>> min_path_sum([[5]])
5
>>> min_path_sum([[1, 2], [3, 4]])
7
\end{lstlisting}
\chapter{Stack Folds}

\section{balanced\_brackets: Matching pairs: valid bracket nesting}
\begin{lstlisting}[style=ex]
Only ()[]{} appear.

Contract:
    balanced(s: str) -> bool

Hint:
    The stack fold, purest form: a plain list IS a stack, and the
    fold's accumulator is one. The twist worth learning: use None as
    the FAILED state -- once a closer mismatches, the poison
    propagates through the rest of the fold untouched (Maybe as
    accumulator). Balanced == the fold ends on an empty stack.
\end{lstlisting}
\begin{lstlisting}[style=ex]
>>> balanced("([{}])")
True
>>> balanced("([)]")
False
>>> balanced("(((")
False
>>> balanced("")
True
>>> balanced("]")
False
\end{lstlisting}

\section{eval\_rpn: Evaluate Reverse Polish notation}
\begin{lstlisting}[style=ex]
Integer operands, operators + - *.

Contract:
    eval_rpn(tokens: list[str]) -> int

Hint:
    One fold over the tokens: operands push; an operator pops two and
    pushes the result. The classic slip: stack[-2] is the LEFT
    operand. The answer is whatever the fold leaves on the stack.
\end{lstlisting}
\begin{lstlisting}[style=ex]
>>> eval_rpn(["2", "3", "+", "4", "*"])
20
>>> eval_rpn(["5", "1", "2", "+", "4", "*", "-"])
-7
>>> eval_rpn(["7"])
7
\end{lstlisting}

\section{next\_greater: First later element that's larger}
\begin{lstlisting}[style=ex]
For each element, the next strictly greater element to its right;
-1 if none.

Contract:
    next_greater(xs: list[int]) -> list[int]

Hint:
    The MONOTONIC stack: fold over enum(xs); the newcomer pops every
    stacked index whose value it beats (that newcomer is their
    answer), then pushes itself. Each index pushes and pops at most
    once -- O(n) despite the inner while. The pro tier of stacks.
\end{lstlisting}
\begin{lstlisting}[style=ex]
>>> next_greater([2, 1, 2, 4, 3])
[4, 2, 4, -1, -1]
>>> next_greater([5, 4, 3])
[-1, -1, -1]
>>> next_greater([])
[]
\end{lstlisting}
\appendix
\part{Appendices}
\chapter{prelude.py, Complete}
\begin{lstlisting}[style=file]
"""prelude.py -- Haskell-style Prelude in pure Python, in pedagogical order.

Pascal discipline: strict define-before-use. Read top to bottom and every
name is already defined when you meet it. Sections build on each other:

   1 combinators          2 pairs                3 arithmetic
   4 list basics          5 higher-order lists   6 folds
   7 scans                8 streams (lazy)       9 slicing & spans
  10 sorting & searching 11 strings             12 containers
  13 nodes               14 grids & windows     15 control

Within a section names run alphabetically, except that a definition is
hoisted ahead of the first name that uses it (foldl before compose, span
before break_), so the discipline holds inside sections too.

Import qualified (import prelude as P): Haskell names shadow builtins by design.
Beyond the Prelude proper: the working parts of Data.List, Data.Map, Data.Maybe,
Data.Function and friends are folded into the same sections -- grouped by what
they do, not by the module they came from. None is Nothing throughout.
foldr = foldl(flip(f), reversed(xs), init): iterative, stack-safe, valid only
because Python is strict.
"""

# ============ 1. combinators ============ (functions about functions)

const     = lambda x: lambda _: x
curry     = lambda f: lambda x: lambda y: f(x, y)
flip      = lambda f: (lambda x, y: f(y, x))          # flip f x y = f y x
fromMaybe = lambda d, x: d if x is None else x        # Data.Maybe: default for every None-returning tool
id        = lambda x: x                               # shadows builtin id() by design
isJust    = lambda x: x is not None                   # Data.Maybe: None test as a handable predicate
isNothing = lambda x: x is None
NOTHING   = object()                                  # Maybe's Nothing: "no arg given"; test with `is`
on        = lambda f, g: lambda x, y: f(g(x), g(y))   # Data.Function: combine via a key -- cmp `on` fst
partial   = lambda f, *bound: lambda *args: f(*bound, *args)
uncurry   = lambda f: lambda p: f(*p)

# ============ 2. pairs ============

fst  = lambda p: p[0]
snd  = lambda p: p[1]
swap = lambda p: (p[1], p[0])   # Data.Tuple

# ============ 3. arithmetic & logic ============

div       = lambda a, b: a // b                                          # floors, like Haskell div
even      = lambda n: n % 2 == 0
gcd       = lambda a, b: abs(a) if b == 0 else gcd(b, a % b)             # >= 0 like Haskell; depth <= 90
hypot     = lambda x, y: (x*x + y*y) ** 0.5

def isqrt(n):  # floor sqrt without floats (Newton)
    if n < 0:
        raise ValueError("isqrt of negative")
    x = n
    y = (x + 1) // 2
    while y < x:
        x, y = y, (y + n // y) // 2
    return x if n else 0

lcm       = lambda a, b: abs(a // gcd(a, b) * b) if a and b else 0       # >= 0 like Haskell
mod       = lambda a, b: a % b
not_      = lambda b: not b                                              # `not` is a Python keyword
odd       = lambda n: n % 2 == 1
otherwise = True
pred      = lambda x: chr(ord(x) - 1) if isinstance(x, str) else x - 1   # Enum: numbers and Chars
quot      = lambda a, b: -(-a // b) if (a < 0) != (b < 0) else a // b    # truncates, like Haskell quot
rem       = lambda a, b: a - b * quot(a, b)
quotRem   = lambda a, b: (quot(a, b), rem(a, b))
signum    = lambda x: (x > 0) - (x < 0)
succ      = lambda x: chr(ord(x) + 1) if isinstance(x, str) else x + 1

# ============ 4. list basics ============

# drop: finite lists only; n<0 drops nothing
drop        = lambda n, xs: list(xs)[max(n, 0):]
elem        = lambda x, xs: x in xs
zip_        = lambda a, b: list(zip(a, b))                                 # shortest wins, any iterable
enum        = lambda xs, start=0: zip_(list(range(start, start + len(xs))), xs)
head        = lambda xs: xs[0]
init        = lambda xs: xs[:-1]
# Data.List isPrefixOf / isSuffixOf: strings and lists alike
isPrefixOf  = lambda p, xs: list(xs[:len(p)]) == list(p)
# NB isSuffixOf slices xs[len(xs)-len(p):], not xs[-len(p):] -- [-0:] is ALL
isSuffixOf  = lambda p, xs: list(xs[len(xs) - len(p):]) == list(p)
last        = lambda xs: xs[-1]
lookup      = lambda k, pairs: next((v for kk, v in pairs if kk == k), None)   # Nothing -> None
notElem     = lambda x, xs: x not in xs
# nub: unique, first occurrence wins (dicts keep order)
nub         = lambda xs: list(dict.fromkeys(xs))
null        = lambda xs: len(xs) == 0
pairwise    = lambda xs: list(zip(xs, xs[1:]))                             # zip xs (tail xs)
replicate   = lambda n, x: [x] * n
splitAt     = lambda n, xs: (xs[:n], xs[n:]) if n >= 0 else (xs[:0], xs)   # n<0 splits at the front
stripPrefix = lambda p, xs: xs[len(p):] if isPrefixOf(p, xs) else None     # Just the rest, or Nothing
tail        = lambda xs: xs[1:]
def transpose(xss):                    # Data.List: rows <-> cols, RAGGED-SAFE like Haskell
    xss = [list(xs) for xs in xss]     # any iterables in (rows may be one-shot); materialise once
    n = max(map(len, xss), default=0)  # short rows just drop out of later columns (no truncation)
    return [[xs[i] for xs in xss if i < len(xs)] for i in range(n)]
unzip       = lambda ps: tuple(map(list, zip(*ps))) if ps else ([], [])
zip3        = lambda a, b, c: list(zip(a, b, c))

# ============ 5. higher-order lists ============

catMaybes = lambda xs: [x for x in xs if x is not None]               # mapMaybe id
concat    = lambda xss: [x for xs in xss for x in xs]
concatMap = lambda f, xs: [y for x in xs for y in f(x)]
cross     = lambda a, b: [(x, y) for x in a for y in b]               # cartesian; `product` is the fold
filter_   = lambda pred, xs: [x for x in xs if pred(x)]
# Data.List find: lazy first-match -- folds can't stop early, find can
find      = lambda pred, xs: next((x for x in xs if pred(x)), None)   # first match, else None
map_      = lambda f, xs: [f(x) for x in xs]
# Data.Maybe mapMaybe: map and drop the Nothings in ONE pass -- the parse-and-filter shape
mapMaybe  = lambda f, xs: [y for y in (f(x) for x in xs) if y is not None]
def partition(pred, xs):                   # (keepers, rest) in ONE pass; pred called once per element
    yes, no = [], []
    for x in xs:
        (yes if pred(x) else no).append(x)
    return (yes, no)
starmap   = lambda f, pairs: [f(*p) for p in pairs]                   # map . uncurry
zipWith   = lambda f, a, b: [f(x, y) for x, y in zip(a, b)]
zipWith3  = lambda f, a, b, c: [f(x, y, z) for x, y, z in zip(a, b, c)]

# ============ 6. folds ============ (collapse a list to one value)

def foldl(f, xs, init=NOTHING):  # THE left fold; Python buried its own in functools as reduce
    it = iter(xs)
    if init is NOTHING:
        try:
            acc = next(it)
        except StopIteration:
            raise TypeError("fold of empty sequence with no initial value")
    else:
        acc = init
    for x in it:
        acc = f(acc, x)
    return acc

foldr   = lambda f, xs, init: foldl(flip(f), reversed(list(xs)), init)
compose = lambda *fns: lambda x: foldr(lambda f, acc: f(acc), fns, x)   # foldr (.) id: RIGHTMOST first
foldl1  = lambda f, xs: foldl(f, xs)
foldr1  = lambda f, xs: foldl(flip(f), reversed(list(xs)))
product = lambda xs: foldl(lambda a, x: a * x, xs, 1)                   # the numeric fold

# enumeration folds -- brute-force licenses for small n (say the bound out loud):
# Control.Monad replicateM: every length-n word over xs -- cross, n times; |xs|^n
replicateM   = lambda n, xs: foldl(lambda acc, _: [s + [x] for s in acc for x in xs], range(n), [[]])
# Data.List subsequences: the powerset; each element DOUBLES acc -- 2^n, fine to n ~ 20
subsequences = lambda xs: foldl(lambda acc, x: acc + [s + [x] for s in acc], list(xs), [[]])

# ============ 7. scans ============ (a fold that keeps its history)

def accumulate(xs, f=None, initial=NOTHING):     # scanl / scanl1
    if f is None:
        f = lambda a, b: a + b
    it = iter(xs)
    if initial is NOTHING:
        try:
            acc = next(it)
        except StopIteration:
            return
    else:
        acc = initial
    yield acc
    for x in it:
        acc = f(acc, x)
        yield acc

scanl  = lambda f, z, xs: list(accumulate(xs, f, initial=z))
scanl1 = lambda f, xs: list(accumulate(xs, f))
scanr  = lambda f, z, xs: list(reversed(scanl(flip(f), z, list(reversed(list(xs))))))
scanr1 = lambda f, xs: list(reversed(scanl1(flip(f), list(reversed(list(xs))))))

# ============ 8. streams ============ (lazy, possibly infinite)

chain     = lambda *iterables: (x for it in iterables for x in it)   # concat, lazily

def count(start=0, step=1):                 # [start, start+step ..]
    n = start
    while True:
        yield n
        n += step

def cycle(xs):
    xs = list(xs)
    if not xs:
        raise ValueError("cycle: empty list")   # Haskell errors too; the alternative is a silent hang
    while True:
        yield from xs

def iterate(f, x):                          # iterate f x = [x, f x, f (f x), ..]
    while True:
        yield x
        x = f(x)

def repeat(x, n=None):
    if n is None:
        while True:
            yield x
    else:
        for _ in range(n):
            yield x

def unfoldr(f, seed):                       # Data.List unfoldr: f(seed) -> None | (value, seed')
    out = []                                # iterate's finite twin: grows a list until f says stop --
    step = f(seed)                          # no more threading (state, acc) tuples through until
    while step is not None:
        val, seed = step
        out.append(val)
        step = f(seed)
    return out

# ============ 9. slicing & spans ============ (finite views of streams)

def span(p, xs):                            # split at first failure, ONE pass; safe on one-shot iters
    xs = list(xs)
    i = next((i for i, x in enumerate(xs) if not p(x)), len(xs))
    return (xs[:i], xs[i:])
break_    = lambda p, xs: span(lambda x: not p(x), xs)
chunksOf  = lambda n, xs: [xs[i:i + n] for i in range(0, len(xs), n)]   # Data.List.Split; short last ok

def dropwhile(pred, xs):
    it = iter(xs)
    for x in it:
        if not pred(x):
            yield x
            break
    yield from it

dropWhile = lambda p, xs: list(dropwhile(p, xs))
inits     = lambda xs: [xs[:i] for i in range(len(xs) + 1)]   # Data.List: every prefix, [] first

def islice(iterable, stop):                 # take
    it = iter(iterable)
    for _ in range(stop):
        try:
            yield next(it)
        except StopIteration:
            return

tails     = lambda xs: [xs[i:] for i in range(len(xs) + 1)]   # every suffix; substrings start here
take      = lambda n, xs: list(islice(iter(xs), n))           # works on infinite streams

def takeuntil(pred, xs):                    # like takewhile(not . pred), but INCLUDES the first hit
    for x in xs:
        yield x
        if pred(x):
            return

def takewhile(pred, xs):
    for x in xs:
        if not pred(x):
            return
        yield x

takeWhile = lambda p, xs: list(takewhile(p, xs))

# ============ 10. sorting & searching ============

def bisect_left(a, x):                      # first index where a[i] >= x
    lo, hi = 0, len(a)
    while lo < hi:
        mid = (lo + hi) // 2
        if a[mid] < x:
            lo = mid + 1
        else:
            hi = mid
    return lo

def bisect_right(a, x):                     # first index where a[i] > x
    lo, hi = 0, len(a)
    while lo < hi:
        mid = (lo + hi) // 2
        if a[mid] <= x:
            lo = mid + 1
        else:
            hi = mid
    return lo

def cmp_to_key(cmp):                        # resurrect Python 2's cmp for sortBy
    class K:
        def __init__(self, x): self.x = x
        def __lt__(self, other): return cmp(self.x, other.x) < 0
    return K

maxOn  = lambda f, xs: max(xs, key=f)      # maximumBy (comparing f) in O(n); first wins ties

def merge(a, b):                            # merge two sorted lists, stable; the heart of merge sort
    out, i, j = [], 0, 0                    # (section 13's merge_lists is this same loop on ListNodes)
    while i < len(a) and j < len(b):
        if a[i] <= b[j]:
            out.append(a[i]); i += 1
        else:
            out.append(b[j]); j += 1
    return out + a[i:] + b[j:]

minOn  = lambda f, xs: min(xs, key=f)      # minimumBy (comparing f): sortOn + head without the sort
sortBy = lambda cmp, xs: sorted(xs, key=cmp_to_key(cmp))   # comparator sort, Python 2 style
sortOn = lambda f, xs: sorted(xs, key=f)   # sortOn (schwartzian, f called once per element)

# ============ 11. strings ============

lines   = str.splitlines
unlines = lambda ls: "".join(l + "\n" for l in ls)
unwords = " ".join
words   = str.split

# ============ 12. containers ============

# multiset ops on Counters (use .get, not [] -- indexing a Counter autovivifies zeros):
bag_diff  = lambda a, b: {k: a[k] - b[k] for k in a if a[k] > b.get(k, 0)}    # clipped at 0
bag_inter = lambda a, b: {k: v for k in a.keys() & b.keys() if (v := min(a.get(k, 0), b.get(k, 0)))}
bag_sub   = lambda a, b: all(b.get(k, 0) >= n for k, n in a.items())          # a <= b (inclusion)
bag_union = lambda a, b: {k: max(a.get(k, 0), b.get(k, 0)) for k in a.keys() | b.keys()}

class defaultdict(dict):
    def __init__(self, factory=None, *args, **kw):
        super().__init__(*args, **kw)
        self.factory = factory
    def __missing__(self, key):
        if self.factory is None:
            raise KeyError(key)
        self[key] = self.factory()
        return self[key]

def Counter(xs):                            # count-by-value; a defaultdict(int) fold
    d = defaultdict(int)
    for x in xs:
        d[x] += 1
    return d

class deque:                                # two-stack; all four ends amortized O(1)
    def __init__(self, xs=()):
        self._in, self._out = [], list(reversed(list(xs)))
    def append(self, x): self._in.append(x)
    def appendleft(self, x): self._out.append(x)
    def pop(self):
        if not self._in:
            self._in, self._out = list(reversed(self._out)), []
        return self._in.pop()
    def popleft(self):
        if not self._out:
            self._out, self._in = list(reversed(self._in)), []
        return self._out.pop()
    def __len__(self): return len(self._in) + len(self._out)
    def __iter__(self): return iter(self._out[::-1] + self._in)

def dsu(n):                                 # union-find, path halving; union -> False if already joined
    parent = list(range(n))
    def find(x):
        while parent[x] != x:
            parent[x] = parent[parent[x]]
            x = parent[x]
        return x
    def union(a, b):
        ra, rb = find(a), find(b)
        if ra == rb:
            return False
        parent[ra] = rb
        return True
    return find, union

def fromListWith(f, pairs):                 # Data.Map fromListWith: THE dict-building fold; f(new, old)
    d = {}                                  # like insertWith -- so grouping with concat REVERSES each
    for k, v in pairs:                      # group (the classic Haskell gotcha); use add for counts
        d[k] = f(v, d[k]) if k in d else v
    return d                                # Counter == fromListWith(add) over (x, 1);
                                            # grouping == fromListWith(++) over (k, [v])

def getpath(t, ks, default=None):           # read along ks WITHOUT autovivifying
    for k in ks:
        if not isinstance(t, dict) or k not in t:
            return default
        t = t[k]
    return t

def groupBy(eq, xs):                        # Data.List groupBy: runs of CONSECUTIVE elements; each new
    out = []                                # element is compared via eq against the run's FIRST element,
    for x in xs:                            # exactly like Haskell (span-based) -- not its neighbor
        if out and eq(out[-1][0], x):
            out[-1].append(x)
        else:
            out.append([x])
    return out

group = lambda xs: groupBy(lambda a, b: a == b, xs)   # Data.List group: runs of equals, [[a]] -- no keys

def heappop(h):                             # min-heap on a plain list: sift-down
    h[0], h[-1] = h[-1], h[0]
    top = h.pop()
    i, n = 0, len(h)
    while True:
        s, l, r = i, 2*i + 1, 2*i + 2
        if l < n and h[l] < h[s]: s = l
        if r < n and h[r] < h[s]: s = r
        if s == i:
            return top
        h[i], h[s] = h[s], h[i]
        i = s

def heappush(h, x):                         # min-heap on a plain list: sift-up
    h.append(x)
    i = len(h) - 1
    while i and h[(i - 1) // 2] > h[i]:
        h[(i - 1) // 2], h[i] = h[i], h[(i - 1) // 2]
        i = (i - 1) // 2

insertWith = lambda f, k, v, d: {**d, k: f(v, d[k]) if k in d else v}   # Data.Map insertWith: PERSISTENT
                                                                        # (fresh dict); f(new, old)
ITree = lambda: defaultdict(ITree)

def paths(t):                               # cata for Tree/ITree: [(keypath, leaf)]
    # recursion depth = tree height (word length / len(nums)) -- well under the ~1000 limit
    if not isinstance(t, dict) or not t:    # leaf = non-dict value, or empty node
        return [([], t)]
    return [([k] + p, leaf) for k, sub in t.items() for p, leaf in paths(sub)]

leaves = lambda t: [l for _, l in paths(t)]

def setpath(t, ks, v):                      # write leaf v at key path ks (autovivifies interior)
    for k in ks[:-1]:
        t = t[k]
    t[ks[-1]] = v                           # NB: clobbers any subtree already at ks

Tree  = lambda depth, leaf: (defaultdict(leaf) if depth == 1
                             else defaultdict(lambda: Tree(depth - 1, leaf)))

unionWith = lambda f, a, b: {**a, **{k: f(a[k], v) if k in a else v for k, v in b.items()}}
                                        # Data.Map unionWith: f(a's value, b's value) on shared keys;
                                        # bag_union == unionWith(max), merge-sum == unionWith(add)

# ============ 13. nodes ============ (the LeetCode givens: binary trees & linked lists)

class ListNode:                             # the LeetCode linked list
    def __init__(self, val=0, next=None):
        self.val, self.next = val, next

class TreeNode:                             # the LeetCode binary tree
    def __init__(self, val=0, left=None, right=None):
        self.val, self.left, self.right = val, left, right

from_list = lambda xs: foldr(ListNode, xs, None)   # build; foldr of ListNode

def has_cycle(node):                        # Floyd: fast laps slow iff a cycle exists
    slow = fast = node
    while fast and fast.next:
        slow, fast = slow.next, fast.next.next
        if slow is fast:
            return True
    return False

# tfold: the tree cata -- f(val, l_acc, r_acc), z for the empty tree; recursion depth = tree height
tfold     = lambda f, z, t: z if t is None else f(t.val, tfold(f, z, t.left), tfold(f, z, t.right))
inorder   = lambda t: tfold(lambda v, l, r: l + [v] + r, [], t)

def levelorder(t):                          # BFS: the deque earning its keep
    out, q = [], deque([t] if t else [])
    while len(q):
        n = q.popleft()
        out.append(n.val)
        if n.left:  q.append(n.left)
        if n.right: q.append(n.right)
    return out

def merge_lists(a, b):                      # merge two sorted lists; dummy-head idiom
    dummy = tail_ = ListNode()
    while a and b:
        if a.val <= b.val:
            tail_.next, a = a, a.next
        else:
            tail_.next, b = b, b.next
        tail_ = tail_.next
    tail_.next = a or b
    return dummy.next

def middle(node):                           # fast/slow: second middle for even lengths
    slow = fast = node
    while fast and fast.next:
        slow, fast = slow.next, fast.next.next
    return slow

postorder = lambda t: tfold(lambda v, l, r: l + r + [v], [], t)
preorder  = lambda t: tfold(lambda v, l, r: [v] + l + r, [], t)

def reverse_list(node):                     # three-pointer; prev is a fold accumulator
    prev = None
    while node:
        node.next, prev, node = prev, node, node.next
    return prev

tdepth    = lambda t: tfold(lambda _, l, r: 1 + max(l, r), 0, t)
to_list   = lambda node: list(unfoldr(lambda n: None if n is None else (n.val, n.next), node))
tsize     = lambda t: tfold(lambda _, l, r: 1 + l + r, 0, t)

# ============ 14. grids & windows ============

def longest_window(xs, valid, add, rem):    # sliding-window skeleton; state lives in the closures
    lo = best = 0
    for hi in range(len(xs)):
        add(xs[hi])
        while not valid():                  # shrink until legal again
            rem(xs[lo])
            lo += 1
        best = max(best, hi - lo + 1)
    return best

def neighbors4(r, c, R, C):                 # in-bounds orthogonal neighbors
    for dr, dc in ((1, 0), (-1, 0), (0, 1), (0, -1)):
        if 0 <= r + dr < R and 0 <= c + dc < C:
            yield r + dr, c + dc

def neighbors8(r, c, R, C):                 # ... plus diagonals
    for dr in (-1, 0, 1):
        for dc in (-1, 0, 1):
            if (dr or dc) and 0 <= r + dr < R and 0 <= c + dc < C:
                yield r + dr, c + dc

# ============ 15. control ============

def memo(f):                                # unbounded memoizer; enough for DP (no eviction by design)
    cache = {}
    def wrapped(*args, **kw):
        key = (args, frozenset(kw.items())) if kw else args
        if key not in cache:
            cache[key] = f(*args, **kw)
        return cache[key]
    wrapped.cache = cache                   # peek at the DP table if curious
    return wrapped

def until(p, f, x):                         # until p f x -- loop, not recursion (unbounded depth)
    while not p(x):
        x = f(x)
    return x

\end{lstlisting}
\backmatter
\end{document}
